\documentclass[12pt, oneside]{book}

% ─── 1. Paquetes ─────────────────────────────────────────────────────────────
% Carga todos los paquetes. Debe ir primero — el resto depende de ellos.
%───────────────────────────────────────────────────────────────────────────────
% PACKAGES.TEX — Carga de paquetes
%───────────────────────────────────────────────────────────────────────────────
%
% REGLA DE ESTE ARCHIVO:
%   Solo carga paquetes y declara opciones de carga.
%   La configuración de cada paquete vive en su archivo de config/ respectivo.
%   Excepción: \sisetup y \pgfplotsset{compat} que son inicialización
%   inmediata post-carga, no configuración de dominio.
%
% ORDEN DE CARGA:
%   El orden importa. Las dependencias están marcadas con (→ después de X).
%   Regla fija: hyperref penúltimo, cleveref último.
%
% MOTOR:
%   Este template requiere LuaLaTeX (catppuccinpalette + listings UTF-8).
%   inputenc NO se carga: LuaLaTeX maneja UTF-8 nativamente.
%   fontenc[T1] se mantiene para fuentes TeX clásicas (Latin Modern).
%   Migración futura opcional: reemplazar fontenc + inconsolata por
%   \usepackage{fontspec} + \setmonofont{Inconsolata}.
%───────────────────────────────────────────────────────────────────────────────


% ─── Encoding y lenguaje ─────────────────────────────────────────────────────
% inputenc eliminado: LuaLaTeX maneja UTF-8 de forma nativa.
\usepackage[T1]{fontenc}
\usepackage[spanish,es-noshorthands]{babel}
\deactivatequoting
\usepackage{csquotes}                      % comillas correctas con babel


% ─── Tipografía ──────────────────────────────────────────────────────────────
\usepackage{microtype}
\usepackage{inconsolata}                   % fuente mono para listings


% ─── Color (→ antes de colortbl y cualquier paquete que use color) ───────────
\usepackage{xcolor}                        % carga explícita: evita conflictos
                                           % de opciones cuando tcolorbox o
                                           % tikz lo cargan implícitamente


% ─── Layout de página (→ antes de fancyhdr) ──────────────────────────────────
\usepackage{geometry}                      % (config en config/geometry.tex)


% ─── Matemáticas ─────────────────────────────────────────────────────────────
\usepackage{amsmath, amssymb, amsthm, mathtools}
\usepackage{bm}                            % bold math (\bm)
\usepackage{mathrsfs}                      % \mathscr: ℋ, ℱ, ℒ (distinto de \mathcal)
\usepackage{physics}                       % \abs, \norm, \tr, \grad, \curl,
                                           % \laplacian, \dv, \pdv, \bra, \ket
                                           % NOTA: redefine \div — ver notation.tex
\usepackage{cancel}                        % cancelar términos en derivaciones
\usepackage{mleftright}                    % \mleft \mright: corrige espaciado de
                                           % delimitadores en ciertos contextos
\usepackage{nicematrix}                    % matrices con coloreado, bloques y
                                           % notación avanzada (ML, álgebra lineal)
% \usepackage{tensor}                      % notación indicial (descomentar si se usa)
\usepackage{stmaryrd}                      % \llbracket, \rrbracket, etc.


% ─── Unidades y números ──────────────────────────────────────────────────────
\usepackage{siunitx}
\sisetup{
  output-decimal-marker = {,},            % separador decimal = coma (español)
  per-mode              = symbol,         % m/s en lugar de m·s⁻¹
  separate-uncertainty  = true,           % valor ± incertidumbre
  group-minimum-digits  = 4,             % 1000 sin separador; 10 000 con
}


% ─── Tablas (→ después de xcolor) ────────────────────────────────────────────
\usepackage{booktabs}                      % \toprule, \midrule, \bottomrule
\usepackage{array}                         % tipos de columna extendidos
\usepackage{tabularx}                      % columnas X de ancho proporcional
\usepackage{longtable}                     % tablas multipágina
\usepackage{multirow}                      % celdas multifila
\usepackage{colortbl}                      % \rowcolor, \cellcolor, \arrayrulecolor
\usepackage{pgfplotstable}                 % carga datos desde .csv a tablas y plots


% ─── Figuras y flotantes ─────────────────────────────────────────────────────
\usepackage{graphicx}                      % \includegraphics
\usepackage{caption}                       % (config en config/captions.tex)
\usepackage{subcaption}                    % subfiguras
\usepackage{float}                         % posicionamiento [H]
\usepackage{wrapfig}                       % figuras embebidas en texto
\usepackage{placeins}                      % \FloatBarrier


% ─── Listas ──────────────────────────────────────────────────────────────────
\usepackage{enumitem}                      % control de listas
\usepackage{pifont}                        % símbolos Zapf Dingbats (\ding)


% ─── Código fuente ───────────────────────────────────────────────────────────
\usepackage{listings}                      % (config en config/listings.tex)


% ─── Algoritmos ──────────────────────────────────────────────────────────────
\usepackage[ruled, vlined, linesnumbered]{algorithm2e}
                                           % (config en config/algorithms.tex)


% ─── Cajas decorativas ───────────────────────────────────────────────────────
% [many] incluye: theorems, breakable, skins, hooks, listings, xparse y más.
% No se necesita \tcbuselibrary adicional.
\usepackage[many]{tcolorbox}               % (config en config/tcolorbox.tex)


% ─── TikZ y diagramas ────────────────────────────────────────────────────────
% pgffor NO se carga explícitamente — tikz lo carga como dependencia interna.
\usepackage{tikz}
\usetikzlibrary{
  arrows.meta,
  positioning,
  shapes.geometric,
  fit,
  backgrounds,
  automata,
  calc,
  decorations.pathreplacing,
  patterns,
  matrix,
}
\usepackage{tikz-cd}                       % diagramas conmutativos
\usepackage{pgfplots}                      % (config en config/pgfplots.tex)
\pgfplotsset{compat=1.18}                  % declaración de versión: va aquí,
                                           % inmediatamente después de cargar
\usepackage{qtree}                         % árboles sintácticos simples
\usepackage{forest}                        % árboles con mayor control


% ─── Espaciado y layout ──────────────────────────────────────────────────────
\usepackage{parskip}                       % párrafos sin sangría, con espacio
\usepackage{setspace}                      % \onehalfspacing, \doublespacing
\usepackage{fancyhdr}                      % encabezados y pies de página
\usepackage{titlesec}                      % (config en config/titles.tex)


% ─── Bibliografía ────────────────────────────────────────────────────────────
\usepackage[
  backend    = biber,
  style      = numeric,                    % [1], [2], ... por orden de aparición
  sorting    = none,
  refsection = chapter,                    % bibliografía independiente por capítulo
  natbib     = true,                       % compatibilidad \citet, \citep
  hyperref   = true,
  backref    = true,                       % muestra página de cita
]{biblatex}


% ─── Notas al pie ────────────────────────────────────────────────────────────
\usepackage[bottom, hang]{footmisc}


% ─── Estructura del documento ────────────────────────────────────────────────
\usepackage{appendix}
\usepackage{pdfpages}
\usepackage[absolute,overlay]{textpos}  % posicionamiento absoluto en página
                                        % absolute: coords relativas a la esquina
                                        % superior izquierda del papel
                                        % overlay:  no afecta el flujo normal


% ─── Programación LaTeX ──────────────────────────────────────────────────────
\usepackage{etoolbox}
% xparse NO se carga: integrado al kernel LaTeX ≥ 2020.


% ─── Hipervínculos (→ PENÚLTIMO) ─────────────────────────────────────────────
\usepackage{hyperref}
\usepackage{bookmark}                      % (→ después de hyperref)
                                           % (config en config/hyperref.tex)


% ─── Referencias inteligentes (→ ÚLTIMO) ─────────────────────────────────────
\usepackage[spanish, noabbrev]{cleveref}

% Flotantes estándar
\crefname{equation}{ec.}{ecs.}
\Crefname{equation}{Ecuación}{Ecuaciones}
\crefname{figure}{fig.}{figs.}
\Crefname{figure}{Figura}{Figuras}
\crefname{table}{tabla}{tablas}
\Crefname{table}{Tabla}{Tablas}
\crefname{chapter}{cap.}{caps.}
\Crefname{chapter}{Capítulo}{Capítulos}
\crefname{section}{sec.}{secs.}
\Crefname{section}{Sección}{Secciones}
\crefname{algorithm}{alg.}{algs.}
\Crefname{algorithm}{Algoritmo}{Algoritmos}
\crefname{appendix}{ap.}{aps.}
\Crefname{appendix}{Apéndice}{Apéndices}
\crefname{listing}{cód.}{códs.}
\Crefname{listing}{Código}{Códigos}

% Entornos teoremáticos propios (definidos en environments/instances.tex)
% El contador se llama {nombre}counter por convención de generator.tex
\crefname{definicioncounter}{def.}{defs.}
\Crefname{definicioncounter}{Definición}{Definiciones}
\crefname{teoremacounter}{teo.}{teos.}
\Crefname{teoremacounter}{Teorema}{Teoremas}
\crefname{proposicioncounter}{prop.}{props.}
\Crefname{proposicioncounter}{Proposición}{Proposiciones}
\crefname{ejemplocounter}{ej.}{ejs.}
\Crefname{ejemplocounter}{Ejemplo}{Ejemplos}
\crefname{estrellacount}{★}{★}
\Crefname{estrellacount}{Estrella Polar}{Estrellas Polares}
\crefname{ejerciciocompcounter}{ej. comp.}{ejs. comp.}
\Crefname{ejerciciocompcounter}{Ejercicio Computacional}{Ejercicios Computacionales}


% ─── 2. Metadatos ────────────────────────────────────────────────────────────
% Variables del documento (\docTitulo, \docAutor, etc.).
% Va antes del tema para que hyperref las use en los metadatos del PDF.
%───────────────────────────────────────────────────────────────────────────────
% METADATA.TEX — Variables del documento
%───────────────────────────────────────────────────────────────────────────────
% Cargado en el paso 2 del preámbulo (main.tex), antes del tema.
% Consumido por: titlepage.tex, preliminary.tex, hyperref.tex.
%
% REGLA: solo \newcommand con valores de texto plano.
% No cargar paquetes ni definir macros matemáticas aquí.
%───────────────────────────────────────────────────────────────────────────────

% ─── Identificación ──────────────────────────────────────────────────────────
\newcommand{\docTitulo}{Título del documento}
\newcommand{\docSubtitulo}{Subtítulo del documento}
\newcommand{\docAutor}{Autor del documento}
\newcommand{\docFecha}{\the\year}        % año de compilación por defecto
\newcommand{\docVersion}{0.1}            % versión semántica del documento
\newcommand{\docEstado}{Borrador}        % Borrador | Versión preliminar | Final

% ─── Portada ─────────────────────────────────────────────────────────────────
\newcommand{\docCita}{``Cita representativa del documento.''}

% ─── Páginas preliminares ────────────────────────────────────────────────────
\newcommand{\docDedicatoria}{A quienes hicieron posible este trabajo.}
\newcommand{\docEpigrafe}{``Epígrafe o cita inicial.''}
\newcommand{\docAgradecimientos}{%
  Quiero expresar mi gratitud a las personas e instituciones
  que apoyaron esta investigación.%
}

% ─── Metadatos PDF (usados por hyperref en config/hyperref.tex) ──────────────
\newcommand{\docKeywords}{palabra clave 1, palabra clave 2}
\newcommand{\docSubject}{Área temática principal}
\newcommand{\docLang}{es}                % código ISO 639-1 del idioma
% \newcommand{\docDOI}{10.xxxx/xxxxx}   % descomentar si el doc tiene DOI

% ─── Recursos bibliográficos ─────────────────────────────────────────────────
%\addbibresource{references/references.bib}
\addbibresource{references/chapter1.bib}
\addbibresource{references/chapter2.bib}
\addbibresource{references/chapter3.bib}
\addbibresource{references/chapter4.bib}
\addbibresource{references/chapter5.bib}
\addbibresource{references/chapter6.bib}
\addbibresource{references/chapter7.bib}
\addbibresource{references/chapter8.bib}

% ─── 3. Tema de color ────────────────────────────────────────────────────────
% Define todos los aliases semánticos (env-*, doc-*, code-*, etc.).
% Descomentar UNO solo en catppuccin-palette.tex, no aquí.
%───────────────────────────────────────────────────────────────────────────────
% CATPPUCCIN-PALETTE.TEX — Sistema de color semántico para documentos académicos
%───────────────────────────────────────────────────────────────────────────────
%
% ARQUITECTURA DE TRES CAPAS:
%
%   Nivel 0 — Ctp* (valores absolutos del flavor, SOLO en este archivo)
%   Nivel 1 — Estructural (roles de UI genéricos, sin prefijo de dominio)
%              surface0, overlay0, maintext, background, ...
%   Nivel 2 — Semántico (roles de dominio del documento, con prefijo)
%              doc-*, link-*, env-*, eq-*, code-*, alg-*, plot-*, ...
%
% CONVENCIÓN DE NOMBRES: {dominio}-{elemento}[-{variante}]
%   Dominios activos:
%     doc      Estructura del documento (capítulos, secciones)
%     link     Hipervínculos (internos, externos, citas)
%     env      Entornos teoremáticos (definición, teorema, ...)
%     eq       Matemáticas semánticas (operadores, incógnitas, ...)
%     code     Código fuente / listings
%     alg      Algoritmos (algorithm2e)
%     plot     Gráficas pgfplots
%     tikz     Elementos gráficos TikZ
%     table    Tablas
%     caption  Captions de flotantes
%     fb       Feedback y estado (éxito, error, dificultad)
%     diff     Diffs y comparaciones
%     rainbow  Secuencia de colores distinguibles
%
% PROTOCOLO DE EXTENSIÓN — para añadir un dominio nuevo:
%   1. Añadir sección aquí con los \colorlet del dominio.
%   2. Crear config/{dominio}-config.tex consumiendo solo aliases.
%   3. Registrar \input en main.tex.
%   Nunca usar Ctp* fuera de este archivo.
%
% VARIANTES DE FONDO: usar \createback{alias} para generar {alias}-back,
%   el fondo claro para tcolorbox según el tema activo.
%───────────────────────────────────────────────────────────────────────────────

% ─── 1. Selección del flavor (descomentar solo una) ──────────────────────────
%\usepackage[Latte,styleAll]{catppuccinpalette}
%\usepackage[Frappe,styleAll]{catppuccinpalette}
%\usepackage[Macchiato,styleAll]{catppuccinpalette}
\usepackage[Mocha,styleAll]{catppuccinpalette}

% ─── 2. Toggle claro/oscuro ──────────────────────────────────────────────────
% Latte          → \darkthemefalse
% Frappé/Macchiato/Mocha → \darkthemetrue
\newif\ifdarktheme
\darkthemetrue

% ─── 3. Color de página y texto base ────────────────────────────────────────
\ifdarktheme
  \pagecolor{CtpBase}
  \AtBeginDocument{\color{CtpText}}
\else
  \pagecolor{white}
  \AtBeginDocument{\color{black}}
\fi

% ─── 4. Macro de variante de fondo ──────────────────────────────────────────
% \createback{alias}  →  define {alias}-back como fondo claro para tcolorbox.
% El resultado adapta luminosidad según el tema activo.
\newcommand{\createback}[1]{%
  \ifdarktheme
    \colorlet{#1-back}{CtpBase!95!#1}%
  \else
    \colorlet{#1-back}{white!97!#1}%
  \fi
}

% =============================================================================
% NIVEL 1 — INFRAESTRUCTURA VISUAL
% =============================================================================
% Mapeo directo del Catppuccin Style Guide: Background Colors + Typography.
% Son el vocabulario visual base — no tienen semántica de dominio.

% ─── Fondos y superficies ────────────────────────────────────────────────────
\colorlet{background}{CtpBase}
\colorlet{secondarybg}{CtpMantle}
\colorlet{crustbg}{CtpCrust}
\colorlet{surface0}{CtpSurface0}
\colorlet{surface1}{CtpSurface1}
\colorlet{surface2}{CtpSurface2}
\colorlet{overlay0}{CtpOverlay0}
\colorlet{overlay1}{CtpOverlay1}
\colorlet{overlay2}{CtpOverlay2}

% ─── Tipografía base ─────────────────────────────────────────────────────────
\colorlet{maintext}{CtpText}
\colorlet{subtext0}{CtpSubtext0}
\colorlet{subtext1}{CtpSubtext1}
\colorlet{subtletext}{CtpOverlay1}

% =============================================================================
% NIVEL 2 — SEMÁNTICA DE DOMINIO
% =============================================================================

% ─── doc: Estructura del documento ──────────────────────────────────────────
% Jerarquía de títulos: contraste decreciente con la profundidad.
% Fundamento style guide: Rainbow Highlights para jerarquía visual.
\colorlet{doc-chapter}{CtpPink}           % máximo contraste, warm
\colorlet{doc-section}{CtpSky}            % cool, segundo nivel
\colorlet{doc-subsection}{CtpLavender}    % suave, tercer nivel
\colorlet{doc-subsubsection}{CtpPeach}    % warm suave, cuarto nivel
\colorlet{doc-paragraph}{CtpSubtext0}     % casi texto, diferencia mínima

% ─── link: Hipervínculos ─────────────────────────────────────────────────────
% Fundamento style guide: Links=Blue, Followed Links=Lavender, Hover=Sky.
\colorlet{link-internal}{CtpBlue}         % refs internas (\ref, \cref, TOC)
\colorlet{link-external}{CtpPeach}        % URLs externas
\colorlet{link-cite}{CtpGreen}            % citas bibliográficas
\colorlet{link-visited}{CtpLavender}      % Followed Links
\colorlet{link-hover}{CtpSky}             % On Hover Links

% ─── env: Entornos teoremáticos ──────────────────────────────────────────────
% Cada entorno tiene frame (color pleno) y back (variante clara, vía \createback).
\colorlet{env-def}{CtpBlue}               % Definición: epistémico, fundacional
\colorlet{env-thm}{CtpMauve}              % Teorema: color más prominente
\colorlet{env-prop}{CtpGreen}             % Proposición: verdad demostrada
\colorlet{env-ex}{CtpPeach}               % Ejemplo: warm, invita a leer
\colorlet{env-obs}{CtpOverlay0}           % Observación: secundaria, no compite
\colorlet{env-phys}{CtpSapphire}          % Conexión física: evoca ciencia
\colorlet{env-star}{CtpYellow}            % Estrella Polar: destaca
\colorlet{env-comp}{CtpPink}              % Ejercicio computacional: activo
\createback{env-def}
\createback{env-thm}
\createback{env-prop}
\createback{env-ex}
\createback{env-obs}
\createback{env-phys}
\createback{env-star}
\createback{env-comp}

% ─── eq: Matemáticas semánticas ──────────────────────────────────────────────
% Coloreado en tres ejes ortogonales. Las macros viven en equation-styles.tex.
% Fundamento: extensión del style guide al dominio matemático.
%
% Eje 1 — Rol sintáctico (estable entre contextos)
\colorlet{eq-op}{CtpBlue}                 % operadores: ∫, ∑, ∂, ∇, Ĥ
\colorlet{eq-fn}{CtpGreen}                % funciones/mapas: f(·), σ(·), V(r)
\colorlet{eq-dm}{CtpTeal}                 % diferenciales/medidas: dx, dt, dW
\colorlet{eq-dom}{CtpSapphire}            % dominios/conjuntos: Ω, ∂Ω, ℝⁿ
%
% Eje 2 — Estatus epistémico (la dimensión que la tipografía no codifica)
\colorlet{eq-unknown}{CtpPeach}           % incógnita / buscado / aprendido
\colorlet{eq-known}{CtpMauve}             % dato / conocido / fijado
\colorlet{eq-const}{CtpLavender}          % constante universal: ℏ, kB, π, c
%
% Eje 3 — Tipo matemático (refuerzo sutil de la tipografía)
\colorlet{eq-vec}{CtpRed}                 % vectores/campos: u, E, B, J
\colorlet{eq-ten}{CtpRosewater}           % tensores/matrices: σᵢⱼ, gμν, W

% ─── code: Código fuente (listings) ─────────────────────────────────────────
% Mapeo directo del Catppuccin Style Guide: Language Defaults (Code Editors).
\colorlet{code-keyword}{CtpMauve}         % Keyword
\colorlet{code-string}{CtpGreen}          % Strings
\colorlet{code-comment}{CtpOverlay2}      % Comments
\colorlet{code-operator}{CtpSky}          % Operators
\colorlet{code-type}{CtpYellow}           % Classes, Interfaces, Types
\colorlet{code-fn}{CtpBlue}               % Methods, Functions
\colorlet{code-mod}{CtpTeal}              % Enum Variants → módulos/namespaces
\colorlet{code-macro}{CtpRosewater}       % Macros, decoradores, pragmas
\colorlet{code-builtin}{CtpRed}           % Builtins, Atoms, Symbols
\colorlet{code-const}{CtpPeach}           % Constants, Numbers
\colorlet{code-param}{CtpMaroon}          % Parameters
\colorlet{code-delim}{CtpOverlay2}        % Braces, Delimiters
\colorlet{code-bg}{CtpBase}               % fondo del bloque de código
\colorlet{code-border}{CtpSurface0}       % frame / rule
\colorlet{code-linenum}{CtpOverlay1}      % Line Numbers

% ─── alg: Algoritmos (algorithm2e) ──────────────────────────────────────────
% Comparte filosofía con code, pero con roles propios del pseudocódigo.
\colorlet{alg-keyword}{CtpMauve}          % palabras clave: While, For, Return
\colorlet{alg-arg}{CtpText}               % argumentos: texto principal
\colorlet{alg-comment}{CtpOverlay2}       % comentarios inline
\colorlet{alg-linenum}{CtpSky}            % números de línea
\colorlet{alg-caption}{CtpBlue}           % caption: ancla visual del algoritmo
\colorlet{alg-line}{CtpOverlay0}          % líneas estructurales: sutiles

% ─── plot: Gráficas (pgfplots) ───────────────────────────────────────────────
% Series: Rainbow Highlights del style guide + extensiones para 8 curvas.
% Ejes e infraestructura usan aliases estructurales de nivel 1.
\colorlet{plot-series-1}{CtpBlue}         % Rainbow #5 (reordenado por contraste)
\colorlet{plot-series-2}{CtpPeach}        % Rainbow #2
\colorlet{plot-series-3}{CtpGreen}        % Rainbow #4
\colorlet{plot-series-4}{CtpMauve}        % extensión
\colorlet{plot-series-5}{CtpRed}          % Rainbow #1
\colorlet{plot-series-6}{CtpTeal}         % extensión
\colorlet{plot-series-7}{CtpYellow}       % Rainbow #3
\colorlet{plot-series-8}{CtpPink}         % extensión
\colorlet{plot-axis}{CtpOverlay0}         % ejes y ticks
\colorlet{plot-grid}{CtpSurface1}         % grid: más sutil que plot-axis
\colorlet{plot-label}{CtpText}            % labels de ejes y títulos
\colorlet{plot-tick-label}{CtpSubtext0}   % números de ticks
\colorlet{plot-legend-bg}{CtpBase}        % fondo de leyenda
\colorlet{plot-legend-border}{CtpSurface0}% borde de leyenda
\colorlet{plot-errorbar}{CtpOverlay0}     % barras de error: sutiles

% ─── tikz: Elementos gráficos TikZ ──────────────────────────────────────────
% Valores por defecto para nodos, aristas y estados de autómata.
% El contexto semántico puede sobreescribir directamente.
\colorlet{tikz-node-bg}{CtpSurface0}      % nodo genérico, fondo
\colorlet{tikz-node-border}{CtpOverlay0}  % nodo genérico, borde
\colorlet{tikz-edge}{CtpOverlay1}         % arista estándar
\colorlet{tikz-arrow}{CtpOverlay2}        % punta de flecha
\colorlet{tikz-label}{CtpSubtext0}        % anotación/label
\colorlet{tikz-accept}{CtpGreen}          % autómata: estado aceptador
\colorlet{tikz-reject}{CtpRed}            % autómata: estado rechazador
\colorlet{tikz-initial}{CtpBlue}          % autómata: estado inicial
\colorlet{tikz-highlight}{CtpYellow}      % nodo destacado en explicación

% ─── table: Tablas ───────────────────────────────────────────────────────────
\colorlet{table-head}{CtpBase}         % fondo encabezado: azul muy sutil
\colorlet{table-head-text}{CtpText}       % texto del encabezado
\colorlet{table-row-alt}{CtpSurface0}  % filas alternadas
\colorlet{table-rule}{CtpOverlay0!50}     % líneas booktabs

% ─── caption: Captions de flotantes ─────────────────────────────────────────
% Cada tipo de flotante tiene identidad visual propia.
% Los labels de caption ecoízan el color del entorno correspondiente.
\colorlet{caption-fig}{CtpBlue}           % label figura: = link-internal
\colorlet{caption-tab}{CtpMauve}          % label tabla: = env-thm
\colorlet{caption-code}{CtpGreen}         % label código: = code-string
\colorlet{caption-alg}{CtpPeach}          % label algoritmo: cálido
\colorlet{caption-sub}{CtpSapphire}       % sub-caption: = env-phys
\colorlet{caption-text}{CtpSubtext0}      % texto descriptivo del caption

% ─── fb: Feedback y estado ───────────────────────────────────────────────────
% Fundamento style guide: Success=Green, Warnings=Yellow, Errors=Red, Info=Teal.
\colorlet{fb-success}{CtpGreen}
\colorlet{fb-warning}{CtpYellow}
\colorlet{fb-error}{CtpRed}
\colorlet{fb-info}{CtpTeal}
\colorlet{fb-difficulty}{CtpLavender}     % indicador de dificultad: neutro
\createback{fb-success}
\createback{fb-warning}
\createback{fb-error}
\createback{fb-info}

% ─── diff: Diffs y comparaciones ─────────────────────────────────────────────
% Fundamento style guide: Diff & Merge.
\colorlet{diff-header}{CtpBlue}
\colorlet{diff-index}{CtpOverlay2}
\colorlet{diff-file}{CtpPink}
\colorlet{diff-hunk}{CtpPeach}
\colorlet{diff-changed-bg}{CtpBlue!15}
\colorlet{diff-inserted-bg}{CtpGreen!15}
\colorlet{diff-removed-bg}{CtpRed!15}
\colorlet{diff-breakpoint}{CtpRed}

% ─── rainbow: Secuencia de colores distinguibles ─────────────────────────────
% Fundamento style guide: Rainbow Highlights (Brackets, Headings, etc.).
\colorlet{rainbow-1}{CtpRed}
\colorlet{rainbow-2}{CtpPeach}
\colorlet{rainbow-3}{CtpYellow}
\colorlet{rainbow-4}{CtpGreen}
\colorlet{rainbow-5}{CtpSapphire}
\colorlet{rainbow-6}{CtpLavender}

% =============================================================================
% ATAJOS TIPOGRÁFICOS
% =============================================================================
\newcommand{\textsuccess}[1]{\textcolor{fb-success}{#1}}
\newcommand{\textwarning}[1]{\textcolor{fb-warning}{#1}}
\newcommand{\texterror}[1]{\textcolor{fb-error}{#1}}
\newcommand{\textinfo}[1]{\textcolor{fb-info}{#1}}
\newcommand{\textkeyword}[1]{\textcolor{code-keyword}{#1}}
\newcommand{\textstring}[1]{\textcolor{code-string}{#1}}
\newcommand{\textcomment}[1]{\textcolor{code-comment}{#1}}

% ─── frontmatter: Páginas preliminares ───────────────────────────────────────
% Alias propios de portada y páginas preliminares.
% No comparten dominio con entornos teoremáticos ni estructura del documento.
\colorlet{titlepage-accent}{CtpMauve}    % línea decorativa y borde de caja de cita
                                          % Mauve: prominente, no reclamado por
                                          % ningún otro rol estructural del doc.


% ─── 4. Hipervínculos ────────────────────────────────────────────────────────
% Después del tema: consume aliases link-* de la paleta.
%───────────────────────────────────────────────────────────────────────────────
% HYPERREF.TEX — Configuración de hipervínculos
%───────────────────────────────────────────────────────────────────────────────
% Cargado DESPUÉS del tema: consume aliases link-* de catppuccin-palette.tex.
% hyperref y bookmark se cargan en packages.tex.
% No usa Ctp* directamente.
%───────────────────────────────────────────────────────────────────────────────

\hypersetup{
  colorlinks  = true,
  linkcolor   = link-internal,   % referencias internas (\ref, \cref, TOC)
  citecolor   = link-cite,       % citas bibliográficas
  urlcolor    = link-external,   % URLs externas
  pdfborder   = {0 0 0},
  breaklinks  = true,
  % ── Metadatos PDF (macros definidas en metadata.tex) ──────────────────────
  pdftitle    = {\docTitulo},
  pdfauthor   = {\docAutor},
  pdfsubject  = {\docSubject},
  pdfkeywords = {\docKeywords},
  pdfcreator  = {LaTeX + Catppuccin template},
}


% ─── 5. Layout de página ─────────────────────────────────────────────────────
% Después de hyperref: fancyhdr puede interactuar con el layout de hyperref.
%───────────────────────────────────────────────────────────────────────────────
% GEOMETRY.TEX — Layout de página, espaciado y encabezados
%───────────────────────────────────────────────────────────────────────────────
% \usepackage{geometry} declarado en packages.tex.
% Paquetes requeridos: geometry, setspace, parskip, fancyhdr.
%───────────────────────────────────────────────────────────────────────────────

% ─── Márgenes ────────────────────────────────────────────────────────────────
% Margen izquierdo más amplio: reserva espacio para encuadernación.
% marginpar: espacio para notas al margen (\\marginpar{}).
\geometry{
  top           = 2.5cm,
  bottom        = 2.5cm,
  left          = 3.5cm,
  right         = 3.0cm,
  headheight    = 15pt,
  headsep       = 20pt,
  footskip      = 25pt,
  marginparwidth= 2cm,
  marginparsep  = 0.5cm,
}

% ─── Espaciado de párrafos ───────────────────────────────────────────────────
\setlength{\parskip}{6pt plus 2pt minus 2pt}
\onehalfspacing

% ─── Encabezados y pies de página ────────────────────────────────────────────
% Estilo principal: sección en izquierda, número de página en derecha.
\pagestyle{fancy}
\fancyhf{}
\fancyhead[L]{\nouppercase{\leftmark}}
\fancyhead[R]{\thepage}
\renewcommand{\headrulewidth}{0.4pt}
\addtolength{\headheight}{3pt}

% Estilo plain: para páginas de inicio de capítulo.
% Solo número de página centrado, sin línea de encabezado.
\fancypagestyle{plain}{%
  \fancyhf{}
  \renewcommand{\headrulewidth}{0pt}
  \fancyfoot[C]{\thepage}
}


% ─── 6. Cajas base ───────────────────────────────────────────────────────────
% Antes de environments/: generator.tex depende de basebox e innerbox.
%───────────────────────────────────────────────────────────────────────────────
% TCOLORBOX.TEX — Estilo base y entornos auxiliares para cajas
%───────────────────────────────────────────────────────────────────────────────
% Define basebox: estilo heredado por todos los entornos de environments/.
% Define boxfig y boxtab: versiones no-flotantes para usar DENTRO de cajas.
%
% PROBLEMA FUNDAMENTAL — flotantes dentro de tcolorbox:
%   tcolorbox construye internamente un \vbox. Cualquier entorno flotante
%   (\begin{figure}, \begin{table}, \begin{acadtable}) invoca \@float, que
%   requiere "outer par mode". Dentro de un \vbox ese modo no existe →
%   error "Not in outer par mode".
%
%   Entornos afectados dentro de cualquier caja basada en basebox:
%     \begin{figure}      → usar \begin{boxfig}
%     \begin{table}       → usar \begin{boxtab}
%     \begin{acadtable}   → usar \begin{boxtab}
%     \begin{lstlisting}[float=...] → quitar la opción float=
%     \begin{longtable}   → incompatible; usar tabularx sin float
%
%   La solución usa \captionof{tipo}{texto} del paquete caption (ya cargado),
%   que genera el caption y actualiza el contador sin invocar el mecanismo
%   de flotante. \phantomsection corrige el anchor del PDF con hyperref.
%
% ANCHO DENTRO DE CAJAS:
%   Dentro de tcolorbox, \linewidth = ancho del área de contenido de la caja
%   (descontando left=6pt + right=6pt). Usar siempre \linewidth, no \textwidth,
%   para tabularx, \includegraphics[width=...], etc.
%───────────────────────────────────────────────────────────────────────────────


% ─── Estilo base: heredado por todos los entornos teoremáticos ───────────────
\tcbset{
  basebox/.style={
    enhanced,
    breakable,
    colback    = surface0,
    colframe   = overlay0!30,
    coltext    = maintext,
    arc        = 0.5pt,
    outer arc  = 0.5pt,
    boxrule    = 0.5pt,
    left       = 6pt,
    right      = 6pt,
    top        = 8pt,
    bottom     = 8pt,
    drop fuzzy shadow = background,
    fonttitle  = \bfseries,
    before title = {\vspace{2pt}},
    after title  = {\vspace{2pt}},
  }
}


% ─── boxfig: figura no-flotante para usar dentro de cajas ────────────────────
% Uso:
%   \begin{boxfig}{fig:label}{Caption de la figura}
%     \includegraphics[width=\linewidth]{ruta/imagen}
%   \end{boxfig}
%
% Nota: el label va DESPUÉS de \captionof para que hyperref ancle
% correctamente. \phantomsection asegura el anchor del PDF.
\NewDocumentEnvironment{boxfig}{ m m }{%
  \begin{center}
}{%
  \phantomsection
  \captionof{figure}{#2}%
  \label{#1}%
  \end{center}
}


% ─── boxtab: tabla no-flotante para usar dentro de cajas ─────────────────────
% Reemplaza \begin{acadtable} y \begin{table} dentro de tcolorbox.
%
% Uso:
%   \begin{boxtab}{tab:label}{Caption de la tabla}
%     \begin{tabularx}{\linewidth}{l Y Z}
%       ...
%     \end{tabularx}
%   \end{boxtab}
%
% Nota: usar \linewidth (no \textwidth) para tabularx dentro de la caja.
% La posición del caption sigue la convención de tablas: arriba del contenido.
\NewDocumentEnvironment{boxtab}{ m m }{%
  \begin{center}
  \phantomsection
  \captionof{table}{#2}%
  \label{#1}%
}{%
  \end{center}
}


% ─── Entornos anidados: cajas dentro de cajas ────────────────────────────────
% tcolorbox permite anidar cajas solo si la caja interna tiene breakable=false.
% Estilo para cajas secundarias dentro de un entorno teoremático:
\tcbset{
  innerbox/.style={
    basebox,
    breakable = false,        % obligatorio para anidar dentro de basebox
    boxrule   = 0.3pt,
    left      = 4pt,
    right     = 4pt,
    top       = 4pt,
    bottom    = 4pt,
    drop fuzzy shadow = none, % sin sombra en cajas internas
  }
}


% ─── 7. Entornos del documento ───────────────────────────────────────────────
% Orden obligatorio: generator → instances → specials.
%───────────────────────────────────────────────────────────────────────────────
% GENERATOR.TEX — Generador de entornos teoremáticos
%───────────────────────────────────────────────────────────────────────────────
% Uso: \crearEntorno{nombre}{color}{título}{conContador?}{opcionesAdicionales}
%
%   #1  nombre del entorno        (ej: definicion)
%   #2  alias semántico de color  (ej: env-def)
%   #3  título base               (ej: Definición)
%   #4  1 si tiene contador, 0 si no
%   #5  opciones adicionales para tcolorbox (puede estar vacío)
%
% Requiere que {#2}-back esté definido en catppuccin-palette.tex vía \createback.
% No usa Ctp* directamente.
%───────────────────────────────────────────────────────────────────────────────
\newcommand{\crearEntorno}[5]{%
  \ifnum#4=1
    \newcounter{#1counter}[chapter]
    \expandafter\renewcommand\csname the#1counter\endcsname{%
      \thechapter.\arabic{#1counter}%
    }
  \fi

  \NewDocumentEnvironment{#1}{ m }{
    \ifnum#4=1
      \refstepcounter{#1counter}%
    \fi
    \begin{tcolorbox}[
      basebox,
      colframe=#2,
      colback=#2-back,
      coltitle = #2-back,
      title={%
        \ifnum#4=1
          \unexpanded{#3 }\csname the#1counter\endcsname\unexpanded{ --- }##1%
        \else
          \unexpanded{#3 --- }##1%
        \fi
      },
      #5
    ]
  }{
    \end{tcolorbox}
  }
}

%───────────────────────────────────────────────────────────────────────────────
% INSTANCES.TEX — Instancias de entornos teoremáticos
%───────────────────────────────────────────────────────────────────────────────
% Consume aliases del dominio env-* definidos en catppuccin-palette.tex.
% Los fondos {env-*-back} se generan automáticamente vía \createback en palette.
% No usa Ctp* directamente.
%───────────────────────────────────────────────────────────────────────────────

% ─── Nivel 1: Fundamentos (con contador) ─────────────────────────────────────
\crearEntorno{definicion}{env-def}{Definición}{1}{}
\crearEntorno{teorema}{env-thm}{Teorema}{1}{}
\crearEntorno{proposicion}{env-prop}{Proposición}{1}{}

% ─── Nivel 2: Apoyo ──────────────────────────────────────────────────────────
\crearEntorno{ejemplo}{env-ex}{Ejemplo}{1}{}
\crearEntorno{observacion}{env-obs}{Observación}{0}{}
\crearEntorno{fisicaconexion}{env-phys}{[F] Conexión física}{0}{%
  fonttitle=\bfseries\small, before upper={\small}}

% ─── Nivel 3: Especiales (con contador) ──────────────────────────────────────
\crearEntorno{estrella}{env-star}{\texorpdfstring{$\bigstar$}{★} Estrella Polar}{1}{%
  fonttitle=\bfseries\small, before upper={\small}}
\crearEntorno{ejerciciocomp}{env-comp}{$\triangleright$ Ejercicio Computacional}{1}{%
  fonttitle=\bfseries\small, before upper={\small}}

%───────────────────────────────────────────────────────────────────────────────
% SPECIALS.TEX — Entornos especiales: dificultad, notas bib, ejercicios
%───────────────────────────────────────────────────────────────────────────────
% Consume aliases fb-* y aliases estructurales de catppuccin-palette.tex.
% No usa Ctp* directamente.
%───────────────────────────────────────────────────────────────────────────────

% ─── Indicador de dificultad ─────────────────────────────────────────────────
\newcommand{\estrellita}[1]{%
  \foreach \i in {1,...,5}{%
    \ifnum\i>#1
      {\small\ding{73}}% estrella vacía (☆)
    \else
      {\small\ding{72}}% estrella rellena (★)
    \fi
    \ifnum\i<5\thinspace\fi
  }%
}

\newcommand{\dificultad}[2]{%
  \begin{tcolorbox}[
    enhanced,
    breakable,
    colframe=fb-difficulty,
    colback=background,
    coltext=maintext,
    arc=0.5pt,
    outer arc=0.5pt,
    boxrule=0.5pt,
    left=2pt,
    right=2pt,
    top=2pt,
    bottom=2pt,
    drop fuzzy shadow=background,
    fontupper=\small
  ]
  \textbf{Nivel:}\quad
  Matemática:~\estrellita{#1}\qquad
  Física:~\estrellita{#2}
  \end{tcolorbox}
}

% ─── Notas bibliográficas ─────────────────────────────────────────────────────
\NewDocumentEnvironment{notasbib}{}
{
  \section*{Notas bibliográficas}
  \addcontentsline{toc}{section}{\texorpdfstring{Notas bibliográficas}{Notas bibliograficas}}
  \small
  \begin{tcolorbox}[
    basebox,
    colback=surface0,
    colframe=overlay0!30,
    left=8pt,
    right=8pt,
    top=6pt,
    bottom=6pt,
  ]
}
{
  \end{tcolorbox}
  \normalsize
}

% ─── Ejercicios simples (sin caja) ───────────────────────────────────────────
\newcounter{ejercicio}[chapter]
\renewcommand{\theejercicio}{\thechapter.\arabic{ejercicio}}

\newenvironment{ejercicio}{%
  \refstepcounter{ejercicio}%
  \par\medskip\noindent
  \textbf{Ejercicio \theejercicio.}\quad
}{\par\medskip}


% ─── 8. Configuraciones de dominio ───────────────────────────────────────────
% Después del tema y de los entornos. Orden entre sí es libre.
%───────────────────────────────────────────────────────────────────────────────
% TITLES-CONFIG.TEX — Jerarquía de títulos con colores semánticos Catppuccin
%───────────────────────────────────────────────────────────────────────────────
% Consume aliases del dominio doc-* definidos en catppuccin-palette.tex.
% No usa Ctp* directamente.
% titlesec se carga en packages.tex.
%───────────────────────────────────────────────────────────────────────────────

% ─── Capítulo ────────────────────────────────────────────────────────────────
\titleformat{\chapter}[display]
  {\normalfont\Huge\bfseries\color{doc-chapter}}
  {\chaptertitlename\ \thechapter}{20pt}{\Huge\color{doc-chapter}}

% ─── Sección ─────────────────────────────────────────────────────────────────
\titleformat{\section}
  {\normalfont\Large\bfseries\color{doc-section}}
  {\thesection}{1em}{}

% ─── Subsección ──────────────────────────────────────────────────────────────
\titleformat{\subsection}
  {\normalfont\large\bfseries\color{doc-subsection}}
  {\thesubsection}{1em}{}

% ─── Subsubsección ───────────────────────────────────────────────────────────
\titleformat{\subsubsection}
  {\normalfont\normalsize\bfseries\color{doc-subsubsection}}
  {\thesubsubsection}{1em}{}

%───────────────────────────────────────────────────────────────────────────────
% CAPTIONS-CONFIG.TEX — Estilo de captions con colores semánticos Catppuccin
%───────────────────────────────────────────────────────────────────────────────
% Consume aliases del dominio caption-* definidos en catppuccin-palette.tex.
% No usa Ctp* directamente.
%───────────────────────────────────────────────────────────────────────────────

% ─── Figuras ─────────────────────────────────────────────────────────────────
\captionsetup{
  font        = small,
  labelfont   = {bf, color=caption-fig},     % "Figura 1.1" en bold + azul
  textfont    = {color=caption-text},         % texto del caption en gris suave
  format      = hang,                         % sangría francesa
  labelsep    = endash,                       % "Figura 1.1 –"
  skip        = 8pt,
  justification = justified,
}

% ─── Subfiguras ──────────────────────────────────────────────────────────────
\captionsetup[sub]{
  font        = footnotesize,
  labelfont   = {bf, color=caption-sub},     % "(a)" en sapphire
  textfont    = {color=caption-text},
  labelsep    = space,
  skip        = 4pt,
}

% ─── Tablas ──────────────────────────────────────────────────────────────────
\captionsetup[table]{
  position    = above,                        % caption arriba de la tabla
  labelfont   = {bf, color=caption-tab},     % "Tabla 1.1" en mauve
  skip        = 6pt,
}

% ─── Algoritmos ──────────────────────────────────────────────────────────────
% algorithm2e gestiona su propio caption vía \SetAlCapSty en algorithms.tex.
% No se necesita configuración aquí.

% ─── Listings (código fuente) ─────────────────────────────────────────────────
\captionsetup[lstlisting]{
  labelfont   = {bf, color=caption-code},    % "Código 1.1" en verde
  textfont    = {color=caption-text},
  format      = hang,
  labelsep    = endash,
  skip        = 6pt,
}

% ─── Nombres en español ──────────────────────────────────────────────────────
\renewcommand{\lstlistingname}{Código}
\renewcommand{\lstlistlistingname}{Índice de códigos}

%───────────────────────────────────────────────────────────────────────────────
% TABLES.TEX — Estilo global de tablas con colores semánticos Catppuccin
%───────────────────────────────────────────────────────────────────────────────
% Consume aliases del dominio table-* definidos en catppuccin-palette.tex.
% No usa Ctp* directamente ni define colorlets propios.
%
% ENTORNOS DISPONIBLES:
%   acadtable   → tabla flotante (fuera de tcolorbox)
%   boxtab      → tabla no-flotante para DENTRO de tcolorbox (en tcolorbox.tex)
%   fittable    → reescala automática si la tabla excede el ancho disponible
%───────────────────────────────────────────────────────────────────────────────


% ─── Espaciado global de booktabs ────────────────────────────────────────────
\setlength{\heavyrulewidth}{0.08em}
\setlength{\lightrulewidth}{0.04em}
\setlength{\cmidrulewidth}{0.03em}
\renewcommand{\arraystretch}{1.25}

% ─── Alineación vertical: m (middle) en lugar de p (top) ────────────────────
% tabularx usa p{} por defecto → desalineación vertical al mezclar con l/c/r.
\renewcommand{\tabularxcolumn}[1]{m{#1}}

% ─── Tipos de columna ────────────────────────────────────────────────────────
% L/C/R: ancho fijo con wrapping y centrado vertical
\newcolumntype{L}[1]{>{\raggedright\arraybackslash}m{#1}}
\newcolumntype{C}[1]{>{\centering\arraybackslash}m{#1}}
\newcolumntype{R}[1]{>{\raggedleft\arraybackslash}m{#1}}

% Y/Z: columnas proporcionales para tabularx (centrada / derecha)
% X ya existe en tabularx (izquierda, proporcional, ahora con m{})
\newcolumntype{Y}{>{\centering\arraybackslash}X}
\newcolumntype{Z}{>{\raggedleft\arraybackslash}X}

% ─── Color de reglas booktabs ─────────────────────────────────────────────────
\AtBeginDocument{\arrayrulecolor{table-rule}}

% ─── Encabezado de columna ───────────────────────────────────────────────────
\newcommand{\thead}[1]{%
  \bfseries\color{table-head-text}#1%
}

% Colorear la fila de encabezado completa vía \rowcolor
\newcommand{\headerrow}{\rowcolor{table-head}}


% ─── fittable: reescala si excede el ancho disponible ────────────────────────
% Funciona tanto en texto normal (\textwidth) como dentro de tcolorbox
% (\linewidth). Usa \linewidth en ambos contextos: fuera de una caja
% \linewidth == \textwidth; dentro, \linewidth es el ancho de la caja.
%
% Uso:
%   \begin{fittable}
%     \begin{tabular}{llll} ... \end{tabular}
%   \end{fittable}
\newsavebox{\fittablebox}
\NewDocumentEnvironment{fittable}{}{%
  \begin{lrbox}{\fittablebox}%
}{%
  \end{lrbox}%
  \ifdim\wd\fittablebox>\linewidth
    \resizebox{\linewidth}{!}{\usebox{\fittablebox}}%
  \else
    \usebox{\fittablebox}%
  \fi
}


% ─── acadtable: tabla flotante (SOLO fuera de tcolorbox) ─────────────────────
% Produce \begin{table}[placement] → flotante.
% DENTRO de un entorno tcolorbox (definicion, ejemplo, etc.) usar boxtab
% de config/tcolorbox.tex en su lugar.
%
% Uso:
%   \begin{acadtable}[htbp]{tab:label}{Caption text}
%     \begin{tabularx}{\textwidth}{l Y Z}
%       ...
%     \end{tabularx}
%   \end{acadtable}
\NewDocumentEnvironment{acadtable}{ O{htbp} m m }{%
  \begin{table}[#1]%
    \centering
    \captionof{table}{#3}\label{#2}%
}{%
  \end{table}%
}


% ─── Renombrar "Cuadro" → "Tabla" en español ─────────────────────────────────
\addto\captionsspanish{\renewcommand{\tablename}{Tabla}}


% ─── Nota al pie de tabla ────────────────────────────────────────────────────
\newcommand{\tablenote}[2]{%
  \par\vspace{2pt}%
  {\footnotesize\textsuperscript{#1}\,#2}%
}


%───────────────────────────────────────────────────────────────────────────────
% NOTA: tabularx NO puede envolverse en \NewDocumentEnvironment porque
% escanea el cuerpo literalmente buscando \end{tabularx}.
% Usar siempre directamente:
%   \begin{tabularx}{\textwidth}{l Y Z}   ← fuera de caja
%   \begin{tabularx}{\linewidth}{l Y Z}   ← dentro de tcolorbox
%───────────────────────────────────────────────────────────────────────────────

%───────────────────────────────────────────────────────────────────────────────
% LISTINGS-CATPPUCCIN.TEX — Syntax highlighting con colores semánticos Catppuccin
%───────────────────────────────────────────────────────────────────────────────
% Consume aliases del dominio code-* definidos en catppuccin-palette.tex.
% No usa Ctp* directamente.
%
% Mapa de clases de keyword → alias semántico:
%
%   Clase         Alias          Rol
%   ─────         ─────          ───
%   keywords[1]   code-keyword   Palabras clave del lenguaje
%   keywords[2]   code-builtin   Builtins / atoms / symbols
%   keywords[3]   code-const     Constantes especiales (Inf, NaN, ...)
%   keywords[4]   code-operator  Operadores como keywords (not, and, or)
%   keywords[5]   code-fn        Funciones / métodos estándar
%   keywords[6]   code-type      Tipos / clases / interfaces
%   keywords[7]   code-mod       Módulos / namespaces / enums
%   keywords[8]   code-macro     Macros / decoradores / atributos / pragmas
%
% Colores automáticos (no requieren keywords):
%   Strings        code-string
%   Comentarios    code-comment  (itálica)
%   Números línea  code-linenum
%   Fondo          code-bg       (desactivado por defecto → transparente)
%   Texto base     maintext
%   Bordes         code-border
%───────────────────────────────────────────────────────────────────────────────


% ═════════════════════════════════════════════════════════════════════════════
% ESTILO BASE
% ═════════════════════════════════════════════════════════════════════════════

\lstdefinestyle{catppuccin-base}{
  % ── Layout ─────────────────────────────────────────────────────────────────
  breaklines=true,
  breakatwhitespace=true,
  showstringspaces=false,
  tabsize=2,
  captionpos=b,
  xleftmargin=10pt,
  xrightmargin=5pt,
  aboveskip=10pt,
  belowskip=10pt,
  frame=lines,
  framesep=5pt,
  % ── Soporte UTF-8 ──────────────────────────────────────────────────────────
  extendedchars=true,
  inputencoding=utf8,
  literate=
    {á}{{\\'a}}1  {é}{{\\'e}}1  {í}{{\\'i}}1  {ó}{{\\'o}}1  {ú}{{\\'u}}1
    {Á}{{\\'A}}1  {É}{{\\'E}}1  {Í}{{\\'I}}1  {Ó}{{\\'O}}1  {Ú}{{\\'U}}1
    {ñ}{{\~n}}1   {Ñ}{{\~N}}1   {ü}{{\"u}}1   {Ü}{{\"U}}1
    {¡}{{!`}}1    {¿}{{?`}}1
    {°}{{\textdegree}}1
    {²}{{\textsuperscript{2}}}1  {³}{{\textsuperscript{3}}}1
    {×}{{\texttimes}}1  {÷}{{\textdiv}}1
    {±}{{\textpm}}1     {≠}{{$\neq$}}1
    {≤}{{$\leq$}}1      {≥}{{$\geq$}}1
    {∞}{{$\infty$}}1    {π}{{$\pi$}}1
    {→}{{$\rightarrow$}}1 {←}{{$\leftarrow$}}1
    {⟨}{{$\langle$}}1   {⟩}{{$\rangle$}}1
    {λ}{{$\lambda$}}1   {α}{{$\alpha$}}1
    {β}{{$\beta$}}1     {γ}{{$\gamma$}}1
    {δ}{{$\delta$}}1    {ε}{{$\varepsilon$}}1
    {θ}{{$\theta$}}1    {σ}{{$\sigma$}}1
    {μ}{{$\mu$}}1       {ρ}{{$\rho$}}1
    {τ}{{$\tau$}}1      {φ}{{$\varphi$}}1
    {ω}{{$\omega$}}1    {Ω}{{$\Omega$}}1
    {∈}{{$\in$}}1       {∉}{{$\notin$}}1
    {⊂}{{$\subset$}}1   {∪}{{$\cup$}}1
    {∩}{{$\cap$}}1      {∅}{{$\emptyset$}}1
    {√}{{$\sqrt{}$}}1   {∑}{{$\sum$}}1
    {∏}{{$\prod$}}1     {∫}{{$\int$}}1
    {…}{{$\ldots$}}1    {·}{{$\cdot$}}1,
  % ── Colores de infraestructura ─────────────────────────────────────────────
  rulecolor=\color{code-border},
  backgroundcolor={},                           % transparente (code-bg disponible)
  basicstyle=\ttfamily\small\color{maintext},
  numberstyle=\tiny\color{code-linenum},
  identifierstyle=\color{maintext},
  % ── Sintaxis por clase ─────────────────────────────────────────────────────
  keywordstyle=\color{code-keyword}\bfseries,   % [1] keywords
  keywordstyle=[2]\color{code-builtin},          % [2] builtins / atoms
  keywordstyle=[3]\color{code-const},            % [3] constantes especiales
  keywordstyle=[4]\color{code-operator},         % [4] operadores keyword
  keywordstyle=[5]\color{code-fn},               % [5] funciones estándar
  keywordstyle=[6]\color{code-type},             % [6] tipos / clases
  keywordstyle=[7]\color{code-mod},              % [7] módulos / enums
  keywordstyle=[8]\color{code-macro},            % [8] macros / decoradores
  stringstyle=\color{code-string},
  commentstyle=\color{code-comment}\itshape,
  % ── Strings genéricos ──────────────────────────────────────────────────────
  morestring=[b][\color{code-string}]",
  morestring=[b][\color{code-string}]',
}

% Estilo por defecto global
\lstset{style=catppuccin-base}


% ═════════════════════════════════════════════════════════════════════════════
%  PYTHON
% ═════════════════════════════════════════════════════════════════════════════

\lstdefinestyle{catppuccin-python}{
  style=catppuccin-base,
  language=Python,
  morekeywords={[1]and, as, assert, async, await, break, class, continue,
    def, del, elif, else, except, finally, for, from, global, if, import,
    in, is, lambda, nonlocal, not, or, pass, raise, return, try, while,
    with, yield},
  morekeywords={[2]True, False, None, self, cls, NotImplemented, Ellipsis,
    __all__, __name__, __main__, __init__, __repr__, __str__, __call__,
    __enter__, __exit__, __iter__, __next__, __getitem__, __setitem__,
    __len__, __eq__, __hash__, __new__, __del__},
  morekeywords={[3]POSITIVE_INFINITY, NEGATIVE_INFINITY, nan, inf},
  morekeywords={[5]abs, all, any, bin, bool, breakpoint, callable, chr,
    compile, complex, delattr, dict, dir, divmod, enumerate, eval, exec,
    filter, float, format, frozenset, getattr, globals, hasattr, hash, help,
    hex, id, input, int, isinstance, issubclass, iter, len, list, locals,
    map, max, memoryview, min, next, object, oct, open, ord, pow, print,
    property, range, repr, reversed, round, set, setattr, slice, sorted,
    staticmethod, str, sum, super, tuple, type, vars, zip,
    classmethod, dataclass, abstractmethod},
  morekeywords={[6]int, float, str, bool, list, dict, set, tuple, bytes,
    bytearray, complex, frozenset, memoryview, range, type, object,
    Exception, ValueError, TypeError, KeyError, IndexError, AttributeError,
    RuntimeError, StopIteration, GeneratorExit, OSError, IOError},
  morekeywords={[7]numpy, scipy, torch, tensorflow, pandas, matplotlib,
    os, sys, math, random, itertools, functools, collections, typing,
    pathlib, json, re, abc, dataclasses, enum},
  morekeywords={[8]property, staticmethod, classmethod, abstractmethod,
    dataclass, overload, override, deprecated},
  morecomment=[l]{\#},
  morecomment=[s]{"""}{"""},
  morecomment=[s]{'''}{'''},
}


% ═════════════════════════════════════════════════════════════════════════════
%  C
% ═════════════════════════════════════════════════════════════════════════════

\lstdefinestyle{catppuccin-c}{
  style=catppuccin-base,
  language=C,
  morekeywords={[1]auto, break, case, const, continue, default, do, else,
    enum, extern, for, goto, if, inline, register, restrict, return,
    sizeof, static, struct, switch, typedef, union, volatile, while,
    _Alignas, _Alignof, _Atomic, _Bool, _Complex, _Generic,
    _Imaginary, _Noreturn, _Static_assert, _Thread_local},
  morekeywords={[2]NULL, EOF, stdin, stdout, stderr, errno,
    EXIT_SUCCESS, EXIT_FAILURE, RAND_MAX, BUFSIZ},
  morekeywords={[3]true, false, INFINITY, NAN, HUGE_VAL,
    INT_MAX, INT_MIN, UINT_MAX, LONG_MAX, LONG_MIN,
    FLT_MAX, FLT_MIN, DBL_MAX, DBL_MIN, SIZE_MAX},
  morekeywords={[5]printf, fprintf, sprintf, snprintf, scanf, fscanf,
    malloc, calloc, realloc, free, memcpy, memmove, memset, memcmp,
    strlen, strcpy, strncpy, strcat, strcmp, strncmp, strstr, strchr,
    fopen, fclose, fread, fwrite, fgets, fputs, fseek, ftell, rewind,
    abs, atoi, atof, atol, strtol, strtod, qsort, bsearch,
    assert, exit, abort, atexit, system, getenv},
  morekeywords={[6]void, char, short, int, long, float, double, signed,
    unsigned, size_t, ptrdiff_t, intptr_t, uintptr_t, FILE,
    int8_t, int16_t, int32_t, int64_t,
    uint8_t, uint16_t, uint32_t, uint64_t,
    bool, complex, imaginary},
  morekeywords={[8]include, define, undef, ifdef, ifndef, endif,
    if, elif, else, pragma, error, warning, line},
  morecomment=[l]{//},
  morecomment=[s]{/*}{*/},
}


% ═════════════════════════════════════════════════════════════════════════════
%  C++
% ═════════════════════════════════════════════════════════════════════════════

\lstdefinestyle{catppuccin-cpp}{
  style=catppuccin-c,
  language=C++,
  morekeywords={[1]alignas, alignof, and, and_eq, bitand, bitor, catch,
    class, compl, concept, consteval, constexpr, constinit,
    co_await, co_return, co_yield, decltype, delete, dynamic_cast,
    explicit, export, final, friend, mutable, namespace, new, noexcept,
    not, not_eq, operator, or, or_eq, override, private, protected,
    public, reinterpret_cast, requires, static_assert, static_cast,
    template, this, throw, try, typeid, typename, using, virtual,
    xor, xor_eq},
  morekeywords={[2]nullptr, this, std, npos},
  morekeywords={[3]true, false, nullptr},
  morekeywords={[5]make_shared, make_unique, make_pair, make_tuple,
    move, forward, swap, begin, end, size, empty, push_back,
    emplace_back, insert, erase, find, count, sort, transform,
    accumulate, reduce, for_each, copy, fill, min, max,
    cout, cin, cerr, endl, getline},
  morekeywords={[6]auto, string, string_view, vector, array, deque,
    list, forward_list, set, multiset, map, multimap,
    unordered_set, unordered_map, stack, queue, priority_queue,
    pair, tuple, optional, variant, any, span,
    shared_ptr, unique_ptr, weak_ptr, function,
    thread, mutex, atomic, future, promise,
    istream, ostream, fstream, stringstream, iostream,
    iterator, const_iterator, reverse_iterator},
  morekeywords={[7]std, chrono, filesystem, ranges, views,
    this_thread, regex, literals},
}


% ═════════════════════════════════════════════════════════════════════════════
%  RUST
% ═════════════════════════════════════════════════════════════════════════════

\lstdefinelanguage{rust}{
  sensitive=true,
  morecomment=[l]{//},
  morecomment=[s]{/*}{*/},
  morestring=[b]",
  morestring=[b]',
  alsoletter={!},
}

\lstdefinestyle{catppuccin-rust}{
  style=catppuccin-base,
  language=rust,
  morekeywords={[1]as, async, await, break, const, continue, crate, dyn,
    else, enum, extern, fn, for, if, impl, in, let, loop, match, mod,
    move, mut, pub, ref, return, self, Self, static, struct, super,
    trait, type, unsafe, use, where, while, yield,
    abstract, become, box, do, final, macro, override, priv,
    try, typeof, unsized, virtual},
  morekeywords={[2]true, false, Some, None, Ok, Err, self, Self},
  morekeywords={[3]INFINITY, NEG_INFINITY, NAN, MAX, MIN,
    EPSILON, PI, E, TAU},
  morekeywords={[5]println!, print!, eprintln!, eprint!, format!,
    vec!, assert!, assert_eq!, assert_ne!, debug_assert!,
    panic!, todo!, unimplemented!, unreachable!,
    unwrap, expect, map, and_then, or_else, filter,
    collect, iter, into_iter, enumerate, zip, fold, reduce,
    push, pop, len, is_empty, contains, insert, remove, get,
    clone, to_string, to_owned, as_ref, as_mut, into, from},
  morekeywords={[6]bool, char, f32, f64, i8, i16, i32, i64, i128, isize,
    u8, u16, u32, u64, u128, usize, str, String,
    Vec, HashMap, HashSet, BTreeMap, BTreeSet, LinkedList, VecDeque,
    Box, Rc, Arc, Cell, RefCell, Mutex, RwLock,
    Option, Result, Iterator, IntoIterator, Display, Debug,
    Clone, Copy, Send, Sync, Sized, Drop, Fn, FnMut, FnOnce,
    From, Into, TryFrom, TryInto, AsRef, AsMut,
    Default, PartialEq, Eq, PartialOrd, Ord, Hash,
    Cow, Pin, Future, Stream, Error},
  morekeywords={[7]std, core, alloc, collections, fmt, io, fs, net,
    path, env, thread, sync, cell, mem, ptr, ops, iter,
    serde, tokio, rayon, anyhow, thiserror, clap, log, tracing},
  morekeywords={[8]derive, cfg, test, allow, warn, deny, forbid,
    inline, must_use, deprecated, doc, feature,
    macro_rules!, macro_export, proc_macro, proc_macro_derive},
}


% ═════════════════════════════════════════════════════════════════════════════
%  GO
% ═════════════════════════════════════════════════════════════════════════════

\lstdefinelanguage{go}{
  sensitive=true,
  morecomment=[l]{//},
  morecomment=[s]{/*}{*/},
  morestring=[b]",
  morestring=[b]',
  morestring=[b]`,
}

\lstdefinestyle{catppuccin-go}{
  style=catppuccin-base,
  language=go,
  morekeywords={[1]break, case, chan, const, continue, default, defer,
    else, fallthrough, for, func, go, goto, if, import, interface,
    map, package, range, return, select, struct, switch, type, var},
  morekeywords={[2]true, false, nil, iota},
  morekeywords={[3]MaxFloat32, MaxFloat64, MaxInt, MinInt,
    MaxInt8, MaxInt16, MaxInt32, MaxInt64,
    MaxUint, MaxUint8, MaxUint16, MaxUint32, MaxUint64},
  morekeywords={[5]append, cap, clear, close, complex, copy, delete,
    imag, len, make, max, min, new, panic, print, println,
    real, recover},
  morekeywords={[6]bool, byte, complex64, complex128, error, float32,
    float64, int, int8, int16, int32, int64, rune, string,
    uint, uint8, uint16, uint32, uint64, uintptr, any, comparable},
  morekeywords={[7]fmt, io, os, net, http, json, strings, strconv,
    math, sync, context, time, errors, log, testing, reflect,
    sort, bytes, bufio, regexp, encoding, crypto, filepath},
}


% ═════════════════════════════════════════════════════════════════════════════
%  JULIA
% ═════════════════════════════════════════════════════════════════════════════

\lstdefinelanguage{julia}{
  sensitive=true,
  morecomment=[l]{\#},
  morecomment=[s]{\#=}{=\#},
  morestring=[b]",
  morestring=[b]',
  alsoletter={!?},
}

\lstdefinestyle{catppuccin-julia}{
  style=catppuccin-base,
  language=julia,
  morekeywords={[1]abstract, baremodule, begin, break, catch, const,
    continue, do, else, elseif, end, export, finally, for, function,
    global, if, import, let, local, macro, module, mutable, new,
    primitive, quote, return, struct, try, type, using, while, where},
  morekeywords={[2]true, false, nothing, missing, undef, Inf, NaN},
  morekeywords={[3]pi, π, ℯ, im, Inf, Inf16, Inf32, Inf64,
    NaN, NaN16, NaN32, NaN64, ARGS, PROGRAM_FILE, DEPOT_PATH},
  morekeywords={[5]println, print, show, display, repr, string, sizeof,
    length, size, ndims, eltype, typeof, isa, convert, promote,
    push!, pop!, append!, insert!, deleteat!, splice!, empty!,
    sort!, sort, reverse!, reverse, unique, filter, filter!,
    map, map!, reduce, foldl, foldr, foreach, zip, enumerate,
    collect, keys, values, pairs, haskey, get, get!, merge,
    parse, tryparse, read, readline, write, open, close,
    rand, randn, zeros, ones, fill, fill!, copy, deepcopy,
    error, throw, rethrow, @assert, @show, @info, @warn, @error},
  morekeywords={[6]Any, Union, Nothing, Missing, UndefInitializer,
    Number, Real, Integer, AbstractFloat, Rational, Complex,
    Bool, Int, Int8, Int16, Int32, Int64, Int128, BigInt,
    UInt, UInt8, UInt16, UInt32, UInt64, UInt128,
    Float16, Float32, Float64, BigFloat,
    Char, String, SubString, Symbol, Regex,
    Array, Vector, Matrix, BitArray, BitVector, BitMatrix,
    Tuple, NamedTuple, Pair, Dict, Set, IdDict,
    AbstractArray, AbstractVector, AbstractMatrix, AbstractDict,
    IO, IOBuffer, IOStream, Ref, Ptr, Type,
    Function, Method, Module, Expr, LineNumberNode,
    Task, Channel, Condition, ReentrantLock, Threads,
    Exception, ErrorException, ArgumentError, BoundsError,
    DomainError, InexactError, MethodError, OverflowError,
    TypeError, UndefVarError, DivideError, StackOverflowError},
  morekeywords={[7]Base, Core, Main, LinearAlgebra, SparseArrays,
    Statistics, Random, Distributed, SharedArrays, Threads,
    Pkg, Test, Logging, Printf, Dates, REPL, InteractiveUtils,
    Plots, Makie, DifferentialEquations, Flux, Zygote, Optim,
    DataFrames, CSV, JuMP, ForwardDiff, ReverseDiff},
  morekeywords={[8]@time, @elapsed, @allocated, @benchmark,
    @show, @assert, @info, @warn, @error, @debug,
    @inline, @noinline, @simd, @threads, @spawn, @async,
    @generated, @eval, @macroexpand, @doc, @test, @testset,
    @inbounds, @views, @boundscheck, @label, @goto,
    @enum, @kwdef, @NamedTuple},
}


% ═════════════════════════════════════════════════════════════════════════════
%  MATLAB / OCTAVE
% ═════════════════════════════════════════════════════════════════════════════

\lstdefinestyle{catppuccin-matlab}{
  style=catppuccin-base,
  language=Matlab,
  morekeywords={[1]break, case, catch, classdef, continue, else, elseif,
    end, enumeration, events, for, function, global, if, methods,
    otherwise, parfor, persistent, properties, return, spmd, switch,
    try, while, arguments},
  morekeywords={[2]true, false, inf, Inf, nan, NaN, eps, pi, i, j,
    end, nargin, nargout, varargin, varargout, ans},
  morekeywords={[3]eps, realmin, realmax, intmax, intmin, flintmax},
  morekeywords={[5]abs, acos, acosh, asin, asinh, atan, atan2, atanh,
    ceil, cos, cosh, exp, floor, log, log2, log10, mod, rem,
    round, sin, sinh, sqrt, tan, tanh, sign, power, fix,
    length, size, numel, ndims, reshape, repmat, squeeze, permute,
    cat, horzcat, vertcat, zeros, ones, eye, rand, randn, randi,
    linspace, logspace, meshgrid, colon,
    diag, trace, det, inv, eig, svd, qr, lu, chol, norm, rank,
    kron, cross, dot, null, orth, pinv, expm, logm, sqrtm,
    fft, ifft, fft2, ifft2, conv, conv2, deconv, filter,
    plot, subplot, figure, hold, xlabel, ylabel, title, legend,
    grid, axis, xlim, ylim, colorbar, imagesc, contour, surf, mesh,
    disp, fprintf, sprintf, input, error, warning, assert,
    save, load, fopen, fclose, fread, fwrite, fscanf, feof,
    sort, unique, find, min, max, sum, prod, cumsum, mean, std, var,
    diff, gradient, trapz, integral, ode45, ode23, fsolve, fzero,
    struct, cell, fieldnames, isfield, rmfield, cellfun, arrayfun,
    strcmp, strcat, strsplit, strtrim, num2str, str2num, str2double,
    class, isa, isnumeric, ischar, islogical, isempty, isequal,
    deal, nargchk, inputParser, addRequired, addOptional, parse},
  morekeywords={[6]double, single, int8, int16, int32, int64,
    uint8, uint16, uint32, uint64, logical, char, string,
    cell, struct, table, timetable, categorical, datetime, duration,
    function_handle, containers.Map, gpuArray, sparse, tall},
  morekeywords={[7]Simulink, Control, Signal, Image, Statistics,
    Optimization, Symbolic, Parallel, DeepLearning},
  morecomment=[l]{\%},
  morecomment=[s]{\%\{}{\%\}},
}


% ═════════════════════════════════════════════════════════════════════════════
%  FORTRAN (F90/F95/2003/2008/2018)
% ═════════════════════════════════════════════════════════════════════════════

\lstdefinestyle{catppuccin-fortran}{
  style=catppuccin-base,
  language=[08]Fortran,
  sensitive=false,
  morekeywords={[1]program, end, subroutine, function, module, submodule,
    use, only, implicit, none, contains, interface, abstract, procedure,
    call, return, stop, error, if, then, else, elseif, endif,
    do, enddo, while, cycle, exit, forall, where, elsewhere,
    select, case, default, associate, block, critical,
    go, to, goto, continue, allocate, deallocate, nullify,
    open, close, read, write, print, format, rewind, backspace,
    inquire, namelist, data, save, common, equivalence,
    import, class, type, extends, generic, deferred, non_overridable,
    pass, nopass, final, enum, enumerator, bind, value, volatile,
    asynchronous, contiguous, codimension, concurrent, sync,
    pure, elemental, impure, recursive, result},
  morekeywords={[2].true., .false., .and., .or., .not., .eqv., .neqv.,
    .eq., .ne., .lt., .le., .gt., .ge.},
  morekeywords={[3]huge, tiny, epsilon, radix, precision, range,
    maxexponent, minexponent, kind},
  morekeywords={[5]abs, acos, aimag, aint, asin, atan, atan2, cos, cosh,
    dble, exp, int, log, log10, max, min, mod, nint, real, sign,
    sin, sinh, sqrt, tan, tanh, char, cmplx, conjg,
    matmul, dot_product, transpose, reshape, shape, size, lbound,
    ubound, maxval, minval, sum, product, count, any, all, merge,
    pack, unpack, spread, maxloc, minloc, trim,
    adjustl, adjustr, len_trim, repeat, scan, verify,
    selected_int_kind, selected_real_kind, allocated, associated,
    present, null, move_alloc, cpu_time, system_clock,
    date_and_time, random_number, random_seed},
  morekeywords={[6]integer, real, double, precision, complex, character,
    logical, type, class, dimension, allocatable, pointer, target,
    intent, in, out, inout, optional, parameter, external, intrinsic},
  morekeywords={[7]iso_fortran_env, iso_c_binding, ieee_arithmetic,
    ieee_exceptions, ieee_features, omp_lib, mpi, mpi_f08},
  morekeywords={[8]!$omp, !$acc, !DEC$, !DIR$, !GCC$},
  morecomment=[l]{!},
}


% ═════════════════════════════════════════════════════════════════════════════
%  ASSEMBLER (x86-64, sintaxis Intel / NASM)
% ═════════════════════════════════════════════════════════════════════════════

\lstdefinelanguage{x86asm}{
  sensitive=false,
  morecomment=[l]{;},
  morestring=[b]",
  morestring=[b]',
}

\lstdefinestyle{catppuccin-asm}{
  style=catppuccin-base,
  language=x86asm,
  morekeywords={[1]mov, movzx, movsx, movsxd, lea, push, pop, pushf, popf,
    add, sub, mul, imul, div, idiv, inc, dec, neg, adc, sbb,
    and, or, xor, not, shl, shr, sar, sal, rol, ror, rcl, rcr,
    test, cmp, bt, bts, btr, btc, bsf, bsr,
    jmp, je, jne, jz, jnz, jg, jge, jl, jle, ja, jae, jb, jbe,
    jc, jnc, jo, jno, js, jns, jp, jnp, jecxz, jrcxz,
    call, ret, enter, leave, int, iret, syscall, sysret,
    loop, loope, loopne, rep, repe, repne,
    nop, hlt, wait, lock, xchg, cmpxchg, xadd,
    cmove, cmovne, cmovg, cmovge, cmovl, cmovle,
    cmova, cmovae, cmovb, cmovbe,
    sete, setne, setg, setge, setl, setle, seta, setae, setb, setbe,
    cbw, cwde, cdqe, cwd, cdq, cqo,
    movs, stos, lods, cmps, scas,
    in, out, cli, sti, clc, stc, cmc, cld, std},
  morekeywords={[2]movss, movsd, movaps, movups, movapd, movupd, movdqa,
    addss, addsd, addps, addpd, subss, subsd, subps, subpd,
    mulss, mulsd, mulps, mulpd, divss, divsd, divps, divpd,
    sqrtss, sqrtsd, sqrtps, sqrtpd, maxss, maxsd, minss, minsd,
    cmpss, cmpsd, comiss, comisd, ucomiss, ucomisd,
    cvtsi2ss, cvtsi2sd, cvtss2si, cvtsd2si, cvtss2sd, cvtsd2ss,
    shufps, shufpd, unpckhps, unpcklps, blendps, blendpd,
    vaddps, vaddpd, vmulps, vmulpd, vsubps, vsubpd, vdivps, vdivpd,
    vfmadd132ps, vfmadd213ps, vfmadd231ps},
  morekeywords={[5]section, segment, global, extern, default, align,
    db, dw, dd, dq, dt, do, dy, dz,
    resb, resw, resd, resq, rest, reso, resy, resz,
    equ, times, incbin, org, bits, use16, use32, use64,
    .text, .data, .bss, .rodata, .section},
  morekeywords={[6]rax, rbx, rcx, rdx, rsi, rdi, rbp, rsp,
    r8, r9, r10, r11, r12, r13, r14, r15,
    eax, ebx, ecx, edx, esi, edi, ebp, esp,
    ax, bx, cx, dx, si, di, bp, sp,
    al, ah, bl, bh, cl, ch, dl, dh,
    sil, dil, bpl, spl, r8b, r9b, r10b, r11b, r12b, r13b, r14b, r15b,
    r8w, r9w, r10w, r11w, r12w, r13w, r14w, r15w,
    r8d, r9d, r10d, r11d, r12d, r13d, r14d, r15d,
    cs, ds, es, fs, gs, ss, rip, eip, ip, rflags, eflags, flags,
    xmm0, xmm1, xmm2, xmm3, xmm4, xmm5, xmm6, xmm7,
    xmm8, xmm9, xmm10, xmm11, xmm12, xmm13, xmm14, xmm15,
    ymm0, ymm1, ymm2, ymm3, ymm4, ymm5, ymm6, ymm7,
    ymm8, ymm9, ymm10, ymm11, ymm12, ymm13, ymm14, ymm15,
    zmm0, zmm1, zmm2, zmm3, zmm4, zmm5, zmm6, zmm7,
    cr0, cr2, cr3, cr4, dr0, dr1, dr2, dr3, dr6, dr7},
  morekeywords={[8]macro, endmacro, %macro, %endmacro, %define, %undef,
    %ifdef, %ifndef, %endif, %elif, %else, %if, %assign, %rep, %endrep,
    %include, %push, %pop, %repl, %error, %warning, struc, endstruc,
    BITS, SECTION, GLOBAL, EXTERN},
}


% ═════════════════════════════════════════════════════════════════════════════
%  MATHEMATICA / WOLFRAM LANGUAGE
% ═════════════════════════════════════════════════════════════════════════════

\lstdefinelanguage{mathematica}{
  sensitive=true,
  morecomment=[s]{(*}{*)},
  morestring=[b]",
  alsoletter={$},
}

\lstdefinestyle{catppuccin-mathematica}{
  style=catppuccin-base,
  language=mathematica,
  morekeywords={[1]Module, Block, With, Do, For, While, If, Which, Switch,
    Return, Break, Continue, Goto, Label, Throw, Catch, Check,
    CompoundExpression, Set, SetDelayed, Rule, RuleDelayed,
    Function, Slot, SlotSequence, Apply, Map, MapThread, Scan,
    Table, Array, Range, Nest, NestList, Fold, FoldList, FixedPoint,
    Hold, HoldForm, HoldAll, HoldFirst, HoldRest, Evaluate, ReleaseHold,
    Condition, PatternTest, Alternatives, Repeated, Except, Blank,
    BlankSequence, BlankNullSequence, Optional, Default, Pattern,
    Begin, End, BeginPackage, EndPackage, Needs, Get},
  morekeywords={[2]True, False, Null, None, All, Full, Automatic,
    Infinity, DirectedInfinity, Indeterminate, Undefined,
    $Failed, $Aborted, $MachineEpsilon, $MachinePrecision,
    $Context, $ContextPath, $Path, $Version},
  morekeywords={[3]Pi, E, I, Degree, EulerGamma, GoldenRatio,
    Catalan, Glaisher, Khinchin, MachinePrecision, MachineEpsilon},
  morekeywords={[5]Print, Echo, Assert, Abort, Quit, Exit,
    Plus, Times, Power, Sqrt, Exp, Log, Sin, Cos, Tan,
    ArcSin, ArcCos, ArcTan, Sinh, Cosh, Tanh,
    Abs, Sign, Re, Im, Conjugate, Arg, Mod, Quotient,
    Floor, Ceiling, Round, IntegerPart, FractionalPart, Max, Min,
    D, Dt, Integrate, NIntegrate, Sum, Product, Limit,
    Series, Normal, Coefficient, CoefficientList,
    Solve, NSolve, DSolve, NDSolve, FindRoot, Minimize, Maximize,
    NMinimize, NMaximize, LinearSolve, Eigenvalues, Eigenvectors,
    Simplify, FullSimplify, Expand, Factor, Apart, Together, Cancel,
    List, Part, First, Last, Rest, Most, Take, Drop, Select, Cases,
    Position, Length, Dimensions, Flatten, Partition, Join, Union,
    Intersection, Complement, Sort, Reverse, RotateLeft, RotateRight,
    Transpose, Inverse, Det, Tr, Dot, Cross, Outer, Inner,
    MatrixForm, TableForm, TeXForm, InputForm, FullForm,
    Plot, Plot3D, ListPlot, ContourPlot, DensityPlot, ParametricPlot,
    Show, Graphics, Style, RGBColor, Hue, Opacity, Thickness,
    Import, Export, ReadList, WriteString, OpenRead, OpenWrite, Close,
    StringJoin, StringReplace, StringSplit, StringQ, NumberQ, IntegerQ,
    Head, MatchQ, FreeQ, MemberQ, DeleteCases, Replace, ReplaceAll},
  morekeywords={[6]Integer, Real, Rational, Complex, String, Symbol,
    List, Association, SparseArray, NumericArray,
    Graph, Region, MeshRegion, BooleanRegion,
    Integers, Reals, Rationals, Complexes, Booleans, Primes,
    Algebraics, PositiveReals, NonNegativeReals, PositiveIntegers},
  morekeywords={[7]System, Global, Internal, Developer, Experimental,
    GeneralUtilities, Notation, CCompilerDriver, NeuralNetworks,
    MachineLearning, ResourceFunction},
}


% ═════════════════════════════════════════════════════════════════════════════
% ALIASES DE CONVENIENCIA
% ═════════════════════════════════════════════════════════════════════════════
\lstdefinestyle{python}{style=catppuccin-python}
\lstdefinestyle{c}{style=catppuccin-c}
\lstdefinestyle{cpp}{style=catppuccin-cpp}
\lstdefinestyle{rust}{style=catppuccin-rust}
\lstdefinestyle{go}{style=catppuccin-go}
\lstdefinestyle{julia}{style=catppuccin-julia}
\lstdefinestyle{matlab}{style=catppuccin-matlab}
\lstdefinestyle{fortran}{style=catppuccin-fortran}
\lstdefinestyle{asm}{style=catppuccin-asm}
\lstdefinestyle{mathematica}{style=catppuccin-mathematica}


% ═════════════════════════════════════════════════════════════════════════════
%  LATEX / TEX
% ═════════════════════════════════════════════════════════════════════════════
% NOTA CRÍTICA: ningún keyword puede llevar backslash dentro de morekeywords
% ni de texcs. Los comandos \par, \vspace, \noindent, etc. son primitivas
% activas que LaTeX ejecuta al parsear el argumento — causa errores de
% compilación ("Paragraph ended before...", "Illegal unit of measure", etc.)
%
% texcs acepta nombres SIN backslash: listings añade el \ al colorear.
% morekeywords[2..6] también usan nombres sin backslash.
%
% Mapa semántico:
%   texcs   Estructura      begin end chapter section label cite input ...
%   [2]     Definición      newcommand renewcommand DeclareMathOperator ...
%   [3]     Formato         textbf textit emph centering medskip ...
%   [4]     Matemáticas     frac sqrt sum int mathbb hat alpha ...
%   [5]     Paquetes        amsmath tikz tcolorbox biblatex ...
%   [6]     Especiales      LaTeX TeX today toprule captionof ...

\lstdefinestyle{catppuccin-latex}{
  style        = catppuccin-base,
  language     = [LaTeX]TeX,
  % ── Estructura (texcs: sin backslash, listings añade \) ───────────────────
  texcs        = {
    begin, end,
    chapter, section, subsection, subsubsection, paragraph, subparagraph,
    label, ref, cref, Cref, crefrange,
    cite, citet, citep, footcite, nocite,
    input, include, includeonly, usepackage,
    documentclass, maketitle, makeatletter, makeatother,
    tableofcontents, listoffigures, listoftables,
    appendix, printbibliography, addbibresource, bibliography,
    part, frontmatter, mainmatter, backmatter,
    addcontentsline, phantomsection,
  },
  % ── Definición de comandos (sin backslash) ─────────────────────────────────
  morekeywords = {[2]
    newcommand, renewcommand, providecommand,
    newenvironment, renewenvironment,
    newcounter, setcounter, addtocounter, refstepcounter,
    DeclareMathOperator,
    DeclareRobustCommand,
    NewDocumentCommand, NewDocumentEnvironment,
    colorlet, definecolor,
    tcbset, tcolorboxenvironment,
    pgfplotsset, pgfplotscreateplotcyclelist,
    usetikzlibrary, tikzset,
    lstdefinestyle, lstdefinelanguage, lstset,
    SetKwSty, SetArgSty, SetCommentSty, SetNlSty,
    geometry, hypersetup, sisetup,
    setlength, addtolength,
    pagestyle, fancyhf, fancyhead, fancyfoot, fancypagestyle,
    crearEntorno, createback, addplot,
  },
  % ── Formato de texto (sin backslash) ──────────────────────────────────────
  morekeywords = {[3]
    textbf, textit, texttt, textrm, textsf,
    emph, underline, textcolor, color,
    textsc, textsuperscript, textsubscript,
    small, footnotesize, scriptsize, tiny,
    large, Large, LARGE, huge, Huge, normalsize,
    bfseries, itshape, ttfamily, rmfamily, sffamily,
    centering, raggedright, raggedleft,
    medskip, bigskip, smallskip, vfill, hfill,
  },
  % ── Matemáticas (sin backslash) ────────────────────────────────────────────
  morekeywords = {[4]
    frac, dfrac, tfrac, cfrac,
    sqrt, sum, prod, int, oint, iint, iiint,
    lim, sup, inf, max, min,
    left, right, mleft, mright,
    mathbb, mathcal, mathscr, mathbf, mathit,
    mathrm, mathfrak, mathsf,
    hat, bar, tilde, vec, dot, ddot, widehat,
    partial, nabla, forall, exists,
    leq, geq, neq, approx, equiv, sim,
    cdot, times, otimes, oplus,
    ldots, cdots, vdots, ddots,
    binom, text, operatorname,
    alpha, beta, gamma, delta, epsilon, varepsilon,
    theta, lambda, mu, pi, sigma, tau, phi, omega,
    Gamma, Delta, Theta, Lambda, Pi, Sigma, Omega,
  },
  % ── Paquetes (sin backslash) ───────────────────────────────────────────────
  morekeywords = {[5]
    amsmath, amssymb, amsthm, mathtools,
    physics, bm, cancel, nicematrix, stmaryrd,
    tikz, pgfplots, tcolorbox, listings,
    hyperref, cleveref, biblatex,
    booktabs, tabularx, array, colortbl,
    graphicx, caption, subcaption, float,
    enumitem, xcolor, microtype,
    algorithm2e, many, babel,
    siunitx, pgfplotstable, forest, qtree,
    titlesec, fancyhdr, setspace, parskip,
    appendix, footmisc, etoolbox, textpos,
    catppuccinpalette, mathrsfs, mleftright,
  },
  % ── Especiales (sin backslash) ─────────────────────────────────────────────
  morekeywords = {[6]
    LaTeX, LaTeXe, TeX, BibTeX, BibLaTeX,
    today, thepage, thechapter, thesection,
    hline, toprule, midrule, bottomrule, cmidrule,
    rowcolor, cellcolor, arrayrulecolor,
    caption, captionof, includegraphics,
    scalebox, resizebox, textbackslash,
    texorpdfstring,
    dificultad, headerrow, thead,
    eqop, eqfn, eqdm, eqdom, equ, eqk, eqcon, eqvec, eqten,
  },
  % ── Colores por clase ──────────────────────────────────────────────────────
  texcsstyle   = {\color{code-keyword}\bfseries},  % estructura
  keywordstyle = {[2]\color{code-fn}},              % definición
  keywordstyle = {[3]\color{code-operator}},        % formato
  keywordstyle = {[4]\color{code-macro}},           % matemáticas
  keywordstyle = {[5]\color{code-mod}},             % paquetes
  keywordstyle = {[6]\color{code-builtin}},         % especiales
  morecomment  = [l]{\%},
}

\lstdefinestyle{latex}{style=catppuccin-latex}
\lstdefinestyle{tex}{style=catppuccin-latex}

%───────────────────────────────────────────────────────────────────────────────
% ALGORITHMS.TEX — Estilo de algoritmos con colores semánticos Catppuccin
%───────────────────────────────────────────────────────────────────────────────
% Consume aliases del dominio alg-* definidos en catppuccin-palette.tex.
% No usa Ctp* directamente.
%
% DECISIÓN DE DISEÑO — sin wrapper tcolorbox:
%   algorithm2e con [ruled, vlined, linesnumbered] es un contenedor visual
%   completo (líneas de delimitación, caption, numeración, líneas de bloque).
%   Envolverlo en tcolorbox superpondría dos sistemas visuales sin ganancia.
%   El theming completo se logra únicamente vía los comandos \Set*Sty y los
%   parches de color de líneas — sin restricciones de breakable ni anidamiento.
%───────────────────────────────────────────────────────────────────────────────

% ─── Estilos de texto ────────────────────────────────────────────────────────
\newcommand{\algkw}[1]{\textcolor{alg-keyword}{#1}}     % palabras clave
\newcommand{\algarg}[1]{\textcolor{alg-arg}{#1}}         % argumentos
\newcommand{\algcomment}[1]{\textcolor{alg-comment}{#1}} % comentarios
\newcommand{\algnum}[1]{\textcolor{alg-linenum}{#1}}     % números de línea
\newcommand{\algcap}[1]{\textcolor{alg-caption}{#1}}     % caption

% ─── Aplicar estilos a algorithm2e ───────────────────────────────────────────
\SetKwSty{algkw}
\SetArgSty{algarg}
\SetCommentSty{algcomment}
\SetNlSty{algnum}{}{}
\SetAlCapSty{algcap}

% ─── Grosor de líneas ────────────────────────────────────────────────────────
\setlength{\algoheightrule}{0.6pt}       % líneas superior e inferior
\setlength{\algotitleheightrule}{0.5pt}  % línea bajo el caption

% ─── Color de líneas estructurales ───────────────────────────────────────────
% Parches internos de algorithm2e para colorear las líneas ruled y vlined
% con alg-line. Separado de alg-linenum (números) para control independiente.
\makeatletter
\newcommand{\algolinecolor}{alg-line}

\def\@algocf@pre@ruled{%
  {\color{\algolinecolor}\hrule width\algocf@ruledwidth
   height\algoheightrule depth0pt}%
  \kern\interspacetitleruled}

\def\@algocf@post@ruled{%
  \kern\interspacealgoruled
  {\color{\algolinecolor}\hrule width\algocf@ruledwidth
   height\algoheightrule}}

\renewcommand{\algocf@caption@ruled}{%
  \box\algocf@capbox
  \kern\interspacetitleruled
  {\color{\algolinecolor}\hrule width\algocf@ruledwidth
   height\algotitleheightrule depth0pt}%
  \kern\interspacealgoruled}
\makeatother

% ─── Línea vertical izquierda (vlined) ───────────────────────────────────────
\makeatletter
\renewcommand{\algocf@Vline}[1]{%
  \strut\par\nointerlineskip%
  \algocf@push{\skiprule}%
  \hbox{{\color{\algolinecolor}\vrule}%
    \vtop{\algocf@push{\skiptext}%
      \vtop{\algocf@addskiptotal #1}\Hlne}}\vskip\skiphlne%
  \algocf@pop{\skiprule}%
  \nointerlineskip}

\renewcommand{\algocf@Vsline}[1]{%
  \strut\par\nointerlineskip%
  \algocf@bblockcode%
  \algocf@push{\skiprule}%
  \hbox{{\color{\algolinecolor}\vrule}%
    \vtop{\algocf@push{\skiptext}%
      \vtop{\algocf@addskiptotal #1}}}%
  \algocf@pop{\skiprule}%
  \algocf@eblockcode%
}
\makeatother

%───────────────────────────────────────────────────────────────────────────────
% EQUATION-STYLES.TEX — Coloreado semántico de ecuaciones (taxonomía híbrida)
%───────────────────────────────────────────────────────────────────────────────
%
% Sistema de "syntax highlighting" para matemáticas en tres ejes ortogonales:
%
%   1. ROL SINTÁCTICO — estable entre contextos
%      Operador, función, diferencial, dominio: siempre el mismo color.
%
%   2. ESTATUS EPISTÉMICO — la dimensión que la tipografía NO codifica
%      Incógnita vs dato/parámetro vs constante universal.
%      En f(x;θ): en física x es incógnita y θ dato; en ML es al revés.
%      El color hace explícita esta inversión sin texto adicional.
%
%   3. TIPO MATEMÁTICO — refuerzo sutil de la tipografía
%      Vector, tensor/matriz: la tipografía (bold, doble barra) ya lo indica;
%      el color refuerza cuando distinguir rango es semánticamente importante.
%
% Los colores se definen en catppuccin-palette.tex bajo el prefijo eq-*.
% Este archivo define únicamente el toggle y las macros.
%
% Toggle global: \eqcolorstrue / \eqcolorsfalse
% Uso en grupo local para versión B/N: {\eqcolorsfalse \begin{equation}...}
%───────────────────────────────────────────────────────────────────────────────

% ─── Toggle global ───────────────────────────────────────────────────────────
\newif\ifeqcolors
\eqcolorstrue     % \eqcolorsfalse para versión impresión B/N

% ═════════════════════════════════════════════════════════════════════════════
% MACROS
% ═════════════════════════════════════════════════════════════════════════════
\makeatletter

% ─── Comando interno ─────────────────────────────────────────────────────────
% Aplica color solo si el toggle está activo.
\newcommand{\@eqc}[2]{%
  \ifeqcolors\textcolor{#1}{#2}\else#2\fi
}

% ── Eje 1: Rol sintáctico ────────────────────────────────────────────────────
%   Macro      Alias de color    Cubre
%   ─────      ─────────────     ─────
%   \eqop      eq-op             Operadores: ∫, ∑, ∏, ∂, ∇, ∇², ∇×, lím, Ĥ
%   \eqfn      eq-fn             Funciones/mapas: f(·), L(θ), σ(·), V(r), ReLU
%   \eqdm      eq-dm             Diferenciales/medidas: dx, dt, dW, dS, dΩ
%   \eqdom     eq-dom            Dominios/conjuntos: Ω, ∂Ω, ℝⁿ, H, S²
\newcommand{\eqop}[1]{\@eqc{eq-op}{#1}}
\newcommand{\eqfn}[1]{\@eqc{eq-fn}{#1}}
\newcommand{\eqdm}[1]{\@eqc{eq-dm}{#1}}
\newcommand{\eqdom}[1]{\@eqc{eq-dom}{#1}}

% ── Eje 2: Estatus epistémico ────────────────────────────────────────────────
%   Macro      Alias de color    Cubre
%   ─────      ─────────────     ─────
%   \equ       eq-unknown        Incógnita / buscado / aprendido
%   \eqk       eq-known          Conocido / dato / fijado
%   \eqcon     eq-const          Constante universal / estructural: ℏ, kB, π
\newcommand{\equ}[1]{\@eqc{eq-unknown}{#1}}
\newcommand{\eqk}[1]{\@eqc{eq-known}{#1}}
\newcommand{\eqcon}[1]{\@eqc{eq-const}{#1}}

% ── Eje 3: Tipo matemático ───────────────────────────────────────────────────
%   Macro      Alias de color    Cubre
%   ─────      ─────────────     ─────
%   \eqvec     eq-vec            Campos vectoriales: u, E, B, J, f
%   \eqten     eq-ten            Tensores/matrices: σᵢⱼ, gμν, W, Σ
\newcommand{\eqvec}[1]{\@eqc{eq-vec}{#1}}
\newcommand{\eqten}[1]{\@eqc{eq-ten}{#1}}

% ═════════════════════════════════════════════════════════════════════════════
% MACROS COMPUESTOS
% ═════════════════════════════════════════════════════════════════════════════

% Derivada parcial coloreada
% Uso: \eqpdv{ρ}{t}  →  ∂ρ/∂t
% Numerador: incógnita. Denominador: variable independiente (dato del problema).
\newcommand{\eqpdv}[2]{%
  \eqop{\frac{\partial \equ{#1}}{\partial \eqk{#2}}}%
}

% Gradiente coloreado
% Uso: \eqgrad{u}  →  ∇u
\newcommand{\eqgrad}[1]{%
  \eqop{\nabla}\equ{#1}%
}

% Divergencia coloreada
% Uso: \eqdvg{\mathbf{u}}  →  ∇·u
% Nombre \eqdvg para no colisionar con \div del paquete physics.
\newcommand{\eqdvg}[1]{%
  \eqop{\nabla \cdot} \equ{#1}%
}

% SDE coloreada
% Uso: \eqSDE{X}{-\theta X}{\sigma}
% Estado = incógnita, drift = función, coef. difusión = conocido.
\newcommand{\eqSDE}[3]{%
  \eqdm{d}\equ{#1} = \eqfn{#2}\,\eqdm{dt} + \eqk{#3}\,\eqdm{dW}%
}

\makeatother

%───────────────────────────────────────────────────────────────────────────────
% PGFPLOTS.TEX — Estilos de gráficas con colores semánticos Catppuccin
%───────────────────────────────────────────────────────────────────────────────
% Consume aliases del dominio plot-* definidos en catppuccin-palette.tex.
% No usa Ctp* directamente.
% NOTA: \pgfplotsset{compat=1.18} se declara en packages.tex inmediatamente
%       después de cargar pgfplots — es inicialización de versión, no config.
%───────────────────────────────────────────────────────────────────────────────

% ─── Cycle list Catppuccin ───────────────────────────────────────────────────
% 8 series, ordenadas por contraste visual decreciente.
% Corresponde a plot-series-1 ... plot-series-8 del palette.
\pgfplotscreateplotcyclelist{catppuccin}{
  {plot-series-1, mark=*,         thick},
  {plot-series-2, mark=square*,   thick},
  {plot-series-3, mark=triangle*, thick},
  {plot-series-4, mark=diamond*,  thick},
  {plot-series-5, mark=pentagon*, thick},
  {plot-series-6, mark=o,         thick},
  {plot-series-7, mark=square,    thick},
  {plot-series-8, mark=triangle,  thick},
}

% ─── Estilo global de ejes ───────────────────────────────────────────────────
\pgfplotsset{
  cycle list name=catppuccin,
  every axis/.append style={
    axis line style  = {plot-axis},
    tick style       = {plot-axis},
    label style      = {color=plot-label},
    tick label style = {color=plot-tick-label},
    title style      = {color=plot-label, font=\bfseries},
    legend style     = {
      draw=plot-legend-border,
      fill=plot-legend-bg,
      text=maintext,
      font=\small,
      rounded corners=1pt,
    },
    grid style = {
      plot-grid,
      line width=0.3pt,
      dashed,
    },
  },
  every error bar/.append style={
    thin,
    plot-errorbar,
  },
}

% ─── Estilos con nombre ───────────────────────────────────────────────────────

% Gráfica limpia para publicación
\pgfplotsset{
  catppuccin-clean/.style={
    grid=major,
    xlabel style={font=\small},
    ylabel style={font=\small},
    tick label style={font=\footnotesize},
    legend pos=north east,
    legend cell align=left,
  },
}

% Gráfica con fondo coloreado (para presentaciones o notas)
\pgfplotsset{
  catppuccin-filled/.style={
    catppuccin-clean,
    axis background/.style={fill=secondarybg},
  },
}

% ─── Colormaps ────────────────────────────────────────────────────────────────
% Secuencial: magnitudes positivas (Base → series frías → series cálidas)
\pgfplotsset{
  colormap={catppuccin-heat}{
    color=(background)
    color=(plot-series-1)
    color=(plot-series-4)
    color=(plot-series-5)
    color=(plot-series-7)
  },
  % Divergente: residuos con signo (series-1 ←→ maintext ←→ series-5)
  colormap={catppuccin-diverge}{
    color=(plot-series-1)
    color=(maintext)
    color=(plot-series-5)
  },
}


% ─── 9. Notación matemática ──────────────────────────────────────────────────
% Al final de la preámbula: puede usar cualquier paquete ya cargado.
%───────────────────────────────────────────────────────────────────────────────
% NOTATION.TEX — Vocabulario matemático del documento
%───────────────────────────────────────────────────────────────────────────────
%
% Este archivo define los objetos matemáticos propios del documento.
% NO configura paquetes — eso vive en config/equations.tex y demás.
%
% RELACIÓN CON OTROS ARCHIVOS:
%   physics       →  \abs, \norm, \tr, \grad, \div, \curl, \laplacian,
%                    \dv, \pdv, \bra, \ket, \braket, \comm, \acomm, etc.
%   equations.tex →  \eqop, \equ, \eqk, \eqfn, \eqdm, \eqdom (coloreado semántico)
%                    \eqpdv, \eqgrad, \eqdvg, \eqSDE  (compuestos coloreados)
%   notation.tex  →  todo lo demás: objetos del dominio, operadores propios,
%                    distribuciones, shortcuts de uso frecuente.
%
% REGLA DE EXTENSIÓN:
%   Añadir secciones por dominio. Nunca mezclar dominios dentro de una sección.
%   Si una sección supera ~30 comandos, considerar un archivo propio.
%
% CONFLICTOS CONOCIDOS CON physics:
%   \div   →  physics lo redefine como divergencia. Usar \eqdvg (coloreado)
%             o \divergence (physics). Para división: / o \frac{}{}.
%   \rank  →  physics define \rank. Para sobrescribir: ver nota al final.
%───────────────────────────────────────────────────────────────────────────────


% =============================================================================
% ÁLGEBRA Y ANÁLISIS
% =============================================================================

% ─── Conjuntos numéricos ─────────────────────────────────────────────────────
\newcommand{\R}{\mathbb{R}}
\newcommand{\N}{\mathbb{N}}
\newcommand{\Z}{\mathbb{Z}}
\newcommand{\C}{\mathbb{C}}
\newcommand{\Q}{\mathbb{Q}}
\newcommand{\F}{\mathbb{F}}          % cuerpo genérico (álgebra abstracta)

% ─── Operadores matriciales y lineales ───────────────────────────────────────
% \tr, \rank ya existen en physics. Ver nota al final para sobrescribir \rank.
\newcommand{\T}{^{\top}}             % transpuesta: A\T
\newcommand{\HH}{^{\dagger}}         % adjunto hermítico (†): A\HH
\newcommand{\Id}{\mathbf{I}}         % matriz identidad
\DeclareMathOperator{\diag}{diag}    % \diag(a,b,c) o \diag(A)
\DeclareMathOperator{\spec}{spec}    % espectro de un operador

% ─── Optimización ────────────────────────────────────────────────────────────
\DeclareMathOperator*{\argmin}{arg\,min}
\DeclareMathOperator*{\argmax}{arg\,max}
\DeclareMathOperator*{\minimize}{minimize}
\DeclareMathOperator*{\maximize}{maximize}

% ─── Aliases de physics (nombres cortos para uso frecuente) ──────────────────
% Mantener solo los que se usan en el documento con alta frecuencia.
% Para el resto: usar directamente \laplacian, \divergence, etc.
\newcommand{\Lap}{\laplacian}        % alias → \laplacian (∇²)
\newcommand{\Div}{\divergence}       % alias → \divergence (∇·)


% =============================================================================
% PROBABILIDAD Y ESTADÍSTICA
% =============================================================================

% ─── Operadores de esperanza y probabilidad ───────────────────────────────────
% \mleft[\mright] corrige el espaciado que \left[\right] introduce en
% ciertos contextos (ver paquete mleftright).
\newcommand{\E}[1]{\mathbb{E}\!\mleft[#1\mright]}
\newcommand{\Prob}[1]{\mathbb{P}\!\mleft(#1\mright)}
\newcommand{\Ep}[1]{\langle #1 \rangle}   % valor esperado estilo física ⟨·⟩

% ─── Operadores estadísticos ─────────────────────────────────────────────────
\DeclareMathOperator{\KL}{KL}
\DeclareMathOperator{\Var}{Var}
\DeclareMathOperator{\Cov}{Cov}
\DeclareMathOperator{\Ent}{H}        % entropía: H(X)
\DeclareMathOperator{\MI}{I}         % información mutua: I(X;Y)

% ─── Distribuciones ──────────────────────────────────────────────────────────
\newcommand{\Normal}{\mathcal{N}}    % \Normal(\mu, \sigma^2)
\newcommand{\Uniform}{\mathcal{U}}   % \Uniform(a, b)
\newcommand{\Bernoulli}{\mathrm{Ber}}
\newcommand{\Poisson}{\mathrm{Poi}}
\newcommand{\Categorical}{\mathrm{Cat}}


% =============================================================================
% MACHINE LEARNING
% =============================================================================

% ─── Objetos del modelo ──────────────────────────────────────────────────────
\newcommand{\Loss}{\mathcal{L}}      % función de pérdida ℒ
\newcommand{\Lscr}{\mathscr{L}}      % pérdida en notación script (alternativa)
\newcommand{\D}{\mathcal{D}}         % dataset 𝒟
\newcommand{\Params}{\Theta}         % espacio de parámetros Θ
\newcommand{\params}{\theta}         % parámetros θ

% ─── Funciones de activación ─────────────────────────────────────────────────
\DeclareMathOperator{\softmax}{softmax}
\DeclareMathOperator{\sigmoid}{\sigma}   % \sigmoid(x) = σ(x)
\DeclareMathOperator{\relu}{ReLU}
\DeclareMathOperator{\gelu}{GELU}

% ─── Score function y difusión ───────────────────────────────────────────────
% \score toma la variable como argumento opcional (default: \bm{x})
% Uso: \score  →  ∇_x log p(x)
%      \score[\bm{z}]  →  ∇_z log p(z)
\newcommand{\score}[1][\bm{x}]{\grad_{#1}\log p}

% Ecuación diferencial estocástica — versión sin color (inline, B/N)
% Para la versión con coloreado semántico usar \eqSDE de equations.tex.
% \SDE{X}{a(X,t)}{b(t)}  →  dX = a(X,t) dt + b(t) dW
\newcommand{\SDE}[3]{d#1 = #2\,dt + #3\,dW}


% =============================================================================
% FÍSICA COMPUTACIONAL
% =============================================================================

% ─── Mecánica cuántica ───────────────────────────────────────────────────────
% \bra, \ket, \braket, \comm, \acomm ya existen en physics.
\newcommand{\Ham}{\hat{H}}           % operador Hamiltoniano Ĥ
\newcommand{\Hscr}{\mathscr{H}}      % espacio de Hilbert ℋ
\newcommand{\Fscr}{\mathscr{F}}      % transformada de Fourier ℱ (operador)

% ─── Mecánica estadística ─────────────────────────────────────────────────────
\newcommand{\kB}{k_{\mathrm{B}}}     % constante de Boltzmann k_B
\newcommand{\Zpart}{Z}               % función de partición Z
\newcommand{\FP}{\mathcal{L}_{\mathrm{FP}}}  % operador Fokker-Planck ℒ_FP

% ─── Física de reactores ─────────────────────────────────────────────────────
\newcommand{\keff}{k_{\mathrm{eff}}} % factor de multiplicación efectivo


% =============================================================================
% SÍMBOLOS AUXILIARES
% =============================================================================
\newcommand{\OK}{\ensuremath{\checkmark}}   % ✓
\newcommand{\nOK}{\ensuremath{\times}}      % ✗


% =============================================================================
% NOTAS DE COMPATIBILIDAD
% =============================================================================
%
% ── Sobrescribir \rank de physics ────────────────────────────────────────────
% physics define \rank como operador. Para usar "rango" en español
% o cambiar el comportamiento, descomentar:
% \let\rank\relax
% \DeclareMathOperator{\rank}{rango}
%
% ── \D puede colisionar ──────────────────────────────────────────────────────
% Si algún paquete cargado define \D (poco común), usar \Dataset en su lugar:
% \let\D\relax
% \newcommand{\Dataset}{\mathcal{D}}
%
% ── \sigmoid vs variable σ ───────────────────────────────────────────────────
% \sigmoid usa \sigma internamente. Si σ ya está ocupada como variable en
% una sección, usar \sigma(x) directamente en lugar de \sigmoid(x).


% =============================================================================
\begin{document}
% =============================================================================

%───────────────────────────────────────────────────────────────────────────────
% TITLEPAGE.TEX — Portada genérica
%───────────────────────────────────────────────────────────────────────────────
% Consume: \docTitulo, \docSubtitulo, \docAutor, \docEstado, \docFecha,
%          \docCita (de metadata.tex)
%          titlepage-accent (de themes/catppuccin-palette.tex)
%
% NOTA — imagen del footer:
%   figs/gray0_ctp_on_line.svg debe convertirse a PDF antes de compilar:
%   inkscape --export-filename=figs/gray0_ctp_on_line.pdf \
%             figs/gray0_ctp_on_line.svg
%   LuaLaTeX no acepta SVG directamente con \includegraphics.
%───────────────────────────────────────────────────────────────────────────────
\begin{titlepage}
  \centering
  \vspace*{2cm}

  % ─── Título principal ──────────────────────────────────────────────────────
  {\fontsize{48}{60}\selectfont\bfseries
   \docTitulo\par}
  \vspace{0.8cm}

  % ─── Línea decorativa ──────────────────────────────────────────────────────
  {\color{titlepage-accent}\rule{0.6\textwidth}{0.5pt}}\par
  \vspace{1.2cm}

  % ─── Subtítulo ─────────────────────────────────────────────────────────────
  {\large\docSubtitulo\par}
  \vspace{2cm}

  % ─── Caja de cita ──────────────────────────────────────────────────────────
  \begin{tcolorbox}[
    enhanced,
    blanker,
    colback   = surface0,
    colframe  = titlepage-accent,
    coltext   = titlepage-accent,
    width     = 0.75\textwidth,
    halign    = center,
    left=6pt, right=6pt, top=6pt, bottom=6pt,
    arc       = 2pt,
    boxrule   = 0.3pt,
  ]
    \itshape\small\docCita
  \end{tcolorbox}

  \vfill

  % ─── Autor, estado y fecha ─────────────────────────────────────────────────
  {\normalsize\docAutor\par}
  \medskip
  {\normalsize\docEstado\ ---\ \docFecha\par}
  \smallskip
  {\footnotesize\color{overlay0}Compilado con \LaTeX}

  % ─── Footer: imagen decorativa ────────────────────────────────────────────
  % textblock* usa coordenadas absolutas en la unidad del segundo argumento.
  % (0cm, 0.93\paperheight) posiciona la esquina superior izquierda del bloque
  % al 93% de la altura del papel desde el borde superior.
  % overlay (opción de textpos) garantiza que no desplace el contenido normal.
  \begin{textblock*}{\paperwidth}(0cm, 0.93\paperheight)
    \noindent
    \includegraphics[width=0.7\paperwidth]{figs/gray0_ctp_on_line.pdf}
  \end{textblock*}

\end{titlepage}

%───────────────────────────────────────────────────────────────────────────────
% PÁGINAS PRELIMINARES: DEDICATORIA, EPÍGRAFE Y AGRADECIMIENTOS
%───────────────────────────────────────────────────────────────────────────────
\ifx\docDedicatoria\empty\else
  \cleardoublepage
  \vspace*{\fill}
  \begin{center}
    \itshape\large \docDedicatoria
  \end{center}
  \vspace*{\fill}
  \newpage
\fi

\ifx\docEpigrafe\empty\else
  \cleardoublepage
  \vspace*{0.3\textheight}
  \begin{flushright}
    \itshape\large \docEpigrafe
  \end{flushright}
  \vspace*{\fill}
  \newpage
\fi

\ifx\docAgradecimientos\empty\else
  \chapter*{Agradecimientos}
  \addcontentsline{toc}{chapter}{Agradecimientos}
  \docAgradecimientos
  \newpage
\fi
%───────────────────────────────────────────────────────────────────────────────
% PREFACIO (texto completo)
%───────────────────────────────────────────────────────────────────────────────
\chapter*{Prefacio}
\addcontentsline{toc}{chapter}{Prefacio}

\begin{flushright}
\itshape
``El aprendizaje automático es física aplicada a altas dimensiones.''
\end{flushright}
\bigskip

El universo no computa con certeza. Un reactor nuclear no produce exactamente
la misma traza de potencia dos veces...

% ... (todo el resto del prefacio) ...

\vspace{1.2cm}
\begin{flushright}
\textit{J.\,A. Chala-Casanova}\\[2pt]
\the\year
\end{flushright}

\tableofcontents

%───────────────────────────────────────────────────────────────────────────────
\chapter{Primeros pasos}
\label{ch:primeros-pasos}
%───────────────────────────────────────────────────────────────────────────────

Esta guía se escribe a sí misma: cada elemento que aparece en estas páginas
es una instancia real del componente que describe. Los entornos coloreados
que ves al leer son exactamente los que usarás al escribir. Las ecuaciones
están coloreadas con el sistema semántico del template. El código de listings
usa los estilos definidos en \texttt{config/listings.tex}.

El template \texttt{latex-catppuccin-academic} está diseñado para documentos
académicos de física computacional, matemáticas y machine learning: notas
de clase, tesis, apuntes avanzados. Su principio central es que la riqueza
visual no es decoración --- es información. El color de un entorno dice qué
rol epistémico tiene. El color de un símbolo dice si es operador, incógnita
o constante. El sistema crece con cada documento que lo usa.

%───────────────────────────────────────────────────────────────────────────────
\section{Requisitos del sistema}
\label{sec:requisitos}
%───────────────────────────────────────────────────────────────────────────────

\begin{definicion}{Motor de compilación: LuaLaTeX}
\textbf{LuaLaTeX}~\cite{luatex2023} es un motor TeX que extiende pdfLaTeX
con soporte nativo para UTF-8, acceso a fuentes del sistema vía
\texttt{fontspec}, y un intérprete Lua embebido. Este template lo requiere
por dos razones concretas: el paquete \texttt{catppuccinpalette} que gestiona
la paleta de colores~\cite{catppuccin2021} está escrito con primitivas
LuaLaTeX, y el soporte completo de caracteres Unicode en bloques de código
(\texttt{listings}) depende del motor.
\end{definicion}

\begin{observacion}{Distribución y versión mínima}
Se recomienda \textbf{TeX Live 2022} o superior (o MiKTeX equivalente).
El paquete \texttt{catppuccinpalette} es relativamente reciente --- versiones
anteriores de la distribución pueden no incluirlo. Si
\texttt{tlmgr install catppuccinpalette} devuelve error, actualiza con
\texttt{tlmgr update -{}-all}.
\end{observacion}

Los paquetes que carga el template están listados y comentados en
\texttt{config/packages.tex}. Incluyen entre otros \texttt{pgfplots}~\cite{pgfplots2023} para gráficas y \texttt{tcolorbox}~\cite{tcolorbox2024} para cajas. Todos forman parte de TeX Live y no requieren
instalaciones adicionales.

%───────────────────────────────────────────────────────────────────────────────
\section{Estructura del proyecto}
\label{sec:estructura}
%───────────────────────────────────────────────────────────────────────────────

El template sigue una arquitectura modular estricta: un archivo por
responsabilidad, agrupados en carpetas por tipo.

\begin{center}
\ttfamily\scriptsize\color{maintext}
\begin{verbatim}
.
|-- main.tex
|-- metadata.tex
|-- config/
|   |-- packages.tex    |-- geometry.tex    |-- hyperref.tex
|   |-- tcolorbox.tex   |-- titles.tex      |-- captions.tex
|   |-- tables.tex      |-- listings.tex    |-- algorithms.tex
|   |-- equations.tex   |-- pgfplots.tex    \-- notation.tex
|-- environments/
|   |-- generator.tex   |-- instances.tex   \-- specials.tex
|-- themes/
|   \-- catppuccin-palette.tex
|-- frontmatter/
|   |-- titlepage.tex   |-- preliminary.tex  \-- preface.tex
|-- chapters/           |-- backmatter/       |-- figs/
|-- references/
|   \-- references.bib  (+ chapter1.bib, chapter2.bib, ...)
\end{verbatim}
\end{center}

\begin{observacion}{Principio de nomenclatura}
El nombre de cada archivo refleja su sujeto, no su tipo. \texttt{tables.tex}
en \texttt{config/} configura tablas; \texttt{equations.tex} configura el
coloreado semántico; \texttt{notation.tex} define el vocabulario matemático.
El sufijo \texttt{-config} es redundante cuando la carpeta ya comunica el
contexto.
\end{observacion}

%───────────────────────────────────────────────────────────────────────────────
\section{Compilación}
\label{sec:compilacion}
%───────────────────────────────────────────────────────────────────────────────

La secuencia de compilación típica para un documento con bibliografía:

\begin{lstlisting}[style=latex, frame=none, numbers=none]
lualatex main.tex   # primera pasada: genera .aux y .bcf
biber main          # procesa la bibliografia
lualatex main.tex   # segunda pasada: inserta referencias
lualatex main.tex   # tercera pasada: estabiliza TOC y cross-refs
\end{lstlisting}

La gestión de bibliografía usa \texttt{biblatex}~\cite{biblatex2024} con backend \texttt{biber}. Para documentos sin bibliografía, dos pasadas son suficientes. Los archivos
intermedios (\texttt{.aux}, \texttt{.log}, \texttt{.bcf}, \texttt{.synctex.gz})
no forman parte del fuente y pueden excluirse con \texttt{.gitignore}.

\begin{observacion}{Compilación continua}
La mayoría de editores (TeXstudio, VS Code con LaTeX Workshop, Neovim con
vimtex) permiten compilación automática al guardar. Configurar la cadena
como \texttt{lualatex -> biber -> lualatex x2}. En línea de comandos,
\texttt{latexmk -lualatex -interaction=nonstopmode main.tex} determina
automáticamente el número de pasadas necesarias.
\end{observacion}

%───────────────────────────────────────────────────────────────────────────────
\section{Eligiendo un flavor}
\label{sec:flavor}
%───────────────────────────────────────────────────────────────────────────────

\begin{definicion}{Flavor de Catppuccin}
Catppuccin~\cite{catppuccin-styleguide} define cuatro \textbf{flavors}:
variantes de la misma paleta de 26 colores con distintas temperaturas de
fondo. \textbf{Latte} es el único flavor claro; \textbf{Frappé},
\textbf{Macchiato} y \textbf{Mocha} son progresivamente más oscuros. Todos
los colores semánticos del template son invariantes entre flavors --- el azul
siempre es el color de referencias internas, el malva siempre es el de
teoremas --- pero sus valores hexadecimales se adaptan para mantener contraste
y legibilidad en cada temperatura de fondo.
\end{definicion}

El flavor se elige en un único lugar:
\texttt{themes/catppuccin-palette.tex}. Descomentar exactamente una línea
y ajustar el toggle de tema oscuro:

\begin{lstlisting}[style=latex, frame=none, numbers=none]
\usepackage[Latte,styleAll]{catppuccinpalette}       % claro
%\usepackage[Frappe,styleAll]{catppuccinpalette}
%\usepackage[Macchiato,styleAll]{catppuccinpalette}
%\usepackage[Mocha,styleAll]{catppuccinpalette}

\darkthemefalse    % Latte
%\darkthemetrue    % Frappe, Macchiato o Mocha
\end{lstlisting}

El toggle \texttt{darktheme} controla el color de fondo de página y la
generación de fondos de cajas: en tema claro son \texttt{white!97!color},
en tema oscuro son \texttt{CtpBase!95!color}. Esto garantiza que los
entornos teoremáticos sean legibles en cualquier flavor.

\begin{observacion}{Sistema de color en tres capas}
\texttt{Ctp*} (valores absolutos del flavor) solo aparece en
\texttt{themes/catppuccin-palette.tex}. Los archivos de config consumen
aliases semánticos (\texttt{env-def}, \texttt{doc-chapter},
\texttt{code-keyword}). Los archivos de contenido consumen comandos
(\texttt{\textbackslash eqop}, \texttt{\textbackslash begin\{teorema\}}).
Ningún archivo viola la capa que le corresponde. Cuando ese invariante
se mantiene, cambiar el flavor del documento entero es descomentar
una línea.
\end{observacion}

%───────────────────────────────────────────────────────────────────────────────
\section{Configurando los metadatos}
\label{sec:metadatos}
%───────────────────────────────────────────────────────────────────────────────

Todos los datos específicos del documento viven en \texttt{metadata.tex}.
Este archivo se carga antes del tema para que \texttt{hyperref} pueda usar
las variables al generar los metadatos del PDF.

\begin{acadtable}[H]{tab:metadata-campos}{Campos de \texttt{metadata.tex}}
  \begin{tabularx}{\textwidth}{l l X}
    \toprule
    \headerrow
    \thead{Comando} & \thead{Usado en} & \thead{Notas} \\
    \midrule
    \texttt{\textbackslash docTitulo}
      & Portada, PDF title    & Sin \texttt{\textbackslash\textbackslash} \\
    \texttt{\textbackslash docSubtitulo}
      & Portada               & Admite \texttt{\textbackslash\textbackslash} \\
    \texttt{\textbackslash docAutor}
      & Portada, PDF author   & Sin \texttt{\textbackslash\textbackslash} \\
    \texttt{\textbackslash docFecha}
      & Portada               & Default: año de compilación \\
    \texttt{\textbackslash docVersion}
      & Portada               & Versionado semántico \\
    \texttt{\textbackslash docEstado}
      & Portada               & Borrador \textbar{} Final \textbar{} etc. \\
    \midrule
    \texttt{\textbackslash docCita}
      & Portada               & Entre comillas tipográficas \\
    \texttt{\textbackslash docDedicatoria}
      & Preliminary           & Texto breve \\
    \texttt{\textbackslash docEpigrafe}
      & Preliminary           & Entre comillas tipográficas \\
    \texttt{\textbackslash docAgradecimientos}
      & Preliminary           & Párrafo libre \\
    \midrule
    \texttt{\textbackslash docKeywords}
      & PDF keywords          & Separadas por coma \\
    \texttt{\textbackslash docSubject}
      & PDF subject           & Área temática \\
    \texttt{\textbackslash docLang}
      & PDF language          & Código ISO 639-1 \\
    \bottomrule
  \end{tabularx}
\end{acadtable}

\begin{observacion}{Metadatos PDF y modo matemático}
Los campos usados como metadatos PDF (\texttt{\textbackslash docTitulo},
\texttt{\textbackslash docAutor}, \texttt{\textbackslash docKeywords})
no admiten comandos matemáticos ni énfasis tipográfico. \texttt{hyperref}
los pasa al PDF como strings planos. Si el título contiene matemáticas,
usar \texttt{\textbackslash texorpdfstring\{\$E=mc\^{}2\$\}\{E=mc2\}} para
proveer una versión alternativa para el PDF.
\end{observacion}

%───────────────────────────────────────────────────────────────────────────────
\section{El preámbulo y su orden de carga}
\label{sec:preambulo}
%───────────────────────────────────────────────────────────────────────────────

El orden de los \texttt{\textbackslash input} en \texttt{main.tex} no es
arbitrario. Cada paso depende del anterior:

\begin{acadtable}[H]{tab:load-order}{Orden de carga del preámbulo}
  \begin{tabularx}{\textwidth}{c l X}
    \toprule
    \headerrow
    \thead{Paso} & \thead{Archivo} & \thead{Dependencia} \\
    \midrule
    1 & \texttt{config/packages.tex}
      & Carga todos los paquetes. El resto depende de ellos. \\
    2 & \texttt{metadata.tex}
      & Define \texttt{\textbackslash docTitulo} etc. antes de
        que hyperref los necesite. \\
    3 & \texttt{themes/catppuccin-palette.tex}
      & Define aliases de color antes de que alguien los consuma. \\
    4 & \texttt{config/hyperref.tex}
      & Consume aliases \texttt{link-*}. Debe ir después del tema. \\
    5 & \texttt{config/geometry.tex}
      & \texttt{fancyhdr} puede interactuar con anchors de hyperref. \\
    6 & \texttt{config/tcolorbox.tex}
      & Define \texttt{basebox}~\cite{tcolorbox2024}. Los entornos lo necesitan. \\
    7 & \texttt{environments/*.tex}
      & Orden obligatorio: generator $\to$ instances $\to$ specials. \\
    8 & \texttt{config/titles.tex} \ldots
      & Configs de dominio. Entre sí el orden es libre. \\
    9 & \texttt{config/notation.tex}
      & Puede usar cualquier paquete ya cargado. \\
    \bottomrule
  \end{tabularx}
\end{acadtable}

%───────────────────────────────────────────────────────────────────────────────
\printbibliography[heading=subbibliography]

%───────────────────────────────────────────────────────────────────────────────
\chapter{Estructura del documento}
\label{ch:estructura}
%───────────────────────────────────────────────────────────────────────────────

Un documento académico de largo aliento tiene estructura en múltiples escalas:
la jerarquía visual de títulos que orienta al lector, el sistema de referencias
cruzadas que da cohesión al argumento, el frontmatter que contextualiza la obra,
y la bibliografía que ancla cada afirmación. Este capítulo documenta cómo el
template gestiona cada una de esas escalas y cómo usarlas correctamente.

%───────────────────────────────────────────────────────────────────────────────
\section{Jerarquía de títulos}
\label{sec:titulos}
%───────────────────────────────────────────────────────────────────────────────

El template configura cuatro niveles de título mediante
\texttt{titlesec}~\cite{titlesec2021}, cada uno con un color semántico
distinto definido en la paleta bajo el prefijo \texttt{doc-*}:

\begin{acadtable}[H]{tab:titulo-niveles}{Jerarquía de títulos y sus colores}
  \begin{tabularx}{\textwidth}{l l l X}
    \toprule
    \headerrow
    \thead{Comando} & \thead{Alias de color} & \thead{Color Ctp} & \thead{Uso} \\
    \midrule
    \texttt{\textbackslash chapter}
      & \texttt{doc-chapter} & Pink
      & Unidad temática mayor \\
    \texttt{\textbackslash section}
      & \texttt{doc-section} & Sky
      & Subtema dentro del capítulo \\
    \texttt{\textbackslash subsection}
      & \texttt{doc-subsection} & Lavender
      & División de una sección \\
    \texttt{\textbackslash subsubsection}
      & \texttt{doc-subsubsection} & Peach
      & Detalle dentro de subsección \\
    \bottomrule
  \end{tabularx}
\end{acadtable}

La progresión de colores no es arbitraria: contraste decreciente con la
profundidad. El capítulo ancla visualmente la página con el Pink más
prominente; la subsubsección aporta apenas un tono Peach que diferencia
sin competir. El lector percibe la jerarquía antes de leer.

\begin{observacion}{Cambiar colores de títulos}
Los colores se modifican exclusivamente en
\texttt{themes/catppuccin-palette.tex}, cambiando el alias correspondiente.
Por ejemplo, para que las secciones usen Mauve en lugar de Sky:
\texttt{\textbackslash colorlet\{doc-section\}\{CtpMauve\}}.
Nunca editar \texttt{config/titles.tex} para poner colores directamente.
\end{observacion}

\subsection{Títulos sin numeración}
\label{subsec:titulos-nonumber}

Los comandos con asterisco (\texttt{\textbackslash section*},
\texttt{\textbackslash subsection*}) suprimen la numeración pero
mantienen el color y la tipografía del nivel correspondiente. Son útiles
para secciones estructurales como prefacios o glosarios que no deben
aparecer numeradas en el texto pero sí en la tabla de contenidos:

\begin{lstlisting}[style=latex, frame=none, numbers=none]
\section*{Notación}
\addcontentsline{toc}{section}{Notación}  % añadir al TOC manualmente
\end{lstlisting}

\subsubsection{Párrafos como cuarto nivel}
\label{subsubsec:parrafos}

El cuarto nivel visible en el texto es \texttt{\textbackslash subsubsection}.
Existe un quinto nivel, \texttt{\textbackslash paragraph}, que aparece como
texto inline sin salto de línea propio. Su uso debe ser esporádico: si
necesitas un quinto nivel estructural con frecuencia, la estructura del
documento probablemente necesita reorganización.

%───────────────────────────────────────────────────────────────────────────────
\section{Sistema de referencias cruzadas}
\label{sec:referencias-cruzadas}
%───────────────────────────────────────────────────────────────────────────────

El template usa \texttt{cleveref}~\cite{cleveref2022} para referencias
inteligentes. A diferencia de \texttt{\textbackslash ref}, que produce
solo el número, \texttt{\textbackslash cref} produce el tipo y el número
correctos en el idioma configurado:

\begin{definicion}{Referencias cruzadas con \texttt{cleveref}}
\texttt{\textbackslash cref\{etiqueta\}} produce la referencia completa
en minúscula: \textit{ec.~(2.1)}, \textit{fig.~3.2}, \textit{def.~1.4}.
\texttt{\textbackslash Cref\{etiqueta\}} produce la versión capitalizada
para inicio de oración: \textit{Ecuación~(2.1)}, \textit{Figura~3.2}.
Para rangos, \texttt{\textbackslash crefrange\{etiq1\}\{etiq2\}} produce
\textit{ecs.~(2.1)--(2.3)} automáticamente.
\end{definicion}

\subsection{Convención de etiquetas}
\label{subsec:etiquetas}

Una convención consistente de prefijos hace el código fuente legible y
evita colisiones entre etiquetas de distintos archivos:

\begin{acadtable}[H]{tab:label-prefijos}{Prefijos de etiquetas recomendados}
  \begin{tabularx}{\textwidth}{l l X}
    \toprule
    \headerrow
    \thead{Tipo} & \thead{Prefijo} & \thead{Ejemplo} \\
    \midrule
    Capítulo      & \texttt{ch:}   & \texttt{ch:estructura} \\
    Sección       & \texttt{sec:}  & \texttt{sec:referencias-cruzadas} \\
    Subsección    & \texttt{subsec:} & \texttt{subsec:etiquetas} \\
    Ecuación      & \texttt{eq:}   & \texttt{eq:schrodinger} \\
    Figura        & \texttt{fig:}  & \texttt{fig:convergencia} \\
    Tabla         & \texttt{tab:}  & \texttt{tab:label-prefijos} \\
    Algoritmo     & \texttt{alg:}  & \texttt{alg:gradient} \\
    Código        & \texttt{lst:}  & \texttt{lst:perceptron} \\
    Apéndice      & \texttt{ap:}   & \texttt{ap:galeria} \\
    Definición    & \texttt{def:}  & \texttt{def:hilbert} \\
    Teorema       & \texttt{thm:}  & \texttt{thm:riesz} \\
    \bottomrule
  \end{tabularx}
\end{acadtable}

\subsection{Referencias a entornos propios}
\label{subsec:cref-entornos}

Los entornos teoremáticos del template crean contadores propios. Para que
\texttt{cleveref} los reconozca, sus contadores están registrados en
\texttt{config/packages.tex}. El patrón de etiqueta es consistente con
los prefijos de la \cref{tab:label-prefijos}: poner el
\texttt{\textbackslash label} dentro del entorno, después del título:

\begin{lstlisting}[style=latex, frame=none, numbers=none]
\begin{definicion}{Espacio de Hilbert}
\label{def:hilbert}
Un espacio de Hilbert es...
\end{definicion}

% Más adelante en el texto:
Como se establece en la \cref{def:hilbert}, un espacio de Hilbert...
\end{lstlisting}

\begin{observacion}{Posición del \texttt{\textbackslash label}}
En entornos de \texttt{tcolorbox}, el \texttt{\textbackslash label} debe
ir \textit{dentro} del entorno, no en el argumento del título. El título
es un argumento de texto, no un entorno flotante. Poner el label fuera
produce un anchor de hyperref incorrecto.
\end{observacion}

%───────────────────────────────────────────────────────────────────────────────
\section{Frontmatter}
\label{sec:frontmatter}
%───────────────────────────────────────────────────────────────────────────────

El frontmatter agrupa las páginas previas al contenido. El template lo
divide en tres archivos en \texttt{frontmatter/}, cada uno con una
responsabilidad única:

\paragraph{\texttt{titlepage.tex} --- Portada.}
Genera la portada completa a partir de las variables de \texttt{metadata.tex}.
No requiere edición directa salvo cambios de diseño estructural. Los contenidos
se controlan exclusivamente desde \texttt{metadata.tex}.

\paragraph{\texttt{preliminary.tex} --- Páginas preliminares.}
Muestra dedicatoria, epígrafe y agradecimientos en páginas separadas.
Cada sección es condicional: si la variable correspondiente está vacía
(\texttt{\textbackslash docDedicatoria\{\}}), la página no se genera.
Esto permite activar o desactivar secciones sin editar el archivo.

\paragraph{\texttt{preface.tex} --- Prefacio.}
Capítulo sin numerar (\texttt{\textbackslash chapter*}) añadido al TOC
manualmente. A diferencia de los preliminares, el prefacio tiene contenido
narrativo propio --- motivación del trabajo, agradecimientos extendidos,
contexto histórico --- y se edita directamente como cualquier capítulo.

\begin{observacion}{Paginación del frontmatter}
Por convención tipográfica, el frontmatter usa numeración romana
(\texttt{\textbackslash pagenumbering\{roman\}}) y el cuerpo del documento
numeración arábiga (\texttt{\textbackslash pagenumbering\{arabic\}}).
El template aplica esto en \texttt{main.tex}. Si quieres que el TOC empiece
en la página~i, asegúrate de que el cambio de numeración ocurra antes de
\texttt{\textbackslash tableofcontents}.
\end{observacion}

%───────────────────────────────────────────────────────────────────────────────
\section{Tabla de contenidos}
\label{sec:toc}
%───────────────────────────────────────────────────────────────────────────────

\texttt{\textbackslash tableofcontents} genera el TOC automáticamente a
partir de los títulos del documento. El nivel de profundidad se controla
con \texttt{tocdepth}:

\begin{lstlisting}[style=latex, frame=none, numbers=none]
\setcounter{tocdepth}{2}    % muestra hasta \subsection
\setcounter{tocdepth}{3}    % muestra hasta \subsubsection (default)
\end{lstlisting}

El template configura \texttt{secnumdepth} en 3 por defecto --- todos los
niveles hasta \texttt{\textbackslash subsubsection} están numerados. Si una
sección con asterisco debe aparecer en el TOC (como el prefacio), se añade
manualmente:

\begin{lstlisting}[style=latex, frame=none, numbers=none]
\chapter*{Prefacio}
\addcontentsline{toc}{chapter}{Prefacio}
\end{lstlisting}

%───────────────────────────────────────────────────────────────────────────────
\section{Bibliografía por capítulo}
\label{sec:bibliografia}
%───────────────────────────────────────────────────────────────────────────────

\begin{definicion}{Refsection de \texttt{biblatex}}
Con la opción \texttt{refsection=chapter}, \texttt{biblatex} trata cada
\texttt{\textbackslash chapter} (o cada \texttt{\textbackslash include})
como una sección bibliográfica independiente. Las citas de un capítulo
solo aparecen en la bibliografía de ese capítulo.
\texttt{\textbackslash printbibliography[heading=subbibliography]} al
final de cada capítulo imprime solo las referencias citadas en ese capítulo,
con numeración reiniciada.
\end{definicion}

El patrón a seguir en cada capítulo nuevo es:

\begin{lstlisting}[style=latex, frame=none, numbers=none]
% 1. Crear references/chapterN.bib con las referencias del capitulo
% 2. En metadata.tex, anadir:
\addbibresource{references/chapterN.bib}

% 3. Al final del archivo chapterN.tex:
\printbibliography[heading=subbibliography]
\end{lstlisting}

\begin{observacion}{Referencias compartidas entre capítulos}
Si una referencia aparece en varios capítulos, debe estar en el
\texttt{.bib} del capítulo donde se cita por primera vez, o en el
\texttt{references.bib} global. \texttt{biblatex} admite múltiples
\texttt{\textbackslash addbibresource} y no duplica entradas aunque
la misma clave aparezca en varios archivos.
\end{observacion}

%───────────────────────────────────────────────────────────────────────────────
\section{Apéndices}
\label{sec:apendices}
%───────────────────────────────────────────────────────────────────────────────

Los apéndices~\cite{appendix2020} usan el entorno \texttt{appendices}
provisto por el paquete homónimo. Dentro de él, cada apéndice es un
\texttt{\textbackslash chapter} normal cuya numeración cambia
automáticamente a letras (A, B, C\ldots):

\begin{lstlisting}[style=latex, frame=none, numbers=none]
% En main.tex, después de los capítulos:
\begin{appendices}
  %───────────────────────────────────────────────────────────────────────────────
\chapter{Style Guide de Catppuccin Académico}
\label{ap:styleguide}
%───────────────────────────────────────────────────────────────────────────────

Este apéndice es el documento normativo del sistema de color del template.
Define la arquitectura, las reglas de decisión, el inventario completo de
aliases y el protocolo de extensión. Todo color nuevo que se añada al
template debe pasar por este documento primero.

%───────────────────────────────────────────────────────────────────────────────
\section{Arquitectura de tres capas}
\label{ap:sec:arquitectura}
%───────────────────────────────────────────────────────────────────────────────

\begin{definicion}{Las tres capas del sistema de color}
\label{ap:def:tres-capas}
El sistema de color opera en tres niveles de abstracción con una regla
de visibilidad estricta:

\textbf{Nivel 0 — Paleta cruda (\texttt{Ctp*}):} Valores hexadecimales
absolutos del flavor activo. Existen 26 colores: \texttt{CtpRosewater},
\texttt{CtpFlamingo}, \texttt{CtpPink}, \texttt{CtpMauve}, \texttt{CtpRed},
\texttt{CtpMaroon}, \texttt{CtpPeach}, \texttt{CtpYellow}, \texttt{CtpGreen},
\texttt{CtpTeal}, \texttt{CtpSky}, \texttt{CtpSapphire}, \texttt{CtpBlue},
\texttt{CtpLavender}, \texttt{CtpText}, \texttt{CtpSubtext1},
\texttt{CtpSubtext0}, \texttt{CtpOverlay2}, \texttt{CtpOverlay1},
\texttt{CtpOverlay0}, \texttt{CtpSurface2}, \texttt{CtpSurface1},
\texttt{CtpSurface0}, \texttt{CtpBase}, \texttt{CtpMantle},
\texttt{CtpCrust}.

\textbf{Nivel 1 — Aliases estructurales:} Roles de UI genéricos sin
semántica de dominio. Ejemplos: \texttt{background}, \texttt{surface0},
\texttt{overlay0}, \texttt{maintext}. Describen propiedades visuales de
elementos de interfaz.

\textbf{Nivel 2 — Aliases semánticos:} Roles de dominio del documento.
Ejemplos: \texttt{env-def}, \texttt{doc-chapter}, \texttt{code-keyword}.
Describen la intención del elemento en el documento.
\end{definicion}

La regla de visibilidad es \textbf{absoluta}:

\begin{estrella}{Invariante de capas}
\label{ap:star:invariante}
\texttt{Ctp*} solo aparece en \texttt{themes/catppuccin-palette.tex}.\\
Los archivos de \texttt{config/} consumen aliases semánticos.\\
Los archivos de contenido (\texttt{chapters/}, \texttt{environments/})
consumen comandos del template.

Cuando este invariante se mantiene, cambiar el flavor del documento
entero es descomentar una línea en \texttt{catppuccin-palette.tex}.
Cuando se viola, el cambio de flavor requiere buscar y editar en
múltiples archivos.
\end{estrella}

%───────────────────────────────────────────────────────────────────────────────
\section{Árbol de decisión para asignar colores}
\label{ap:sec:decision}
%───────────────────────────────────────────────────────────────────────────────

Cuando se necesita un color nuevo en cualquier archivo de config o de
contenido, la secuencia de preguntas es:

\begin{enumerate}[label=\textbf{\arabic*.}, leftmargin=2em]
  \item \textbf{¿Estoy en \texttt{catppuccin-palette.tex}?}
    Si sí, puedo usar \texttt{Ctp*} libremente para definir aliases.
    Si no, nunca uso \texttt{Ctp*}: avanzo al paso 2.

  \item \textbf{¿Ya existe un alias semántico para este rol?}
    Buscar en la paleta con el prefijo del dominio (\texttt{env-},
    \texttt{code-}, \texttt{tikz-}, etc.). Si existe, lo consumo.
    Si no existe, avanzo al paso 3.

  \item \textbf{¿El color tiene semántica de dominio o es infraestructura?}
    Si es semántica de dominio (``el borde de un estado aceptador de
    autómata'', ``el color de una keyword''), añado un alias semántico
    nuevo en la paleta bajo el prefijo correspondiente.
    Si es infraestructura genérica (``un fondo de superficie'', ``texto
    secundario''), uso un alias estructural existente (\texttt{surface0},
    \texttt{subtext0}).

  \item \textbf{¿Es un cambio independiente de otros elementos?}
    Si el color de este elemento siempre cambia junto con otro, comparten
    alias. Si pueden cambiar independientemente, tienen aliases separados.
\end{enumerate}

\begin{observacion}{La regla de granularidad mínima}
Un elemento nuevo merece su propio alias semántico si y solo si su
cambio de color es \textit{independiente} de otros elementos. Si el
label de figura y el label de tabla siempre deben ser del mismo color,
comparten \texttt{caption-label}. Si pueden diferir (como es el caso en
este template), tienen \texttt{caption-fig} y \texttt{caption-tab}
separados.
\end{observacion}

%───────────────────────────────────────────────────────────────────────────────
\section{Convención de nombres}
\label{ap:sec:nombres}
%───────────────────────────────────────────────────────────────────────────────

\begin{definicion}{Convención kebab-case}
\label{ap:def:kebab}
Todos los aliases siguen la convención:
\[
  \texttt{\{dominio\}-\{elemento\}[-\{variante\}]}
\]
en kebab-case. El dominio identifica la categoría; el elemento describe
el rol dentro del dominio; la variante (\texttt{frame}, \texttt{back},
\texttt{text}) solo aparece cuando el elemento tiene múltiples roles de
color. Los aliases estructurales (Nivel 1) no llevan prefijo de dominio.
\end{definicion}

\begin{acadtable}[H]{ap:tab:dominios}{Dominios activos y sus prefijos}
  \begin{tabularx}{\textwidth}{l l X}
    \toprule
    \headerrow
    \thead{Prefijo} & \thead{Dominio} & \thead{Cubre} \\
    \midrule
    (sin prefijo)    & Infraestructura     & Fondos, superficies, texto base \\
    \texttt{doc-}    & Documento           & Jerarquía de títulos \\
    \texttt{link-}   & Hipervínculos       & Referencias, URLs, citas \\
    \texttt{env-}    & Entornos            & Definición, teorema, ejemplo, etc. \\
    \texttt{eq-}     & Ecuaciones          & Coloreado semántico de matemáticas \\
    \texttt{code-}   & Código fuente       & Listings, syntax highlighting \\
    \texttt{alg-}    & Algoritmos          & algorithm2e, pseudocódigo \\
    \texttt{plot-}   & Gráficas            & pgfplots, series, ejes \\
    \texttt{tikz-}   & Diagramas TikZ      & Nodos, aristas, autómatas \\
    \texttt{table-}  & Tablas              & Encabezados, filas, reglas \\
    \texttt{caption-}& Captions            & Labels por tipo de flotante \\
    \texttt{fb-}     & Feedback            & Éxito, error, advertencia, info \\
    \texttt{diff-}   & Diffs               & Cambios insertados/eliminados \\
    \texttt{rainbow-}& Secuencia           & N colores distinguibles ordenados \\
    \texttt{frontmatter-} & Portada        & Elementos de la página de título \\
    \bottomrule
  \end{tabularx}
\end{acadtable}

%───────────────────────────────────────────────────────────────────────────────
\section{Inventario completo de aliases}
\label{ap:sec:inventario}
%───────────────────────────────────────────────────────────────────────────────

\subsection{Infraestructura visual (Nivel 1)}
\label{ap:subsec:infraestructura}

\begin{acadtable}[H]{ap:tab:infra}{Aliases estructurales: fondos, superficies y tipografía}
  \begin{tabularx}{\textwidth}{l l X}
    \toprule
    \headerrow
    \thead{Alias} & \thead{Color Ctp} & \thead{Fundamento Style Guide} \\
    \midrule
    \texttt{background}   & Base      & Background Pane \\
    \texttt{secondarybg}  & Mantle    & Secondary Panes \\
    \texttt{crustbg}      & Crust     & Secondary Panes (más oscuro) \\
    \texttt{surface0/1/2} & Surface 0/1/2 & Surface Elements \\
    \texttt{overlay0/1/2} & Overlay 0/1/2 & Overlays \\
    \texttt{maintext}     & Text      & Body Copy \\
    \texttt{subtext0/1}   & Subtext 0/1 & Sub-Headlines, Labels \\
    \texttt{subtletext}   & Overlay 1 & Subtle — texto terciario \\
    \bottomrule
  \end{tabularx}
\end{acadtable}

\subsection{Estructura del documento (\texttt{doc-})}

\begin{acadtable}[H]{ap:tab:doc}{Aliases de jerarquía de títulos}
  \begin{tabularx}{\textwidth}{l l X}
    \toprule
    \headerrow
    \thead{Alias} & \thead{Color Ctp} & \thead{Razón} \\
    \midrule
    \texttt{doc-chapter}       & Pink      & Máximo contraste; warm; ancla el capítulo \\
    \texttt{doc-section}       & Sky       & Cool; segundo nivel; contrasta con Pink \\
    \texttt{doc-subsection}    & Lavender  & Suave; tercer nivel \\
    \texttt{doc-subsubsection} & Peach     & Warm suave; cuarto nivel \\
    \texttt{doc-paragraph}     & Subtext0  & Casi texto; diferencia mínima \\
    \bottomrule
  \end{tabularx}
\end{acadtable}

\subsection{Hipervínculos (\texttt{link-})}

Mapeo directo del Catppuccin Style Guide: Links=Blue, Followed=Lavender, Hover=Sky.

\begin{acadtable}[H]{ap:tab:link}{Aliases de hipervínculos}
  \begin{tabularx}{\textwidth}{l l X}
    \toprule
    \headerrow
    \thead{Alias} & \thead{Color Ctp} & \thead{Uso en hyperref} \\
    \midrule
    \texttt{link-internal} & Blue     & \texttt{linkcolor}: refs internas, TOC \\
    \texttt{link-external} & Peach    & \texttt{urlcolor}: URLs externas \\
    \texttt{link-cite}     & Green    & \texttt{citecolor}: citas bibliográficas \\
    \texttt{link-visited}  & Lavender & Followed Links \\
    \texttt{link-hover}    & Sky      & On Hover Links \\
    \bottomrule
  \end{tabularx}
\end{acadtable}

\subsection{Entornos teoremáticos (\texttt{env-})}

Cada entorno tiene \texttt{env-X} (frame) y \texttt{env-X-back} (fondo,
generado por \texttt{\textbackslash createback\{env-X\}}).

\begin{acadtable}[H]{ap:tab:env}{Aliases de entornos teoremáticos}
  \begin{tabularx}{\textwidth}{l l X}
    \toprule
    \headerrow
    \thead{Alias} & \thead{Color Ctp} & \thead{Entorno y razón} \\
    \midrule
    \texttt{env-def}  & Blue     & Definición: epistémico, fundacional \\
    \texttt{env-thm}  & Mauve    & Teorema: el más prominente del sistema \\
    \texttt{env-prop} & Green    & Proposición: verdad demostrada \\
    \texttt{env-ex}   & Peach    & Ejemplo: warm, invita a leer \\
    \texttt{env-obs}  & Overlay0 & Observación: secundaria, no compite \\
    \texttt{env-phys} & Sapphire & Conexión física: evoca ciencia \\
    \texttt{env-star} & Yellow   & Estrella Polar: máxima prominencia \\
    \texttt{env-comp} & Pink     & Ejercicio Computacional: activo \\
    \bottomrule
  \end{tabularx}
\end{acadtable}

\subsection{Matemáticas semánticas (\texttt{eq-})}

Tres ejes ortogonales. Ver \cref{ch:ecuaciones} para la documentación completa.

\begin{acadtable}[H]{ap:tab:eq}{Aliases de coloreado semántico de ecuaciones}
  \begin{tabularx}{\textwidth}{l l l X}
    \toprule
    \headerrow
    \thead{Alias} & \thead{Macro} & \thead{Color Ctp} & \thead{Eje y rol} \\
    \midrule
    \texttt{eq-op}      & \texttt{\textbackslash eqop}  & Blue      & Sintáctico: operadores \\
    \texttt{eq-fn}      & \texttt{\textbackslash eqfn}  & Green     & Sintáctico: funciones/mapas \\
    \texttt{eq-dm}      & \texttt{\textbackslash eqdm}  & Teal      & Sintáctico: diferenciales \\
    \texttt{eq-dom}     & \texttt{\textbackslash eqdom} & Sapphire  & Sintáctico: dominios/conjuntos \\
    \midrule
    \texttt{eq-unknown} & \texttt{\textbackslash equ}   & Peach     & Epistémico: incógnita \\
    \texttt{eq-known}   & \texttt{\textbackslash eqk}   & Mauve     & Epistémico: dato/fijado \\
    \texttt{eq-const}   & \texttt{\textbackslash eqcon} & Lavender  & Epistémico: constante universal \\
    \midrule
    \texttt{eq-vec}     & \texttt{\textbackslash eqvec} & Red       & Tipo: vector/campo \\
    \texttt{eq-ten}     & \texttt{\textbackslash eqten} & Rosewater & Tipo: tensor/matriz \\
    \bottomrule
  \end{tabularx}
\end{acadtable}

\subsection{Código fuente (\texttt{code-})}

Mapeo directo del Catppuccin Style Guide: Language Defaults (Code Editors).

\begin{acadtable}[H]{ap:tab:code}{Aliases de syntax highlighting}
  \begin{tabularx}{\textwidth}{l l X}
    \toprule
    \headerrow
    \thead{Alias} & \thead{Color Ctp} & \thead{Clase sintáctica} \\
    \midrule
    \texttt{code-keyword} & Mauve    & Keywords del lenguaje \\
    \texttt{code-string}  & Green    & Strings y literales \\
    \texttt{code-comment} & Overlay2 & Comentarios \\
    \texttt{code-operator}& Sky      & Operadores \\
    \texttt{code-type}    & Yellow   & Tipos / clases / interfaces \\
    \texttt{code-fn}      & Blue     & Funciones / métodos \\
    \texttt{code-mod}     & Teal     & Módulos / namespaces / enums \\
    \texttt{code-macro}   & Rosewater& Macros / decoradores / pragmas \\
    \texttt{code-builtin} & Red      & Builtins / atoms / symbols \\
    \texttt{code-const}   & Peach    & Constantes / números \\
    \texttt{code-param}   & Maroon   & Parámetros \\
    \texttt{code-delim}   & Overlay2 & Delimitadores (\texttt{\{\}[]()}) \\
    \texttt{code-bg}      & Base     & Fondo del bloque \\
    \texttt{code-border}  & Surface0 & Frame / regla \\
    \texttt{code-linenum} & Overlay1 & Números de línea \\
    \bottomrule
  \end{tabularx}
\end{acadtable}

\subsection{Algoritmos, plots, TikZ, tablas, captions y feedback}

Las tablas de aliases para los dominios \texttt{alg-}, \texttt{plot-},
\texttt{tikz-}, \texttt{table-}, \texttt{caption-} y \texttt{fb-}
se presentan en las secciones correspondientes de los capítulos:
\cref{ch:codigo} (algoritmos), \cref{ch:figuras} (plots y captions),
\cref{ch:diagramas} (TikZ), \cref{ch:tablas} (tablas).

%───────────────────────────────────────────────────────────────────────────────
\section{Protocolo de extensión}
\label{ap:sec:extension}
%───────────────────────────────────────────────────────────────────────────────

\begin{definicion}{Protocolo para añadir un dominio nuevo}
\label{ap:def:extension}
Cuando el template necesita soporte para un paquete o tipo de elemento
nuevo, el protocolo es exactamente este — no más, no menos:
\begin{enumerate}[label=\textbf{\arabic*.}, leftmargin=2em]
  \item Añadir una sección comentada en \texttt{themes/catppuccin-palette.tex}
    con los \texttt{\textbackslash colorlet} del nuevo dominio.
  \item Si los entornos necesitan fondos, llamar
    \texttt{\textbackslash createback\{alias\}} para cada uno.
  \item Crear \texttt{config/\{dominio\}.tex} consumiendo únicamente
    los aliases recién definidos. Nunca \texttt{Ctp*}.
  \item Registrar \texttt{\textbackslash input\{config/\{dominio\}.tex\}}
    en \texttt{main.tex} en el paso 8 (configs de dominio).
\end{enumerate}
\end{definicion}

El ejemplo canónico para un dominio \texttt{circuitikz} hipotético:

\begin{lstlisting}[style=latex, frame=none, numbers=none]
% En themes/catppuccin-palette.tex:
% - Circuitos (circuitikz) 
% Roles: pasivos, activos, fuentes, cables, labels
\colorlet{circ-passive}{CtpBlue}    % resistencias, capacitores
\colorlet{circ-active}{CtpGreen}    % transistores, diodos
\colorlet{circ-source}{CtpYellow}   % fuentes de voltaje/corriente
\colorlet{circ-wire}{CtpOverlay1}   % cables y conexiones
\colorlet{circ-label}{CtpSubtext0}  % labels de componentes

% En config/circuitikz.tex:
% Nunca Ctp* aqui
\ctikzset{
  resistors/scale=0.8,
  color=circ-wire,
}
\tikzset{
  component/.style={color=circ-passive, fill=circ-passive!10},
  source/.style   ={color=circ-source},
}

% En main.tex (paso 8):
\input{config/circuitikz.tex}
\end{lstlisting}

%───────────────────────────────────────────────────────────────────────────────
\section{El mecanismo \textbackslash createback}
\label{ap:sec:createback}
%───────────────────────────────────────────────────────────────────────────────

\begin{definicion}{\textbackslash createback}
\label{ap:def:createback}
El macro \texttt{\textbackslash createback\{alias\}} genera
\texttt{\{alias\}-back} como variante de fondo clara para
\texttt{tcolorbox}. La luminosidad se adapta automáticamente al flavor:
\begin{itemize}[leftmargin=2em]
  \item Tema claro (Latte):
    \texttt{\{alias\}-back} = \texttt{white!97!\{alias\}}
  \item Tema oscuro (Frappé, Macchiato, Mocha):
    \texttt{\{alias\}-back} = \texttt{CtpBase!95!\{alias\}}
\end{itemize}
\end{definicion}

Esto garantiza que todos los entornos teoremáticos tengan fondos
consistentes sin importar el flavor. El toggle \texttt{\textbackslash ifdarktheme}
se establece manualmente en \texttt{catppuccin-palette.tex} junto a la
línea de \texttt{\textbackslash usepackage[Flavor]}.

%───────────────────────────────────────────────────────────────────────────────
\section{Fundamento en el Style Guide oficial de Catppuccin}
\label{ap:sec:fundamento}
%───────────────────────────────────────────────────────────────────────────────

Las asignaciones de este template no son arbitrarias. Cada color de dominio
tiene una correspondencia directa con las categorías del Style Guide oficial
de Catppuccin. La siguiente tabla traza esa correspondencia para los dominios
de mayor visibilidad:

\begin{acadtable}[H]{ap:tab:sg-mapping}{Correspondencia con el Style Guide oficial}
  \begin{tabularx}{\textwidth}{l l X}
    \toprule
    \headerrow
    \thead{Alias del template} & \thead{Categoría Style Guide} & \thead{Descripción SG} \\
    \midrule
    \texttt{link-internal}   & Links, URLs      & ``Blue for links'' \\
    \texttt{link-visited}    & Followed Links   & ``Lavender for followed links'' \\
    \texttt{link-hover}      & On Hover Links   & ``Sky for hover'' \\
    \texttt{code-keyword}    & Keyword          & ``Mauve for keywords'' \\
    \texttt{code-string}     & Strings          & ``Green for strings'' \\
    \texttt{code-comment}    & Comments         & ``Overlay 2 for comments'' \\
    \texttt{code-operator}   & Operators        & ``Sky for operators'' \\
    \texttt{code-type}       & Classes, Types   & ``Yellow for types'' \\
    \texttt{code-fn}         & Methods          & ``Blue for methods'' \\
    \texttt{fb-success}      & Success          & ``Green for success'' \\
    \texttt{fb-warning}      & Warnings         & ``Yellow for warnings'' \\
    \texttt{fb-error}        & Errors           & ``Red for errors'' \\
    \texttt{fb-info}         & Information      & ``Teal for information'' \\
    \texttt{env-thm}         & Highest priority & ``Mauve for most prominent'' \\
    \texttt{rainbow-1..6}    & Rainbow Highlights & Secuencia de colores ordenada \\
    \bottomrule
  \end{tabularx}
\end{acadtable}

Las extensiones propias del template (dominios \texttt{eq-}, \texttt{doc-},
\texttt{tikz-}, \texttt{table-}, \texttt{caption-}, \texttt{plot-}) no
tienen correspondencia directa en el Style Guide porque ese no cubre
documentos matemáticos. Las asignaciones se derivaron de los principios
del Style Guide: colores fríos para elementos epistémicos neutros (Blue
para operadores, Sapphire para dominios), colores cálidos para elementos
con carga semántica fuerte (Peach para incógnitas, Red para vectores).

\begin{notasbib}
El Style Guide oficial de Catppuccin está disponible en
\texttt{https://github.com/catppuccin/catppuccin/blob/main/docs/style-guide.md}.
La paleta completa con valores hexadecimales para cada flavor se documenta
en \texttt{https://catppuccin.com/palette}. El paquete
\texttt{catppuccinpalette} para \LaTeX\ se mantiene en CTAN y en
\texttt{https://github.com/catppuccin/latex}.
\end{notasbib}

  %═══════════════════════════════════════════════════════════════════════════════
\chapter{Galería de código fuente}
\label{ap:codigo-galeria}
%═══════════════════════════════════════════════════════════════════════════════

Este apéndice presenta ejemplos de código en los diez lenguajes configurados
en la plantilla. Cada estilo mapea las categorías sintácticas del lenguaje a
la paleta Catppuccin según la tabla del \cref{sec:codigo-breve}. Los estilos
se invocan con \verb|style=python|, \verb|style=rust|, etc.

\dificultad{1}{1}

%───────────────────────────────────────────────────────────────────────────────
\section{Python}
%───────────────────────────────────────────────────────────────────────────────

\begin{lstlisting}[style=python, caption={Descenso de gradiente con momentum de Nesterov}]
import numpy as np
from typing import Callable, Tuple

class SGDMomentum:
    """SGD optimizer with Nesterov momentum."""

    def __init__(self, lr: float = 1e-3, momentum: float = 0.9):
        self.lr = lr
        self.momentum = momentum
        self.velocity = None

    def step(self, params: np.ndarray, grad: np.ndarray) -> np.ndarray:
        if self.velocity is None:
            self.velocity = np.zeros_like(params)
        # Nesterov momentum update
        self.velocity = self.momentum * self.velocity - self.lr * grad
        return params + self.momentum * self.velocity - self.lr * grad

def train(model: Callable, X: np.ndarray, y: np.ndarray,
          epochs: int = 100) -> list:
    optimizer = SGDMomentum(lr=0.01)
    losses = []
    for epoch in range(epochs):
        y_hat = model(X)
        loss = np.mean((y_hat - y) ** 2)
        losses.append(loss)
        if loss < 1e-6:
            print(f"Converged at epoch {epoch}")
            break
    return losses
\end{lstlisting}

%───────────────────────────────────────────────────────────────────────────────
\section{C}
%───────────────────────────────────────────────────────────────────────────────

\begin{lstlisting}[style=c, caption={Producto matriz-vector}]
#include <stdio.h>
#include <stdlib.h>
#include <string.h>
#include <math.h>

void matvec(const double *A, const double *x, double *y,
            size_t m, size_t n) {
    memset(y, 0, m * sizeof(double));
    for (size_t i = 0; i < m; i++)
        for (size_t j = 0; j < n; j++)
            y[i] += A[i * n + j] * x[j];
}

double norm2(const double *v, size_t n) {
    double s = 0.0;
    for (size_t i = 0; i < n; i++)
        s += v[i] * v[i];
    return sqrt(s);
}

int main(void) {
    double A[] = {1.0, 2.0, 3.0, 4.0, 5.0, 6.0};
    double x[] = {1.0, 1.0, 1.0};
    double y[2];
    matvec(A, x, y, 2, 3);
    printf("||y|| = %.6f\n", norm2(y, 2));
    return EXIT_SUCCESS;
}
\end{lstlisting}

%───────────────────────────────────────────────────────────────────────────────
\section{C++}
%───────────────────────────────────────────────────────────────────────────────

\begin{lstlisting}[style=cpp, caption={Integrador de Monte Carlo con templates}]
#include <iostream>
#include <vector>
#include <random>
#include <functional>

template<typename T>
class MonteCarloIntegrator {
private:
    std::mt19937 rng_;
    std::uniform_real_distribution<T> dist_;
public:
    explicit MonteCarloIntegrator(unsigned seed = 42)
        : rng_(seed), dist_(0.0, 1.0) {}

    T integrate(std::function<T(T)> f, T a, T b, size_t N) {
        T sum = 0.0;
        for (size_t i = 0; i < N; ++i) {
            T x = a + (b - a) * dist_(rng_);
            sum += f(x);
        }
        return (b - a) * sum / static_cast<T>(N);
    }
};

int main() {
    MonteCarloIntegrator<double> mc;
    auto result = mc.integrate(
        [](double x) { return std::sin(x); },
        0.0, M_PI, 1'000'000
    );
    std::cout << "Result: " << result << std::endl;
    return 0;
}
\end{lstlisting}

%───────────────────────────────────────────────────────────────────────────────
\section{Rust}
%───────────────────────────────────────────────────────────────────────────────

\begin{lstlisting}[style=rust, caption={Solucionador RK4 con traits}]
trait OdeSolver {
    fn step(&self, f: &dyn Fn(f64, &[f64]) -> Vec<f64>,
            t: f64, y: &[f64], dt: f64) -> Vec<f64>;
}

struct RungeKutta4;

impl OdeSolver for RungeKutta4 {
    fn step(&self, f: &dyn Fn(f64, &[f64]) -> Vec<f64>,
            t: f64, y: &[f64], dt: f64) -> Vec<f64> {
        let n = y.len();
        let k1 = f(t, y);
        let y2: Vec<f64> = (0..n).map(|i| y[i] + 0.5*dt*k1[i]).collect();
        let k2 = f(t + 0.5*dt, &y2);
        let y3: Vec<f64> = (0..n).map(|i| y[i] + 0.5*dt*k2[i]).collect();
        let k3 = f(t + 0.5*dt, &y3);
        let y4: Vec<f64> = (0..n).map(|i| y[i] + dt*k3[i]).collect();
        let k4 = f(t + dt, &y4);
        (0..n).map(|i| {
            y[i] + dt/6.0 * (k1[i] + 2.0*k2[i] + 2.0*k3[i] + k4[i])
        }).collect()
    }
}

fn main() {
    let rk4 = RungeKutta4;
    let harmonic = |_t: f64, y: &[f64]| -> Vec<f64> {
        vec![y[1], -y[0]]
    };
    let mut y = vec![1.0, 0.0];
    for i in 0..1000 {
        y = rk4.step(&harmonic, i as f64 * 0.01, &y, 0.01);
    }
    println!("y(10) = {:.6}, y'(10) = {:.6}", y[0], y[1]);
}
\end{lstlisting}

%───────────────────────────────────────────────────────────────────────────────
\section{Go}
%───────────────────────────────────────────────────────────────────────────────

\begin{lstlisting}[style=go, caption={Problema de Basilea con goroutines}]
package main

import (
    "fmt"
    "math"
    "sync"
)

type Result struct {
    Index int
    Value float64
}

func partialSum(start, end int, ch chan<- Result, wg *sync.WaitGroup) {
    defer wg.Done()
    sum := 0.0
    for k := start; k < end; k++ {
        sum += 1.0 / math.Pow(float64(k), 2)
    }
    ch <- Result{Index: start, Value: sum}
}

func main() {
    const N = 1_000_000
    const workers = 8
    chunk := N / workers

    ch := make(chan Result, workers)
    var wg sync.WaitGroup

    for i := 0; i < workers; i++ {
        wg.Add(1)
        s := i*chunk + 1
        e := s + chunk
        if i == workers-1 { e = N + 1 }
        go partialSum(s, e, ch, &wg)
    }
    go func() { wg.Wait(); close(ch) }()

    total := 0.0
    for r := range ch { total += r.Value }
    fmt.Printf("pi ~ %.10f\n", math.Sqrt(6.0*total))
}
\end{lstlisting}

%───────────────────────────────────────────────────────────────────────────────
\section{Julia}
%───────────────────────────────────────────────────────────────────────────────

\begin{lstlisting}[style=julia, caption={Euler--Maruyama con dispatch múltiple}]
using Random, LinearAlgebra, Statistics

abstract type AbstractSDE end

struct OrnsteinUhlenbeck <: AbstractSDE
    theta::Float64
    mu::Float64
    sigma::Float64
end

drift(sde::OrnsteinUhlenbeck, x, t) = sde.theta * (sde.mu - x)
diffusion(sde::OrnsteinUhlenbeck, x, t) = sde.sigma

function euler_maruyama(sde::AbstractSDE, x0::Float64,
                        dt::Float64, T::Float64; seed::Int=42)
    rng = MersenneTwister(seed)
    N = Int(T / dt)
    x = zeros(N + 1)
    x[1] = x0
    @inbounds for i in 1:N
        dW = sqrt(dt) * randn(rng)
        x[i+1] = x[i] + drift(sde, x[i], (i-1)*dt)*dt +
                         diffusion(sde, x[i], (i-1)*dt)*dW
    end
    return range(0.0, T; length=N+1), x
end

ou = OrnsteinUhlenbeck(2.0, 0.0, 1.0)
t, x = euler_maruyama(ou, 1.0, 1e-3, 10.0)
@info "Stationary variance" exact=ou.sigma^2/(2*ou.theta) numerical=var(x[end-1000:end])
\end{lstlisting}

%───────────────────────────────────────────────────────────────────────────────
\section{MATLAB}
%───────────────────────────────────────────────────────────────────────────────

\begin{lstlisting}[style=matlab, caption={Elementos finitos 1D}]
function [u, x] = fem1d(N, f_rhs, a, b)
% FEM1D  Solves -u''(x) = f(x) on [a,b], u(a)=u(b)=0
%        Linear finite elements with N elements.

h = (b - a) / N;
x = linspace(a, b, N+1)';
K = sparse(N+1, N+1);
F = zeros(N+1, 1);

for e = 1:N
    Ke = (1/h) * [1 -1; -1 1];
    xm = (x(e) + x(e+1)) / 2;
    Fe = h/2 * [f_rhs(xm); f_rhs(xm)];
    idx = [e, e+1];
    K(idx, idx) = K(idx, idx) + Ke;
    F(idx) = F(idx) + Fe;
end

interior = 2:N;
u = zeros(N+1, 1);
u(interior) = K(interior, interior) \ F(interior);

u_exact = @(x) sin(pi*x) / pi^2;
err = sqrt(h * sum((u - u_exact(x)).^2));
fprintf('L2 error: %.2e (h = %.2e)\n', err, h);
end
\end{lstlisting}

%───────────────────────────────────────────────────────────────────────────────
\section{Fortran}
%───────────────────────────────────────────────────────────────────────────────

\begin{lstlisting}[style=fortran, caption={Factorización LU con pivoteo parcial}]
module linear_algebra
  implicit none
  integer, parameter :: dp = selected_real_kind(15, 307)
contains
  subroutine lu_factor(A, n, ipiv, info)
    integer, intent(in)     :: n
    real(dp), intent(inout) :: A(n, n)
    integer, intent(out)    :: ipiv(n), info
    integer  :: i, j, k, imax
    real(dp) :: maxval, temp

    info = 0
    do k = 1, n
      maxval = abs(A(k, k));  imax = k
      do i = k + 1, n
        if (abs(A(i, k)) > maxval) then
          maxval = abs(A(i, k));  imax = i
        end if
      end do
      ipiv(k) = imax
      if (maxval < epsilon(1.0_dp)) then
        info = k;  return
      end if
      if (imax /= k) then
        do j = 1, n
          temp = A(k,j); A(k,j) = A(imax,j); A(imax,j) = temp
        end do
      end if
      do i = k + 1, n
        A(i,k) = A(i,k) / A(k,k)
        do j = k + 1, n
          A(i,j) = A(i,j) - A(i,k) * A(k,j)
        end do
      end do
    end do
  end subroutine lu_factor
end module linear_algebra
\end{lstlisting}

%───────────────────────────────────────────────────────────────────────────────
\section{Assembler (x86-64)}
%───────────────────────────────────────────────────────────────────────────────

\begin{lstlisting}[style=asm, caption={Producto punto con SSE}]
; dot_product_sse(float *a, float *b, int n) -> float
; System V AMD64: rdi=a, rsi=b, edx=n; returns in xmm0
section .text
global dot_product_sse

dot_product_sse:
    xorps   xmm0, xmm0
    test    edx, edx
    jle     .done
    mov     ecx, edx
    shr     ecx, 2
    jz      .remainder
.loop4:
    movups  xmm1, [rdi]
    movups  xmm2, [rsi]
    mulps   xmm1, xmm2
    addps   xmm0, xmm1
    add     rdi, 16
    add     rsi, 16
    dec     ecx
    jnz     .loop4
    ; horizontal reduction
    movaps  xmm1, xmm0
    shufps  xmm1, xmm1, 0x4E
    addps   xmm0, xmm1
    movaps  xmm1, xmm0
    shufps  xmm1, xmm1, 0xB1
    addss   xmm0, xmm1
.remainder:
    and     edx, 3
    jz      .done
.loop1:
    movss   xmm1, [rdi]
    mulss   xmm1, [rsi]
    addss   xmm0, xmm1
    add     rdi, 4
    add     rsi, 4
    dec     edx
    jnz     .loop1
.done:
    ret
\end{lstlisting}

%───────────────────────────────────────────────────────────────────────────────
\section{Mathematica}
%───────────────────────────────────────────────────────────────────────────────

\begin{lstlisting}[style=mathematica, caption={EDP de calor: simbólica + numérica}]
(* 1D heat equation: du/dt = k * d^2u/dx^2 *)

k = 0.01;  L = Pi;

uExact[x_, t_, nTerms_] := Module[{n},
    Sum[(2/(n*Pi)) * Sin[n*x] * Exp[-k*n^2*t],
        {n, 1, nTerms, 2}]
]

sol = NDSolve[{
    D[u[x,t], t] == k * D[u[x,t], {x,2}],
    u[0,t] == 0,  u[L,t] == 0,
    u[x,0] == If[x < L/2, 2*x/L, 2*(1 - x/L)]
}, u, {x, 0, L}, {t, 0, 10}];

error = NIntegrate[
    (u[x,5] /. sol[[1]] - uExact[x, 5, 50])^2,
    {x, 0, L}] // Sqrt;

Print["L2 error at t=5: ", ScientificForm[error, 3]]
\end{lstlisting}

%───────────────────────────────────────────────────────────────────────────────
\section*{Referencia rápida de estilos}
\addcontentsline{toc}{section}{Referencia rápida de estilos}
%───────────────────────────────────────────────────────────────────────────────

\begin{acadtable}{tab:estilos-listings}{Estilos de listings disponibles}
  \begin{tabularx}{\textwidth}{l l Y}
    \toprule
    \headerrow
    \thead{Estilo} & \thead{Lenguaje} & \thead{Uso} \\
    \midrule
    \texttt{python}      & Python         & \verb|style=python| \\
    \rowcolor{table-row-alt}
    \texttt{c}           & C              & \verb|style=c| \\
    \texttt{cpp}         & C++            & \verb|style=cpp| \\
    \rowcolor{table-row-alt}
    \texttt{rust}        & Rust           & \verb|style=rust| \\
    \texttt{go}          & Go             & \verb|style=go| \\
    \rowcolor{table-row-alt}
    \texttt{julia}       & Julia          & \verb|style=julia| \\
    \texttt{matlab}      & MATLAB/Octave  & \verb|style=matlab| \\
    \rowcolor{table-row-alt}
    \texttt{fortran}     & Fortran (2008) & \verb|style=fortran| \\
    \texttt{asm}         & x86-64 (Intel) & \verb|style=asm| \\
    \rowcolor{table-row-alt}
    \texttt{mathematica} & Wolfram        & \verb|style=mathematica| \\
    \bottomrule
  \end{tabularx}
\end{acadtable}

\begin{notasbib}
La configuración de \texttt{listings} sigue la guía de estilo oficial de
Catppuccin (\url{https://catppuccin.com/style-guide}). Los lenguajes Rust
y Go se definen desde cero con \verb|\lstdefinelanguage| porque el paquete
\texttt{listings} no los incluye nativamente.
Para resaltado más avanzado (funciones definidas por el usuario, type
inference), considerar \texttt{minted} con Pygments como alternativa.
\end{notasbib}

\begin{notasbib}
La configuración de \texttt{listings} sigue la guía de estilo oficial de
Catppuccin. Los lenguajes Rust y Go se definen desde cero con
\texttt{\textbackslash lstdefinelanguage} porque \texttt{listings}
no los incluye nativamente.
\end{notasbib}

\end{appendices}

% En appendixA.tex:
\chapter{Style Guide de Catppuccin Académico}
\label{ap:styleguide}
\end{lstlisting}

Las referencias a apéndices usan el mismo sistema que cualquier otro
elemento: \texttt{\textbackslash cref\{ap:styleguide\}} produce
\textit{ap.~A} automáticamente, con el nombre registrado en
\texttt{config/packages.tex}.

\begin{observacion}{Numeración de figuras y tablas en apéndices}
Por defecto, las figuras en el apéndice A se numeran A.1, A.2, etc.
Este comportamiento es automático con el paquete \texttt{appendix}.
No se necesita configuración adicional.
\end{observacion}

%───────────────────────────────────────────────────────────────────────────────
\section{Encabezados y pies de página}
\label{sec:headers}
%───────────────────────────────────────────────────────────────────────────────

El template usa \texttt{fancyhdr}~\cite{fancyhdr2022} configurado en
\texttt{config/geometry.tex}. El estilo por defecto muestra el nombre de
la sección actual a la izquierda y el número de página a la derecha.
Las páginas de inicio de capítulo usan el estilo \texttt{plain}
(solo número de página centrado, sin línea de encabezado).

Para modificar el contenido de los encabezados, editar directamente
\texttt{config/geometry.tex}:

\begin{lstlisting}[style=latex, frame=none, numbers=none]
\fancyhead[L]{\nouppercase{\leftmark}}   % seccion actual, izquierda
\fancyhead[R]{\thepage}                  % numero de pagina, derecha
\fancyhead[C]{}                          % centro vacio
\fancyfoot[C]{}                          % pie vacio
\end{lstlisting}

En documentos de doble cara (\texttt{twoside}), las posiciones L y R se
alternan entre páginas pares e impares. Para ese caso, usar las posiciones
\texttt{LE}, \texttt{RO} (par izquierda, impar derecha) en lugar de
\texttt{L} y \texttt{R}.

%───────────────────────────────────────────────────────────────────────────────
\printbibliography[heading=subbibliography]

%───────────────────────────────────────────────────────────────────────────────
\chapter{Entornos teoremáticos}
\label{ch:entornos}
%───────────────────────────────────────────────────────────────────────────────

Los entornos teoremáticos son la interfaz más visible del template. Cada
uno tiene un rol epistémico preciso: no son simplemente cajas de color
distintos, sino que comunican al lector la naturaleza lógica del contenido.
Una \textbf{Definición} establece vocabulario. Un \textbf{Teorema} es una
verdad demostrada. Un \textbf{Ejemplo} instancia lo abstracto. Confundirlos
degrada el argumento matemático tanto como una demostración incorrecta.

Este capítulo documenta los ocho entornos disponibles, su sintaxis, su rol
semántico, y los casos donde cada uno es el correcto. Al final se cubren los
entornos especiales y el manejo de figuras y tablas dentro de cajas.

%───────────────────────────────────────────────────────────────────────────────
\section{Arquitectura de los entornos}
\label{sec:entornos-arquitectura}
%───────────────────────────────────────────────────────────────────────────────

\begin{definicion}{Generador de entornos}
\label{def:generador}
Todos los entornos teoremáticos se crean con la macro
\texttt{\textbackslash crearEntorno} definida en
\texttt{environments/generator.tex}~\cite{tcolorbox2024b}:
\begin{lstlisting}[style=latex, frame=none, numbers=none]
\crearEntorno{nombre}{color}{Titulo}
{conContador}{opcionesAdicionales}
\end{lstlisting}
El parámetro \texttt{color} es un alias semántico del dominio \texttt{env-*}
definido en la paleta. El fondo de la caja (\texttt{color-back}) se genera
automáticamente mediante \texttt{\textbackslash createback} al definir la paleta.
Ningún archivo de config toca colores \texttt{Ctp*} directamente.
\end{definicion}

La sintaxis de uso es uniforme para todos los entornos: un argumento
obligatorio con el título de la instancia.

\begin{lstlisting}[style=latex, frame=none, numbers=none]
\begin{nombre}{Título de esta instancia}
\label{tipo:etiqueta}   % opcional pero recomendado
Contenido del entorno...
\end{nombre}
\end{lstlisting}

Los entornos con contador incluyen automáticamente el número en el título
(\textit{Definición~1.1 --- Título}). Los entornos sin contador muestran
solo el título.

%───────────────────────────────────────────────────────────────────────────────
\section{Nivel 1: Entornos fundacionales}
\label{sec:entornos-nivel1}
%───────────────────────────────────────────────────────────────────────────────

Los entornos fundacionales establecen la estructura lógica del argumento.
Tienen contador para que puedan referenciarse con precisión.

\subsection{Definición}
\label{subsec:entorno-def}

Color: \texttt{env-def} $\to$ \textcolor{env-def}{\textbf{Blue}}.
Establece vocabulario formal. El contenido de una definición no se discute
ni se demuestra --- simplemente fija el significado de un término dentro
del sistema. Si el lector no conoce el término, la definición es su
contrato de lectura.

\begin{definicion}{Espacio métrico}
\label{def:espacio-metrico}
Un \textbf{espacio métrico} es un par $(\eqdom{X}, \eqfn{d})$ donde $\eqdom{X}$ es un conjunto
y $\eqfn{d} : \eqdom{X} \eqop{\times} \eqdom{X} \eqop{\to} \mathbb{R}_{\geq 0}$ es una \textit{métrica} que
satisface para todo $\eqk{x}, \eqk{y}, \eqk{z} \in \eqdom{X}$:
\begin{enumerate}[label=(\roman*), leftmargin=2em]
  \item $\eqfn{d}(\eqk{x},\eqk{y}) = 0 \iff \eqk{x} = \eqk{y}$ \hfill (identidad)
  \item $\eqfn{d}(\eqk{x},\eqk{y}) = \eqfn{d}(\eqk{y},\eqk{x})$ \hfill (simetría)
  \item $\eqfn{d}(\eqk{x},\eqk{z}) \leq \eqfn{d}(\eqk{x},\eqk{y}) + \eqfn{d}(\eqk{y},\eqk{z})$ \hfill (desigualdad triangular)
\end{enumerate}
\end{definicion}

\subsection{Teorema}
\label{subsec:entorno-thm}

Color: \texttt{env-thm} $\to$ \textcolor{env-thm}{\textbf{Mauve}}.
El color más prominente del sistema, reservado para el resultado más
importante de su categoría. Un teorema es una afirmación no trivial cuya
verdad se ha demostrado formalmente. La prominencia visual señala al lector
que aquí hay un resultado que merece atención sostenida.

\begin{teorema}{Completitud de los espacios $L^p$}
\label{thm:lp-completo}
Sea $(\eqdom{X}, \eqdom{\mathcal{M}}, \eqk{\mu})$ un espacio de medida y $1 \leq \eqk{p} < \infty$.
El espacio $\eqdom{L^p}(\eqdom{X}, \eqk{\mu})$ dotado de la norma
\[
  \|\equ{f}\|_{\eqk{p}} = \left(\eqop{\int}_{\eqdom{X}} |\equ{f}|^{\eqk{p}} \, \eqdm{d\mu}\right)^{1/\eqk{p}}
\]
es un espacio de Banach. En particular, $\eqdom{L^2}(\eqdom{X},\eqk{\mu})$ es un espacio de
Hilbert con producto interior $\langle \equ{f}, \equ{g} \rangle = \eqop{\int}_{\eqdom{X}} \equ{f}\,\overline{\equ{g}}\,\eqdm{d\mu}$.
\end{teorema}

\begin{observacion}{Teorema vs.\ Proposición}
Un \textbf{teorema} es el resultado principal de una línea de argumento ---
el punto al que todo lo demás conduce. Una \textbf{proposición} es un
resultado intermedio correcto y no trivial, pero subordinado. Si tienes
dudas, pregunta: ¿citarías este resultado de forma independiente, sin
contexto? Si sí, es un teorema. Si solo tiene sentido dentro del argumento
local, es una proposición.
\end{observacion}

\subsection{Proposición}
\label{subsec:entorno-prop}

Color: \texttt{env-prop} $\to$ \textcolor{env-prop}{\textbf{Green}}.
Resultado verdadero y demostrable, pero de alcance más local que un teorema.
El verde evoca una afirmación positiva --- algo que se sostiene.

\begin{proposicion}{Continuidad de la norma}
\label{prop:norma-continua}
Sea $(\eqdom{X}, \|\cdot\|)$ un espacio normado. La función $\|\cdot\| : \eqdom{X} \to \mathbb{R}$
es continua respecto a la topología inducida por la norma. Más aún, es
Lipschitz con constante $1$: para todo $\eqk{x}, \eqk{y} \in \eqdom{X}$,
\[
  \big|\|\eqk{x}\| - \|\eqk{y}\|\big| \leq \|\eqk{x} - \eqk{y}\|.
\]
\end{proposicion}

%───────────────────────────────────────────────────────────────────────────────
\section{Nivel 2: Entornos de apoyo}
\label{sec:entornos-nivel2}
%───────────────────────────────────────────────────────────────────────────────

Los entornos de apoyo complementan el argumento sin avanzarlo directamente.
No tienen la misma jerarquía lógica que los fundacionales, pero son
igualmente indispensables para la comprensión.

\subsection{Ejemplo}
\label{subsec:entorno-ej}

Color: \texttt{env-ex} $\to$ \textcolor{env-ex}{\textbf{Peach}}.
Instancia un concepto abstracto en un caso concreto. El tono cálido invita
a la lectura --- los ejemplos deben sentirse accesibles. Un buen ejemplo
no solo ilustra: muestra por qué la hipótesis general es necesaria (el
ejemplo borderline) o por qué el resultado no es trivial.

\begin{ejemplo}{$\ell^2$ como espacio de Hilbert}
\label{ex:l2-hilbert}
El espacio de sucesiones de cuadrado sumable
\[
  \eqdom{\ell^2} = \left\{(\equ{x}_n)_{n=1}^\infty \subset \mathbb{R} \;\middle|\;
  \eqop{\sum}_{n=1}^\infty \equ{x}_n^2 < \infty\right\}
\]
es un espacio de Hilbert~\cite{reed1980} con producto interior
$\langle \equ{x}, \equ{y} \rangle = \eqop{\sum}_{n=1}^\infty \equ{x}_n \equ{y}_n$.
La base ortonormal canónica $\{\eqk{e}_n\}_{n=1}^\infty$, donde $\eqk{e}_n$ tiene
un $1$ en la posición $n$ y ceros en el resto, es una base de Schauder
de $\ell^2$ pero \textit{no} es una base algebraica (Hamel): ningún
conjunto finito genera el espacio.
\end{ejemplo}

\subsection{Observación}
\label{subsec:entorno-obs}

Color: \texttt{env-obs} $\to$ \textcolor{env-obs}{\textbf{Overlay0}}.
El entorno más sutil del sistema --- casi gris. Una observación no avanza
el argumento principal ni instancia un concepto: llama la atención sobre
un detalle que el lector podría pasar por alto. Sirve para matices,
advertencias técnicas, notas de implementación, y condiciones que se
asumen implícitamente. El color discreto señala que puede omitirse en
una primera lectura sin perder el hilo.

\begin{observacion}{Convergencia débil vs. fuerte en $L^2$}
En espacios de Hilbert de dimensión infinita, la convergencia débil
no implica convergencia en norma. La sucesión $\eqk{e}_n \rightharpoonup 0$
débilmente en $\eqdom{\ell^2}$ (porque $\langle \eqk{e}_n, \equ{x} \rangle = \equ{x}_n \to 0$
para todo $\equ{x} \in \eqdom{\ell^2}$), pero $\|\eqk{e}_n\| = 1$ para todo $n$. Este fenómeno
no tiene análogo en dimensión finita y es fuente frecuente de errores en
argumentos de compacidad.
\end{observacion}

\subsection{Conexión física}
\label{subsec:entorno-phys}

Color: \texttt{env-phys} $\to$ \textcolor{env-phys}{\textbf{Sapphire}}.
Relaciona un resultado matemático abstracto con su interpretación o
motivación física. En documentos de física computacional o matemática
aplicada, este entorno ancla la abstracción en la intuición del dominio.
No es un ejemplo --- es un puente entre dos lenguajes.

\begin{fisicaconexion}{Completitud de $L^2$ y mecánica cuántica}
El \cref{thm:lp-completo} para $p=2$ es el fundamento del formalismo
de Hilbert en mecánica cuántica. El espacio de estados de una partícula
cuántica es $\eqdom{L^2}(\eqdom{\mathbb{R}^3})$: las funciones de onda $\equ{\psi}(\eqvec{\mathbf{r}})$
son elementos de $\eqdom{L^2}$, y la norma $\|\equ{\psi}\|_2^2 = \eqop{\int} |\equ{\psi}|^2 \, \eqdm{d^3r}$
es la probabilidad total (normalizada a 1). La completitud garantiza que
los límites de sucesiones de Cauchy de estados físicos son también estados
físicos --- sin ella, la evolución temporal $\equ{\psi}(t) = e^{-\eqcon{i}\eqop{\hat{H}}t/\eqcon{\hbar}}\eqk{\psi_0}$
podría salir del espacio.
\end{fisicaconexion}

%───────────────────────────────────────────────────────────────────────────────
\section{Nivel 3: Entornos especiales}
\label{sec:entornos-nivel3}
%───────────────────────────────────────────────────────────────────────────────

Los entornos especiales tienen propósitos concretos que los diferencian
de los niveles anteriores. Su uso debe ser deliberado.

\subsection{Estrella Polar}
\label{subsec:entorno-estrella}

Color: \texttt{env-star} $\to$ \textcolor{env-star}{\textbf{Yellow}}.
El amarillo llama la atención más que cualquier otro color del sistema.
Reservado para el resultado o idea que orienta toda una unidad temática ---
el punto hacia el que converge el argumento del capítulo o sección, como
una estrella polar orienta la navegación. No es simplemente un teorema
importante: es el \textit{por qué} existe el capítulo. Debe haber a lo
sumo uno por capítulo, y muchos capítulos no lo necesitan.

\begin{estrella}{Teorema espectral en espacios de Hilbert}
\label{star:espectral}
Sea $\eqop{A} : \eqdom{\mathcal{H}} \to \eqdom{\mathcal{H}}$ un operador autoadjunto acotado
en un espacio de Hilbert separable. Existe una medida espectral
$\eqfn{E} : \mathcal{B}(\mathbb{R}) \to \mathcal{L}(\eqdom{\mathcal{H}})$ tal que
\[
  \eqop{A} = \eqop{\int}_{\sigma(\eqop{A})} \equ{\lambda} \, \eqdm{dE}(\equ{\lambda}),
\]
donde la integral converge en la topología de operadores. Todo cálculo
funcional $\eqfn{f}(\eqop{A})$ se define de forma consistente como
$\eqfn{f}(\eqop{A}) = \eqop{\int} \eqfn{f}(\equ{\lambda})\,\eqdm{dE}(\equ{\lambda})$.
\end{estrella}

\subsection{Ejercicio Computacional}
\label{subsec:entorno-comp}

Color: \texttt{env-comp} $\to$ \textcolor{env-comp}{\textbf{Pink}}.
Para tareas que requieren implementación, simulación o verificación
numérica. El tono activo señala que el lector debe hacer algo, no solo
leer. A diferencia de los ejercicios matemáticos ordinarios (que usan
\texttt{\textbackslash ejercicio}), este entorno es para trabajo que
involucra código, datos o computación explícita.

\begin{ejerciciocomp}{Verificación de la desigualdad triangular}
\label{comp:triangular}
Implementa en Python o Julia una función que, dado un conjunto finito de
puntos en $\mathbb{R}^n$ y una función $\eqfn{d} : \eqdom{\mathbb{R}^n} \eqop{\times} \eqdom{\mathbb{R}^n}
\eqop{\to} \mathbb{R}$, verifique si $\eqfn{d}$ es una métrica en el sentido de la
\cref{def:espacio-metrico}. Prueba con:
\begin{enumerate}[label=(\alph*), leftmargin=2em]
  \item $\eqfn{d_2}(\eqk{x},\eqk{y}) = \|\eqk{x}-\eqk{y}\|_2$ (distancia euclídea, debe pasar)
  \item $\eqfn{d_\infty}(\eqk{x},\eqk{y}) = \|\eqk{x}-\eqk{y}\|_\infty$ (norma infinito, debe pasar)
  \item $\eqfn{d}(\eqk{x},\eqk{y}) = \|\eqk{x}-\eqk{y}\|_2^2$ (cuadrado de la euclídea, debe fallar)
\end{enumerate}
¿En qué caso falla el tercer candidato y con qué contraejemplo concreto?
\end{ejerciciocomp}

%───────────────────────────────────────────────────────────────────────────────
\section{El indicador de dificultad}
\label{sec:dificultad}
%───────────────────────────────────────────────────────────────────────────────

El comando \texttt{\textbackslash dificultad\{m\}\{f\}} genera un indicador
visual de dos ejes: dificultad matemática y dificultad física. Cada eje
toma un valor entre 1 y 5. Se coloca al inicio de una sección o subsección
para calibrar las expectativas del lector:

\begin{lstlisting}[style=latex, frame=none, numbers=none]
\dificultad{3}{2}   % matemática: ★★★☆☆  física: ★★☆☆☆
\end{lstlisting}

\dificultad{3}{2}

Los dos ejes son independientes. Un problema de integración numérica puede
tener baja dificultad matemática (cuadratura de Gauss es algorítmica) y
alta dificultad física (elegir el método correcto requiere intuición del
dominio). Separar los dos ejes da información más precisa que un único
número.

%───────────────────────────────────────────────────────────────────────────────
\section{Figuras y tablas dentro de entornos}
\label{sec:floats-en-entornos}
%───────────────────────────────────────────────────────────────────────────────

\begin{definicion}{Restricción de flotantes en tcolorbox}
\label{def:no-floats}
Los entornos \texttt{figure} y \texttt{table} son flotantes: invocan el
mecanismo \texttt{\textbackslash @float} de \LaTeX, que requiere estar
en \textit{outer par mode}. Dentro de una caja \texttt{tcolorbox} ese
modo no existe. El resultado es el error \texttt{Not in outer par mode}.
Los entornos afectados son: \texttt{figure}, \texttt{table},
\texttt{acadtable}, \texttt{lstlisting[float=...]}, y \texttt{longtable}.
\end{definicion}

La solución es usar los entornos no-flotantes \texttt{boxfig} y
\texttt{boxtab} definidos en \texttt{config/tcolorbox.tex}:

\begin{lstlisting}[style=latex, frame=none, numbers=none]
\begin{ejemplo}{Distribución gaussiana}
  \begin{boxfig}{fig:gauss}{Densidad $\mathcal{N}(0,1)$}
    \includegraphics[width=\linewidth]{figs/gauss.pdf}
  \end{boxfig}
  La densidad de probabilidad es $\eqfn{p}(\equ{x}) = \frac{1}{\sqrt{2\eqcon{\pi}}} e^{-\equ{x}^2/2}$.
\end{ejemplo}

\begin{definicion}{Media muestral}
  \begin{boxtab}{tab:media}{Propiedades de la media muestral}
    \begin{tabularx}{\linewidth}{l X}
      \toprule
      \headerrow
      \thead{Propiedad} & \thead{Valor} \\
      \midrule
      Sesgo   & Insesgado: $\eqop{\mathbb{E}}[\equ{\bar{X}}] = \eqk{\mu}$ \\
      Varianza & $\eqop{\mathrm{Var}}(\equ{\bar{X}}) = \eqk{\sigma}^2/\eqk{n}$ \\
      \bottomrule
    \end{tabularx}
  \end{boxtab}
\end{definicion}
\end{lstlisting}

La diferencia con los entornos flotantes es en el ancho: dentro de una
caja siempre se usa \texttt{\textbackslash linewidth} (el ancho disponible
dentro de la caja), nunca \texttt{\textbackslash textwidth}.

\begin{acadtable}[H]{tab:float-comparison}{Entornos flotantes vs.\ no-flotantes}
  \begin{tabularx}{\textwidth}{l l l X}
    \toprule
    \headerrow
    \thead{Situación} & \thead{Figuras} & \thead{Tablas} & \thead{Ancho} \\
    \midrule
    Flujo normal de texto  & \texttt{figure}    & \texttt{acadtable} & \texttt{\textbackslash textwidth} \\
    Dentro de tcolorbox    & \texttt{boxfig}    & \texttt{boxtab}    & \texttt{\textbackslash linewidth} \\
    \bottomrule
  \end{tabularx}
\end{acadtable}

%───────────────────────────────────────────────────────────────────────────────
\section{Notas bibliográficas}
\label{sec:entornos-notasbib}
%───────────────────────────────────────────────────────────────────────────────

Al final de un capítulo, el entorno \texttt{notasbib} sirve para
contextualizar históricamente los resultados presentados: atribuir
resultados a sus autores originales, indicar referencias para profundizar,
o señalar generalizaciones que no cupieron en el cuerpo principal.

\begin{lstlisting}[style=latex, frame=none, numbers=none]
\begin{notasbib}
  El teorema espectral en su forma moderna se debe a von Neumann (1929).
  Para una exposición rigurosa ver Reed y Simon [reed1980],
  Vol.~I, Cap.~VII...
\end{notasbib}
\end{lstlisting}

\begin{notasbib}
El marco general de los espacios $L^p$ fue desarrollado por Riesz
(1910) y Fischer (1907), cuya equivalencia estableció el
\textit{teorema de Riesz-Fischer}. Para una exposición moderna y rigurosa,
Reed y Simon~\cite{reed1980} es la referencia estándar en análisis
funcional orientado a la física matemática. Los espacios métricos en su
generalidad moderna se deben a Fréchet (1906); Munkres~\cite{munkres2000}
ofrece una introducción topológica accesible.
\end{notasbib}

%───────────────────────────────────────────────────────────────────────────────
\printbibliography[heading=subbibliography]

%───────────────────────────────────────────────────────────────────────────────
\chapter{Matemáticas semánticas}
\label{ch:ecuaciones}
%───────────────────────────────────────────────────────────────────────────────

La tipografía matemática estándar codifica una sola dimensión de información:
el tipo matemático. Los vectores van en negrita, los operadores en redonda,
los escalares en itálica. Pero una ecuación contiene al menos tres dimensiones
de información simultáneas que la tipografía convencional no distingue: el
\textit{rol sintáctico} del símbolo (¿es un operador, una función, un dominio?),
su \textit{estatus epistémico} (¿es una incógnita, un dato fijo, una constante
universal?) y su \textit{tipo matemático} (¿es un vector, un tensor, un escalar?).

El sistema de coloreado semántico de este template codifica estas tres
dimensiones mediante color, con un mapa fijo entre rol y tono que es
consistente en todo el documento.

%───────────────────────────────────────────────────────────────────────────────
\section{Regla fundamental: siempre en modo matemático}
\label{sec:regla-matmode}
%───────────────────────────────────────────────────────────────────────────────

\begin{definicion}{Modo matemático obligatorio}
\label{def:matmode}
Todas las macros de coloreado semántico (\texttt{\textbackslash eqop},
\texttt{\textbackslash equ}, \texttt{\textbackslash eqk}, etc.) deben usarse
\textbf{dentro de un entorno matemático}. Fuera de modo matemático, los
símbolos que colorean son glifos de texto, no glifos matemáticos: el color
se aplica sobre el carácter incorrecto o no se aplica en absoluto.
\end{definicion}

La forma correcta e incorrecta es:

\begin{lstlisting}[style=latex, frame=none, numbers=none]
% CORRECTO: dentro de modo matematico
$\eqop{\nabla} \equ{f}$          % inline
\[ \eqop{\nabla}^2 \equ{\phi} = \eqk{\rho} / \eqcon{\varepsilon_0} \]

% INCORRECTO: fuera de modo matematico
\eqop{\nabla} \equ{f}   % el color se aplica sobre texto, no glifo matematico
\end{lstlisting}

Esto aplica en todos los contextos: texto corrido, tablas, títulos de
entornos, labels de ejes en pgfplots, labels de nodos en TikZ. En cada
caso, el wrapper matemático es obligatorio.

%───────────────────────────────────────────────────────────────────────────────
\section{Los tres ejes de coloreado}
\label{sec:tres-ejes}
%───────────────────────────────────────────────────────────────────────────────

\dificultad{2}{3}

\subsection{Eje 1 — Rol sintáctico}
\label{subsec:eje-sintactico}

El rol sintáctico de un símbolo es estable entre contextos: $\nabla$ es
siempre un operador diferencial, independientemente de qué ecuación aparezca
o cuál sea la incógnita.

\begin{acadtable}[H]{tab:eje-sintactico}{Eje 1: Rol sintáctico}
  \begin{tabularx}{\textwidth}{l l l X}
    \toprule
    \headerrow
    \thead{Macro} & \thead{Alias} & \thead{Color} & \thead{Cubre} \\
    \midrule
    \texttt{\textbackslash eqop\{·\}}
      & \texttt{eq-op} & \textcolor{eq-op}{\textbf{Blue}}
      & Operadores: $\eqop{\int}$, $\eqop{\sum}$, $\eqop{\partial}$,
        $\eqop{\nabla}$, $\eqop{\times}$, $\eqop{\delta}(\cdot,\cdot)$ \\
    \texttt{\textbackslash eqfn\{·\}}
      & \texttt{eq-fn} & \textcolor{eq-fn}{\textbf{Green}}
      & Funciones/mapas: $\eqfn{f}(\cdot)$, $\eqfn{\sigma}(\cdot)$,
        $\eqfn{V}(\mathbf{r})$, $\eqfn{\mathcal{L}}$ \\
    \texttt{\textbackslash eqdm\{·\}}
      & \texttt{eq-dm} & \textcolor{eq-dm}{\textbf{Teal}}
      & Diferenciales/medidas: $\eqdm{dx}$, $\eqdm{dt}$,
        $\eqdm{dW}$, $\eqdm{d^3r}$, $\eqdm{d\mu}$ \\
    \texttt{\textbackslash eqdom\{·\}}
      & \texttt{eq-dom} & \textcolor{eq-dom}{\textbf{Sapphire}}
      & Dominios/conjuntos: $\eqdom{\Omega}$, $\eqdom{\partial\Omega}$,
        $\eqdom{\mathbb{R}^n}$, $\eqdom{\Sigma}$, $\eqdom{Q}$ \\
    \bottomrule
  \end{tabularx}
\end{acadtable}

\subsection{Eje 2 — Estatus epistémico}
\label{subsec:eje-epistemico}

El estatus epistémico es la dimensión que la tipografía convencional
\textit{no} codifica. En $f(\mathbf{x};\boldsymbol{\theta})$, $\mathbf{x}$
es dato de entrada y $\boldsymbol{\theta}$ es el parámetro a aprender ---
en inferencia estadística es exactamente al revés. El color hace explícita
esta inversión sin texto adicional.

\begin{acadtable}[H]{tab:eje-epistemico}{Eje 2: Estatus epistémico}
  \begin{tabularx}{\textwidth}{l l l X}
    \toprule
    \headerrow
    \thead{Macro} & \thead{Alias} & \thead{Color} & \thead{Cubre} \\
    \midrule
    \texttt{\textbackslash equ\{·\}}
      & \texttt{eq-unknown} & \textcolor{eq-unknown}{\textbf{Peach}}
      & Incógnita / buscado / aprendido: la cantidad que se desea hallar \\
    \texttt{\textbackslash eqk\{·\}}
      & \texttt{eq-known} & \textcolor{eq-known}{\textbf{Mauve}}
      & Conocido / dato / parámetro fijado: entra al sistema como input \\
    \texttt{\textbackslash eqcon\{·\}}
      & \texttt{eq-const} & \textcolor{eq-const}{\textbf{Lavender}}
      & Constante universal/estructural: $\eqcon{\hbar}$, $\eqcon{k_B}$,
        $\eqcon{\pi}$, $\eqcon{\varepsilon}$ \\
    \bottomrule
  \end{tabularx}
\end{acadtable}

\subsection{Eje 3 — Tipo matemático}
\label{subsec:eje-tipo}

El eje de tipo refuerza lo que la tipografía ya comunica: la negrita indica
vector, la doble barra indica espacio. El color añade un canal redundante
útil en documentos donde la distinción rango-tensor es semánticamente
importante y el lector necesita identificarla de un vistazo.

\begin{acadtable}[H]{tab:eje-tipo}{Eje 3: Tipo matemático}
  \begin{tabularx}{\textwidth}{l l l X}
    \toprule
    \headerrow
    \thead{Macro} & \thead{Alias} & \thead{Color} & \thead{Cubre} \\
    \midrule
    \texttt{\textbackslash eqvec\{·\}}
      & \texttt{eq-vec} & \textcolor{eq-vec}{\textbf{Red}}
      & Vectores/campos: $\eqvec{\mathbf{u}}$, $\eqvec{\mathbf{E}}$,
        $\eqvec{\mathbf{B}}$, $\eqvec{\mathbf{J}}$ \\
    \texttt{\textbackslash eqten\{·\}}
      & \texttt{eq-ten} & \textcolor{eq-ten}{\textbf{Rosewater}}
      & Tensores/matrices: $\eqten{\sigma_{ij}}$, $\eqten{g_{\mu\nu}}$,
        $\eqten{W}$, $\eqten{\Sigma}$ \\
    \bottomrule
  \end{tabularx}
\end{acadtable}

%───────────────────────────────────────────────────────────────────────────────
\section{Ecuaciones completas con los tres ejes}
\label{sec:ecuaciones-completas}
%───────────────────────────────────────────────────────────────────────────────

La riqueza del sistema aparece cuando los tres ejes se aplican
simultáneamente en una ecuación completa. Los ejemplos siguientes
usan física y machine learning para mostrar el sistema bajo carga real.

\subsection{Mecánica cuántica: ecuación de Schrödinger}
\label{subsec:schrodinger}

\begin{fisicaconexion}{Ecuación de Schrödinger~\cite{griffiths2017quantum}}
La ecuación de Schrödinger dependiente del tiempo:
\[
  \eqop{i}\eqcon{\hbar}\,\eqop{\frac{\partial}{\partial t}}\equ{\Psi}
  = \eqop{\hat{H}}\equ{\Psi}
  = \left(\eqop{-}\frac{\eqcon{\hbar}^2}{2\eqk{m}}\eqop{\nabla}^2
    + \eqk{V}(\eqvec{\mathbf{r}},t)\right)\equ{\Psi}
\]
El color comunica: $\eqcon{\hbar}$ y $\eqcon{i}$ son constantes estructurales.
$\eqk{m}$ es la masa, un parámetro fijado del sistema.
$\eqk{V}(\mathbf{r},t)$ es el potencial, conocido.
$\equ{\Psi}(\mathbf{r},t)$ es la función de onda, la incógnita que se busca.
$\eqop{\partial/\partial t}$ y $\eqop{\nabla^2}$ son operadores diferenciales.
\end{fisicaconexion}

\subsection{Electrodinámica: ecuaciones de Maxwell}
\label{subsec:maxwell}

\begin{fisicaconexion}{Ecuaciones de Maxwell~\cite{jackson1999electrodynamics}}
En forma diferencial, con unidades SI:
\begin{align}
  \eqop{\nabla \cdot} \eqvec{\mathbf{E}}
    &= \frac{\eqk{\rho}}{\eqcon{\varepsilon_0}}
  \label{eq:gauss-e} \\
  \eqop{\nabla \cdot} \eqvec{\mathbf{B}}
    &= 0
  \label{eq:gauss-b} \\
  \eqop{\nabla \times} \eqvec{\mathbf{E}}
    &= \eqop{-}\eqop{\frac{\partial}{\partial t}}\eqvec{\mathbf{B}}
  \label{eq:faraday} \\
  \eqop{\nabla \times} \eqvec{\mathbf{B}}
    &= \eqcon{\mu_0}\left(\eqvec{\mathbf{J}}
      + \eqcon{\varepsilon_0}\eqop{\frac{\partial}{\partial t}}
      \eqvec{\mathbf{E}}\right)
  \label{eq:ampere}
\end{align}
Los campos $\eqvec{\mathbf{E}}$ y $\eqvec{\mathbf{B}}$ son incógnitas
vectoriales en rojo. Las fuentes $\eqk{\rho}$ y $\eqvec{\mathbf{J}}$ son
datos del problema en malva. Las constantes $\eqcon{\varepsilon_0}$ y
$\eqcon{\mu_0}$ en lavanda son universales.
\end{fisicaconexion}

\subsection{Machine learning: ELBO variacional}
\label{subsec:elbo}

\begin{ejemplo}{Evidence Lower BOund (ELBO)~\cite{goodfellow2016deep}}
\label{ex:elbo}
En inferencia variacional, el objetivo es maximizar:
\[
  \eqfn{\mathcal{L}}(\eqk{\phi}, \eqk{\theta})
  = \eqop{\mathbb{E}}_{\equ{z} \sim \eqfn{q}_{\eqk{\phi}}}
    \!\left[\log \eqfn{p}_{\eqk{\theta}}(\eqk{x} \mid \equ{z})\right]
  \eqop{-} \eqop{D_{\mathrm{KL}}}\!\left(
    \eqfn{q}_{\eqk{\phi}}(\equ{z} \mid \eqk{x})
    \,\|\, \eqfn{p}(\equ{z})
  \right)
\]
El color hace explícita la estructura: $\eqk{\phi}$ y $\eqk{\theta}$ son los
parámetros del encoder y decoder, datos del sistema de optimización.
$\equ{z}$ es la variable latente, la incógnita que se integra.
$\eqfn{q}_\phi$ y $\eqfn{p}_\theta$ son las distribuciones (funciones del
modelo). $\eqop{\mathbb{E}}$ y $\eqop{D_{\mathrm{KL}}}$ son operadores.
\end{ejemplo}

\subsection{Física computacional: ecuación de Fokker-Planck}
\label{subsec:fokker-planck}

\begin{definicion}{Ecuación de Fokker-Planck}
\label{def:fokker-planck}
La evolución de la densidad de probabilidad $\equ{p}(\eqvec{\mathbf{x}},t)$
de un proceso de difusión gobernado por la SDE
$\eqdm{d}\equ{X} = \eqfn{f}(\equ{X})\eqdm{dt} + \eqk{\sigma}\eqdm{dW}$
está dada por:
\[
  \eqop{\frac{\partial}{\partial t}}\equ{p}
  = \eqop{-}\eqop{\nabla_x \cdot}
    \!\left(\eqfn{f}\,\equ{p}\right)
  + \frac{\eqk{\sigma}^2}{2}\,\eqop{\nabla_x^2}\,\equ{p}
\]
donde $\eqop{\nabla_x \cdot}$ es el operador divergencia y
$\eqop{\nabla_x^2}$ es el Laplaciano, ambos actuando sobre la
variable espacial $\eqvec{\mathbf{x}}$. El coeficiente de difusión
$\eqk{\sigma}^2/2$ es un parámetro conocido.
\end{definicion}

%───────────────────────────────────────────────────────────────────────────────
\section{Macros compuestos}
\label{sec:macros-compuestos}
%───────────────────────────────────────────────────────────────────────────────

\texttt{config/equations.tex} define macros compuestos para patrones
frecuentes en física y machine learning. Usan los ejes semánticos
internamente.

\begin{acadtable}[H]{tab:macros-compuestos}{Macros compuestos de equations.tex}
  \begin{tabularx}{\textwidth}{l X X}
    \toprule
    \headerrow
    \thead{Macro} & \thead{Uso} & \thead{Produce} \\
    \midrule
    \texttt{\textbackslash eqpdv\{f\}\{t\}}
      & Derivada parcial coloreada
      & $\eqpdv{f}{t}$ \\
    \texttt{\textbackslash eqgrad\{u\}}
      & Gradiente coloreado
      & $\eqgrad{u}$ \\
    \texttt{\textbackslash eqdvg\{\textbackslash mathbf\{u\}\}}
      & Divergencia coloreada
      & $\eqdvg{\mathbf{u}}$ \\
    \texttt{\textbackslash eqSDE\{X\}\{f(X)\}\{\textbackslash sigma\}}
      & SDE coloreada completa
      & $\eqSDE{X}{f(X)}{\sigma}$ \\
    \bottomrule
  \end{tabularx}
\end{acadtable}

El macro \texttt{\textbackslash eqpdv} aplica la convención epistémica
natural de las derivadas parciales: el numerador es siempre la incógnita
(lo que cambia, lo que se busca entender) y el denominador es la variable
independiente, dato del problema:

\begin{lstlisting}[style=latex, frame=none, numbers=none]
% Ecuacion de calor: u es la incognita, t y x son variables independientes
\[
  \eqpdv{u}{t} = \eqk{\alpha}\, \eqpdv[2]{u}{x}
\]
\end{lstlisting}

\[
  \eqpdv{u}{t} = \eqk{\alpha}\, \eqop{\frac{\partial^2 \equ{u}}{\partial \eqk{x}^2}}
\]

%───────────────────────────────────────────────────────────────────────────────
\section{Coloreado en contextos distintos}
\label{sec:contextos}
%───────────────────────────────────────────────────────────────────────────────

El sistema funciona en cualquier contexto donde se pueda escribir matemáticas.
La regla del modo matemático (\cref{def:matmode}) aplica en todos.

\subsection{En tablas}
\label{subsec:eq-en-tablas}

El documento de referencia muestra el caso de expresiones regulares con
theming en tablas: $\eqk{a} \eqop{\mid} \eqk{b}$ en una celda de
\texttt{tabularx}. El patrón es simplemente matemáticas normales dentro de la celda:

\begin{acadtable}[H]{tab:regexp-coloreado}{Expresiones regulares con coloreado semántico}
  \begin{tabularx}{\textwidth}{X X}
    \toprule
    \headerrow
    \thead{Expresión regular} & \thead{Lenguaje descrito} \\
    \midrule
    $\eqk{a} \eqop{\mid} \eqk{b}$
      & $\eqdom{\{}\eqk{a}, \eqk{b}\eqdom{\}}$ \\
    $\eqk{a}\eqop{*}$
      & $\eqdom{\{}\eqcon{\varepsilon}, \eqk{a}, \eqk{aa}, \ldots\eqdom{\}}$ \\
    $(\eqk{a} \eqop{\mid} \eqk{b})\eqop{*}\eqk{a}\eqk{b}\eqk{b}$
      & Cadenas sobre $\eqdom{\{a,b\}}$ con sufijo $\eqk{abb}$ \\
    $\eqop{\delta} : \eqdom{Q} \eqop{\times} \eqdom{\Sigma}
      \eqop{\rightarrow} \eqdom{Q}$
      & Función de transición de un DFA \\
    \bottomrule
  \end{tabularx}
\end{acadtable}

\subsection{En algoritmos}
\label{subsec:eq-en-algoritmos}

Dentro de \texttt{algorithm2e}, el theming funciona directamente en modo
matemático. El documento de referencia muestra algoritmos de construcción
de subconjuntos con $\eqk{q'_0} \eqop{\gets} \eqop{E}(\eqdom{\{}\eqk{q_0}\eqdom{\}})$ ---
el estado inicial (conocido) se inicializa con el operador de clausura:

\begin{lstlisting}[style=latex, frame=none, numbers=none]
\begin{algorithm}[H]
  \KwIn{NFA $\eqk{N} = (\eqdom{Q}, \eqdom{\Sigma}, \eqop{\delta},
               \eqk{q_0}, \eqdom{F})$}
  $\eqk{q'_0} \eqop{\gets} \eqop{E}(\eqdom{\{}\eqk{q_0}\eqdom{\}})$\;
  \While{$\equ{T} \eqop{\in} \eqdom{Q'}$ no marcado}{
    $\equ{U} \eqop{\gets} \eqop{E}(\eqop{move}(\equ{T}, \equ{a}))$\;
  }
\end{algorithm}
\end{lstlisting}

\subsection{En TikZ}
\label{subsec:eq-en-tikz}

En nodos y aristas de TikZ, los labels matemáticos van dentro de \texttt{\$...\$}.
El theming es idéntico al de cualquier ecuación:

\begin{lstlisting}[style=latex, frame=none, numbers=none]
\path[->]
  (q0) edge node {$\eqk{a}$} (q1)
  (q0) edge node {$\eqcon{\varepsilon}$} (q2);
\end{lstlisting}

\begin{observacion}{Colores en TikZ: usar aliases, no Ctp*}
El documento de referencia usa \texttt{CtpBlue}, \texttt{CtpBase},
\texttt{CtpText} directamente en los estilos de nodos TikZ.
Esto viola la filosofía del template: los colores \texttt{Ctp*}
deben vivir \textit{únicamente} en \texttt{themes/catppuccin-palette.tex}.
El uso correcto es:
\begin{lstlisting}[style=latex, frame=none, numbers=none]
% INCORRECTO: Ctp* en un archivo de contenido
state/.style={draw=CtpBlue, fill=CtpBase, text=CtpText}

% CORRECTO: aliases semanticos
state/.style={draw=tikz-node-border, fill=background, text=maintext}
\end{lstlisting}
\end{observacion}

%───────────────────────────────────────────────────────────────────────────────
\section{Toggle de impresión en blanco y negro}
\label{sec:toggle-bn}
%───────────────────────────────────────────────────────────────────────────────

\begin{definicion}{Toggle \texttt{eqcolors}}
\label{def:toggle}
El toggle global \texttt{\textbackslash ifeqcolors} controla si las macros
aplican color o pasan el argumento sin modificar. Esto permite generar una
versión B/N del documento sin editar ninguna ecuación:
\end{definicion}

\begin{lstlisting}[style=latex, frame=none, numbers=none]
% En config/equations.tex (activo por defecto):
\eqcolorstrue    % color activado
\eqcolorsfalse   % color desactivado: todas las macros son transparentes

% Para desactivar en un grupo local:
{\eqcolorsfalse
  \begin{equation}
    i\hbar \frac{\partial \Psi}{\partial t} = \hat{H}\Psi
  \end{equation}
}
\end{lstlisting}

Con \texttt{\textbackslash eqcolorsfalse}, \texttt{\textbackslash eqop\{\textbackslash nabla\}}
produce $\nabla$ sin color --- funcionalmente idéntico a escribir
$\nabla$ directamente. El documento compila sin cambios.

%───────────────────────────────────────────────────────────────────────────────
\section{Integración con el paquete \texttt{physics}}
\label{sec:physics-pkg}
%───────────────────────────────────────────────────────────────────────────────

El paquete \texttt{physics}~\cite{physics-pkg} provee comandos de alto nivel
que se integran naturalmente con el coloreado semántico. La norma, el valor
absoluto y los operadores de Dirac se colorean envolviendo el comando de
\texttt{physics} en la macro semántica correcta:

\begin{lstlisting}[style=latex, frame=none, numbers=none]
\norm{\equ{f}}_{\eqdom{L^2}}         % norma de la incognita en L2
\abs{\eqk{z} - \eqk{z_0}}            % valor absoluto de datos
\bra{\eqk{\phi}}\hat{H}\ket{\equ{\psi}}   % elemento de matriz
\end{lstlisting}

\[
  \|\equ{f}\|_{\eqdom{L^2}}^2
  = \eqop{\int}_{\eqdom{\Omega}} \abs{\equ{f}(\eqvec{\mathbf{x}})}^2
    \eqdm{d\mu}(\eqvec{\mathbf{x}})
\]

\begin{observacion}{Conflicto de \texttt{\textbackslash div}}
El paquete \texttt{physics} redefine \texttt{\textbackslash div} como
el operador divergencia ($\nabla\cdot$). Para la divergencia coloreada
usar \texttt{\textbackslash eqdvg\{\}}, definido en
\texttt{config/equations.tex} precisamente para evitar este conflicto.
Para la división aritmética usar \texttt{\textbackslash frac\{\}\{\}} o
el operador \texttt{/}.
\end{observacion}

%───────────────────────────────────────────────────────────────────────────────
\printbibliography[heading=subbibliography]

%───────────────────────────────────────────────────────────────────────────────
\chapter{Código y algoritmos}
\label{ch:codigo}
%───────────────────────────────────────────────────────────────────────────────

El template distingue dos tipos de material computacional con herramientas
distintas. El paquete \texttt{listings}~\cite{listings2020} es para código
fuente ejecutable real: produce syntax highlighting completo y puede importar
archivos directamente desde el sistema de archivos. El paquete
\texttt{algorithm2e}~\cite{algorithm2e2017} es para pseudocódigo matemático
independiente del lenguaje: opera en modo texto con macros de coloreado
semántico y ecuaciones.

La distinción no es estética sino epistémica: un bloque de listings dice
\textit{``este código corre en Python 3.11''}; un bloque de algorithm2e
dice \textit{``este procedimiento existe en el nivel matemático, su
implementación es ortogonal a la exposición''}.

%───────────────────────────────────────────────────────────────────────────────
\section{Listings: código fuente ejecutable}
\label{sec:listings}
%───────────────────────────────────────────────────────────────────────────────

\subsection{Sintaxis básica}
\label{subsec:listings-sintaxis}

El bloque de código más simple usa un estilo predefinido y un caption:

\begin{lstlisting}[style=latex, frame=none, numbers=none]
\begin{lstlisting}[style=python, caption={Nombre descriptivo}, label={lst:etiqueta}]
# codigo aqui
end{lstlisting}   % <- sin backslash en el documento real
\end{lstlisting}

Los estilos disponibles son: \texttt{python}, \texttt{c}, \texttt{cpp},
\texttt{rust}, \texttt{go}, \texttt{julia}, \texttt{matlab},
\texttt{fortran}, \texttt{asm}, \texttt{mathematica}, y \texttt{latex}.
Todos heredan del estilo base \texttt{catppuccin-base} definido en
\texttt{config/listings.tex}.

\subsection{Importar archivos externos}
\label{subsec:lstinputlisting}

\begin{definicion}{lstinputlisting}
\label{def:lstinputlisting}
El comando \texttt{\textbackslash lstinputlisting} importa un archivo de
código directamente desde el sistema de archivos, sin necesidad de copiar
el contenido en el \texttt{.tex}. La carpeta \texttt{code/} del proyecto
es el repositorio canónico de todos los archivos de código.
\end{definicion}

\begin{lstlisting}[style=latex, frame=none, numbers=none]
% Archivo completo
\lstinputlisting[style=python,
  caption={Integrador de Euler para EDOs},
  label={lst:euler}]{code/euler.py}

% Fragmento: solo lineas 15-40
\lstinputlisting[style=python,
  firstline=15, lastline=40,
  caption={Funcion de paso del integrador},
  label={lst:euler-step}]{code/euler.py}

% Con numero de lineas relativo al archivo
\lstinputlisting[style=julia,
  firstnumber=15,
  firstline=15, lastline=40,
  caption={...}]{code/rk4.jl}
\end{lstlisting}

\begin{observacion}{Rutas relativas}
Las rutas en \texttt{\textbackslash lstinputlisting} son relativas al
archivo \texttt{main.tex}, no al archivo \texttt{.tex} del capítulo.
Si \texttt{main.tex} está en la raíz del proyecto, \texttt{code/euler.py}
apunta a \texttt{./code/euler.py} desde la raíz, independientemente de
en qué capítulo aparezca el comando.
\end{observacion}

%───────────────────────────────────────────────────────────────────────────────
\section{Galería de estilos por lenguaje}
\label{sec:galeria}
%───────────────────────────────────────────────────────────────────────────────

Los snippets siguientes son fragmentos reales de física computacional y
machine learning. Cada uno está elegido para ejercitar las distintas clases
de color: keywords estructurales, tipos, funciones, strings, comentarios,
macros y módulos.

\subsection{Python — Descenso de gradiente estocástico}
\label{subsec:python}

\begin{lstlisting}[style=python,
  caption={SGD con momentum para un modelo de regresión logística},
  label={lst:sgd}]
import numpy as np
from typing import Callable

def sgd_momentum(
    grad_fn: Callable[[np.ndarray], np.ndarray],
    theta0:  np.ndarray,
    lr:      float = 1e-3,
    beta:    float = 0.9,
    n_iter:  int   = 1000,
) -> np.ndarray:
    """Descenso de gradiente estocastico con momentum de Polyak."""
    theta = theta0.copy()
    v     = np.zeros_like(theta)   # momento acumulado

    for t in range(1, n_iter + 1):
        g  = grad_fn(theta)        # gradiente estocastico
        v  = beta * v + (1 - beta) * g
        theta -= lr * v            # actualizacion de parametros

    return theta
\end{lstlisting}

\subsection{C — Integrador de Euler para una EDO escalar}
\label{subsec:c}

\begin{lstlisting}[style=c,
  caption={Método de Euler explícito: $\equ{y}' = \eqfn{f}(\eqk{t}, \equ{y})$},
  label={lst:euler-c}]
#include <stdio.h>
#include <math.h>

/* f(t, y) = -lambda * y  =>  solucion exacta: y(t) = y0 * exp(-lambda*t) */
static double f(double t, double y, double lam) {
    (void)t;
    return -lam * y;
}

void euler(double y0, double t0, double tf,
           double lam, int n, double *out_t, double *out_y) {
    double h = (tf - t0) / n;
    out_t[0] = t0;
    out_y[0] = y0;

    for (int i = 0; i < n; ++i) {
        out_y[i+1] = out_y[i] + h * f(out_t[i], out_y[i], lam);
        out_t[i+1] = out_t[i] + h;
    }
}
\end{lstlisting}

\subsection{C++ — Resolución de sistema lineal con LU (Eigen)}
\label{subsec:cpp}

\begin{lstlisting}[style=cpp,
  caption={Factorización LU parcialmente pivotada con Eigen},
  label={lst:lu-cpp}]
#include <Eigen/Dense>
#include <iostream>

using Matrix = Eigen::MatrixXd;
using Vector = Eigen::VectorXd;

// Resuelve A x = b mediante descomposicion LU con pivoteo parcial.
// Lanza std::runtime_error si A es singular numericamente.
Vector solve_lu(const Matrix& A, const Vector& b) {
    auto lu = A.partialPivLu();
    if (std::abs(lu.determinant()) < 1e-12)
        throw std::runtime_error("Matriz singular: det(A) ~ 0");
    return lu.solve(b);
}

int main() {
    Matrix A = Matrix::Random(4, 4);
    Vector b = Vector::Random(4);
    Vector x = solve_lu(A, b);
    std::cout << "Residuo: " << (A * x - b).norm() << "\n";
}
\end{lstlisting}

\subsection{Rust — Integrador RK4 genérico}
\label{subsec:rust}

\begin{lstlisting}[style=rust,
  caption={Runge-Kutta de orden 4 genérico sobre \texttt{f64}},
  label={lst:rk4-rust}]
/// Avance de un paso del metodo RK4.
/// `f`: derivada del sistema  dy/dt = f(t, y)
/// Devuelve y(t + h).
pub fn rk4_step<F>(f: F, t: f64, y: f64, h: f64) -> f64
where
    F: Fn(f64, f64) -> f64,
{
    let k1 = f(t,           y);
    let k2 = f(t + 0.5*h,  y + 0.5*h*k1);
    let k3 = f(t + 0.5*h,  y + 0.5*h*k2);
    let k4 = f(t + h,       y + h*k3);
    y + (h / 6.0) * (k1 + 2.0*k2 + 2.0*k3 + k4)
}

#[cfg(test)]
mod tests {
    use super::*;
    #[test]
    fn decay_exact() {
        // y' = -y, y(0) = 1  =>  y(1) = e^{-1}
        let f = |_t: f64, y: f64| -y;
        let y1 = rk4_step(f, 0.0, 1.0, 1.0);
        assert!((y1 - (-1.0_f64).exp()).abs() < 1e-6);
    }
}
\end{lstlisting}

\subsection{Go — Caminata aleatoria 2D}
\label{subsec:go}

\begin{lstlisting}[style=go,
  caption={Simulación de caminata aleatoria 2D en Go},
  label={lst:rw-go}]
package main

import (
    "fmt"
    "math/rand"
)

type Point struct{ X, Y float64 }

// RandomWalk simula n pasos de una caminata aleatoria 2D.
// Cada paso tiene longitud 1 en una direccion aleatoria.
func RandomWalk(n int, seed int64) []Point {
    rng   := rand.New(rand.NewSource(seed))
    path  := make([]Point, n+1)
    path[0] = Point{0, 0}
    steps := [][2]float64{{1,0},{-1,0},{0,1},{0,-1}}

    for i := range n {
        s       := steps[rng.Intn(4)]
        path[i+1] = Point{path[i].X + s[0], path[i].Y + s[1]}
    }
    return path
}

func main() {
    path := RandomWalk(10_000, 42)
    fmt.Printf("Posicion final: (%.2f, %.2f)\n",
        path[len(path)-1].X, path[len(path)-1].Y)
}
\end{lstlisting}

\subsection{Julia — Método de potencias para el eigenvalor dominante}
\label{subsec:julia}

\begin{lstlisting}[style=julia,
  caption={Iteración de potencias: eigenvalor dominante de $\eqten{A}$},
  label={lst:power-julia}]
using LinearAlgebra

"""
    power_iteration(A, tol=1e-10, maxiter=1000) -> (lambda, v)

Calcula el eigenvalor dominante lambda y su eigenvector normalizado v
de la matriz simetrica A mediante iteracion de potencias.
"""
function power_iteration(A::AbstractMatrix{T};
                         tol::Float64 = 1e-10,
                         maxiter::Int = 1000) where T <: AbstractFloat
    n = size(A, 1)
    v = normalize(randn(n))      # vector inicial aleatorio
    lambda_old = zero(T)

    for _ in 1:maxiter
        w      = A * v           # multiplicacion matriz-vector
        lambda = dot(v, w)       # cociente de Rayleigh
        v      = normalize(w)
        abs(lambda - lambda_old) < tol && break
        lambda_old = lambda
    end
    return lambda_old, v
end
\end{lstlisting}

\subsection{MATLAB — Ecuación de calor 1D por diferencias finitas}
\label{subsec:matlab}

\begin{lstlisting}[style=matlab,
  caption={Ecuación de calor 1D: esquema de Crank-Nicolson [LeVeque 2007]},
  label={lst:heat-matlab}]
function U = heat_crank_nicolson(alpha, L, T, Nx, Nt)
% heat_crank_nicolson  Resuelve u_t = alpha * u_xx en [0,L] x [0,T]
%   Condiciones de frontera: u(0,t) = u(L,t) = 0
%   Condicion inicial:       u(x,0) = sin(pi*x/L)

dx = L / Nx;  dt = T / Nt;
r  = alpha * dt / (2 * dx^2);    % numero de Courant-Friedrich-Lewy

% Matrices tridiagonales del esquema implicito
e = ones(Nx-1, 1);
A = spdiags([-r*e, (1+2*r)*e, -r*e], [-1 0 1], Nx-1, Nx-1);
B = spdiags([ r*e, (1-2*r)*e,  r*e], [-1 0 1], Nx-1, Nx-1);

x = linspace(dx, L-dx, Nx-1)';
U = sin(pi * x / L);             % condicion inicial

for n = 1:Nt
    U = A \ (B * U);             % paso de tiempo Crank-Nicolson
end
end
\end{lstlisting}

\subsection{Fortran — Multiplicación matriz-vector con BLAS}
\label{subsec:fortran}

\begin{lstlisting}[style=fortran,
  caption={Llamada a \texttt{dgemv} de BLAS para $\eqvec{\mathbf{y}} \leftarrow \alpha\eqten{A}\eqvec{\mathbf{x}} + \beta\eqvec{\mathbf{y}}$},
  label={lst:blas-fortran}]
program matvec_blas
  implicit none
  integer,  parameter :: n = 4
  real(8),  parameter :: alpha = 1.0d0, beta = 0.0d0
  real(8)             :: A(n,n), x(n), y(n)
  integer             :: i

  ! Inicializar A y x
  A = reshape([(real(i,8), i=1,n*n)], [n,n])
  x = [(1.0d0, i=1,n)]
  y = 0.0d0

  ! y <- alpha*A*x + beta*y  (BLAS nivel 2)
  call dgemv('N', n, n, alpha, A, n, x, 1, beta, y, 1)

  write(*,'(A,*(F8.2))') 'y = ', y
end program matvec_blas
\end{lstlisting}

\subsection{Assembler x86-64 — Norma euclídea con SIMD}
\label{subsec:asm}

\begin{lstlisting}[style=asm,
  caption={Norma euclídea de un vector float con instrucciones SSE4},
  label={lst:norm-asm}]
; norm_f32(float* v, int n) -> float
; Calcula ||v||_2 = sqrt(sum v[i]^2) usando SSE4.1
; Argumento: rdi = puntero, esi = longitud (multiplo de 4)
section .text
global norm_f32
norm_f32:
    xorps   xmm0, xmm0          ; acumulador = 0.0
.loop:
    movups  xmm1, [rdi]         ; cargar 4 floats
    mulps   xmm1, xmm1          ; xmm1 = v[i]^2
    addps   xmm0, xmm1          ; acumular
    add     rdi, 16             ; avanzar 4 * sizeof(float)
    sub     esi, 4
    jg      .loop               ; continuar si quedan elementos

    haddps  xmm0, xmm0          ; reduccion horizontal
    haddps  xmm0, xmm0
    sqrtss  xmm0, xmm0          ; sqrt del escalar
    ret
\end{lstlisting}

\subsection{Mathematica — Solución simbólica de la ecuación de onda}
\label{subsec:mathematica}

\begin{lstlisting}[style=mathematica,
  caption={Solución simbólica de la ecuación de onda 1D con condiciones de Dirichlet},
  label={lst:wave-mathematica}]
(* Ecuacion de onda: u_tt = c^2 u_xx en [0,L], u(0,t)=u(L,t)=0 *)
With[{c = 1, L = Pi},
  u[x_, t_] = Sum[
    With[{omega = n*c*Pi/L},
      (A[n]*Cos[omega*t] + B[n]*Sin[omega*t]) * Sin[n*Pi*x/L]],
    {n, 1, 10}];

  (* Condicion inicial: u(x,0) = sin(x), u_t(x,0) = 0 *)
  coefs = Table[
    {n, 2/L * Integrate[Sin[x]*Sin[n*Pi*x/L], {x,0,L}]},
    {n, 1, 10}];

  Plot[u[x, 0.5] /. {A[n_] :> coefs[[n,2]], B[n_] :> 0},
       {x, 0, L}, PlotLabel -> "u(x, 0.5)"]
]
\end{lstlisting}

%───────────────────────────────────────────────────────────────────────────────
\section{Algoritmos: pseudocódigo matemático}
\label{sec:algoritmos}
%───────────────────────────────────────────────────────────────────────────────

\begin{definicion}{Cuándo usar algorithm2e vs listings}
\label{def:algo-vs-listings}
\texttt{algorithm2e} es el entorno correcto cuando el procedimiento existe
en el nivel matemático — cuando su descripción no depende de ningún lenguaje
de programación específico y cuando los símbolos del pseudocódigo son objetos
matemáticos con sus propios roles semánticos. Si el código se va a ejecutar
tal cual, usar \texttt{listings}.
\end{definicion}

\subsection{Sintaxis básica de algorithm2e}
\label{subsec:algo-sintaxis}

\begin{lstlisting}[style=latex, frame=none, numbers=none]
\begin{algorithm}[H]
  \DontPrintSemicolon
  \SetKwInOut{Input}{Entrada}\SetKwInOut{Output}{Salida}
  \Input{descripcion de entradas}
  \Output{descripcion de salidas}

  inicializacion\;
  \While{condicion}{
    cuerpo del bucle\;
    \If{condicion}{
      rama verdadera\;
    }
  }
  \Return{resultado}
  \caption{Nombre del algoritmo [clave-bib]}
  \label{alg:etiqueta}
\end{algorithm}
\end{lstlisting}

Los estilos de texto \texttt{\textbackslash algkw\{\}}, \texttt{\textbackslash algarg\{\}},
\texttt{\textbackslash algcomment\{\}} definidos en \texttt{config/algorithms.tex}
se aplican automáticamente vía \texttt{\textbackslash SetKwSty} y compañía.
No es necesario usarlos manualmente en el cuerpo del algoritmo.

\subsection{Theming semántico de ecuaciones en algoritmos}
\label{subsec:algo-ecuaciones}

La intersección más poderosa del template: pseudocódigo matemático donde
cada símbolo lleva su rol semántico. Las macros de \texttt{equations.tex}
funcionan dentro de \texttt{algorithm2e} en modo matemático normal.

\begin{algorithm}[H]
  \DontPrintSemicolon
  \SetKwInOut{Input}{Entrada}\SetKwInOut{Output}{Salida}
  \Input{%
    NFA $\eqk{N} = (\eqdom{Q}, \eqdom{\Sigma}, \eqop{\delta},
    \eqk{q_0}, \eqdom{F})$}
  \Output{%
    DFA $\eqk{D} = (\eqdom{Q'}, \eqdom{\Sigma}, \eqop{\delta'},
    \eqk{q'_0}, \eqdom{F'})$}

  $\eqk{q'_0} \gets \eqop{E}(\eqdom{\{}\eqk{q_0}\eqdom{\}})$\;
  Añadir $\eqk{q'_0}$ a $\eqdom{Q'}$ como estado no marcado\;

  \While{$\exists\, \equ{T} \in \eqdom{Q'}$ no marcado}{
    Marcar $\equ{T}$\;
    \For{cada $\equ{a} \in \eqdom{\Sigma}$}{
      $\equ{U} \gets \eqop{E}(\eqop{move}(\equ{T}, \equ{a}))$\;
      \If{$\equ{U} \notin \eqdom{Q'}$}{
        Añadir $\equ{U}$ a $\eqdom{Q'}$ no marcado\;
      }
      $\eqop{\delta'}(\equ{T}, \equ{a}) \gets \equ{U}$\;
    }
  }
  \Return{$\eqk{D}$}
  \caption{Construcción de subconjuntos: NFA $\to$ DFA~\cite{press2007numerical}}
  \label{alg:subset-construction}
\end{algorithm}

El color comunica la estructura lógica del algoritmo: $\eqdom{Q}$,
$\eqdom{\Sigma}$, $\eqdom{F}$ son conjuntos (dominios, sapphire).
$\eqk{q_0}$, $\eqk{N}$, $\eqk{D}$ son parámetros conocidos (mauve).
$\equ{T}$ y $\equ{U}$ son variables de trabajo (incógnitas locales, peach).
$\eqop{\delta}$ y $\eqop{E}$ son operadores/funciones (blue, green).

\subsection{Algoritmo numérico: gradiente conjugado}
\label{subsec:conjugate-gradient}

\begin{algorithm}[H]
  \DontPrintSemicolon
  \SetKwInOut{Input}{Entrada}\SetKwInOut{Output}{Salida}
  \Input{%
    Matriz $\eqk{A} \in \eqdom{\mathbb{R}^{n\times n}}$ SPD,
    vector $\eqk{b} \in \eqdom{\mathbb{R}^n}$,
    tolerancia $\eqk{\epsilon} > 0$}
  \Output{Solución aproximada $\equ{x} \approx \eqk{A}^{-1}\eqk{b}$}

  $\equ{x} \gets \mathbf{0}$\;
  $\equ{r} \gets \eqk{b} - \eqk{A}\equ{x}$\;
  $\equ{p} \gets \equ{r}$\;
  $\equ{\rho} \gets \equ{r}^\top \equ{r}$\;

  \While{$\sqrt{\equ{\rho}} > \eqk{\epsilon}$}{
    $\equ{q} \gets \eqk{A}\equ{p}$\;
    $\equ{\alpha} \gets \equ{\rho}\, /\, (\equ{p}^\top \equ{q})$\;
    $\equ{x} \gets \equ{x} + \equ{\alpha}\,\equ{p}$\;
    $\equ{r} \gets \equ{r} - \equ{\alpha}\,\equ{q}$\;
    $\equ{\rho_{\mathrm{new}}} \gets \equ{r}^\top \equ{r}$\;
    $\equ{p} \gets \equ{r} + (\equ{\rho_{\mathrm{new}}} / \equ{\rho})\,\equ{p}$\;
    $\equ{\rho} \gets \equ{\rho_{\mathrm{new}}}$\;
  }
  \Return{$\equ{x}$}
  \caption{Método del gradiente conjugado~\cite{higham2002accuracy}}
  \label{alg:cg}
\end{algorithm}

\begin{observacion}{Algoritmos largos y ruptura de página}
\texttt{algorithm2e} no envuelve el entorno \texttt{algorithm} en
\texttt{tcolorbox} (ver \cref{ch:estructura}). Esto permite que algoritmos
largos se dividan naturalmente entre páginas. Usar el placement \texttt{[H]}
solo cuando el algoritmo es corto y debe aparecer junto al texto que lo
introduce. Para algoritmos de más de 20 líneas, dejar el placement vacío y
permitir que \LaTeX{} lo posicione.
\end{observacion}

%───────────────────────────────────────────────────────────────────────────────
\section{Referenciando código y algoritmos}
\label{sec:ref-codigo}
%───────────────────────────────────────────────────────────────────────────────

Con \texttt{cleveref} configurado en \texttt{config/packages.tex}:

\begin{lstlisting}[style=latex, frame=none, numbers=none]
% Referenciar un bloque de listing
Como muestra el \cref{lst:sgd}, el momentum...

% Referenciar un algoritmo
El \cref{alg:cg} converge en a lo sumo $n$ iteraciones...

% Rango de listings
Los \crefrange{lst:euler-c}{lst:rk4-rust} ilustran...
\end{lstlisting}

Los nombres en español ya están configurados: \texttt{\textbackslash cref\{lst:sgd\}}
produce \textit{cód.~1} y \texttt{\textbackslash Cref\{alg:cg\}} produce
\textit{Algoritmo~2}.

%───────────────────────────────────────────────────────────────────────────────
\printbibliography[heading=subbibliography]

%───────────────────────────────────────────────────────────────────────────────
\chapter{Tablas}
\label{ch:tablas}
%───────────────────────────────────────────────────────────────────────────────

Una tabla bien diseñada comunica estructura antes de que el lector lea un
solo número. El template construye tablas sobre tres principios del paquete
\texttt{booktabs}~\cite{booktabs2020}: nunca líneas verticales, reglas
horizontales con grosor diferenciado, y espacio generoso entre filas. El
coloreado de \texttt{catppuccin-palette.tex} se aplica como capa semántica
encima de esa estructura, no como sustituto de ella.

%───────────────────────────────────────────────────────────────────────────────
\section{Anatomía de una tabla}
\label{sec:tabla-anatomia}
%───────────────────────────────────────────────────────────────────────────────

\begin{definicion}{Reglas de booktabs}
\label{def:booktabs}
\texttt{booktabs}~\cite{booktabs2020} define tres reglas horizontales con
grosor decreciente. \texttt{\textbackslash toprule} abre la tabla (0.08em),
\texttt{\textbackslash midrule} separa el encabezado del cuerpo (0.05em),
y \texttt{\textbackslash bottomrule} la cierra (0.08em). Las reglas
intermedias parciales usan \texttt{\textbackslash cmidrule\{i-j\}}.
Nunca se usan líneas verticales (\texttt{|} en el preámbulo de columnas).
\end{definicion}

La estructura canónica de cualquier tabla del template es:

\begin{lstlisting}[style=latex, frame=none, numbers=none]
\begin{acadtable}[H]{tab:etiqueta}{Caption descriptivo}
  \begin{tabularx}{\textwidth}{l Y Z}
    \toprule
    \headerrow
    \thead{Col A} & \thead{Col B} & \thead{Col C} \\
    \midrule
    dato          & dato          & dato          \\
    \bottomrule
  \end{tabularx}
\end{acadtable}
\end{lstlisting}

Tres comandos propios del template:
\texttt{\textbackslash headerrow} aplica \texttt{\textbackslash rowcolor\{table-head\}}
a la fila del encabezado.
\texttt{\textbackslash thead\{\}} pone el texto en bold con color
\texttt{table-head-text}.
\texttt{\textbackslash tablenote\{a\}\{texto\}} añade una nota al pie.

%───────────────────────────────────────────────────────────────────────────────
\section{Tipos de columna}
\label{sec:tipos-columna}
%───────────────────────────────────────────────────────────────────────────────

\texttt{config/tables.tex} define tipos de columna complementarios a los
estándar de \LaTeX\ y \texttt{tabularx}~\cite{tabularx1999}:

\begin{acadtable}[H]{tab:col-tipos}{Tipos de columna disponibles}
  \begin{tabularx}{\textwidth}{l l X}
    \toprule
    \headerrow
    \thead{Tipo} & \thead{Definición} & \thead{Uso} \\
    \midrule
    \texttt{l}, \texttt{c}, \texttt{r}
      & Estándar \LaTeX
      & Ancho fijo por contenido, sin wrapping \\
    \texttt{X}
      & \texttt{tabularx}
      & Proporcional, izquierda, wrapping \\
    \texttt{Y}
      & \texttt{tabularx}
      & Proporcional, centrada, wrapping \\
    \texttt{Z}
      & \texttt{tabularx}
      & Proporcional, derecha, wrapping \\
    \texttt{L\{ancho\}}
      & \texttt{array}
      & Ancho fijo, izquierda, wrapping \\
    \texttt{C\{ancho\}}
      & \texttt{array}
      & Ancho fijo, centrada, wrapping \\
    \texttt{R\{ancho\}}
      & \texttt{array}
      & Ancho fijo, derecha, wrapping \\
    \bottomrule
  \end{tabularx}
\end{acadtable}

\begin{observacion}{Columnas proporcionales X/Y/Z}
Las columnas \texttt{X}, \texttt{Y}, \texttt{Z} solo funcionan dentro de
\texttt{tabularx}. Distribuyen el espacio sobrante de \texttt{\textbackslash textwidth}
(o \texttt{\textbackslash linewidth} dentro de tcolorbox) entre todas las
columnas de ese tipo de forma proporcional. Una tabla con columnas \texttt{lYY}
tiene una columna de ancho natural a la izquierda y dos columnas de igual
ancho proporcional. Todas usan centrado vertical (\texttt{m\{\}}) para
alinear celdas con distinto número de líneas.
\end{observacion}

%───────────────────────────────────────────────────────────────────────────────
\section{Entornos de tabla}
\label{sec:entornos-tabla}
%───────────────────────────────────────────────────────────────────────────────

\subsection{acadtable — tabla flotante}
\label{subsec:acadtable}

\texttt{acadtable} es un wrapper de \texttt{table} con caption y label
integrados. Es el entorno por defecto para tablas en flujo normal de texto:

\begin{lstlisting}[style=latex, frame=none, numbers=none]
% Placement por defecto: [htbp]
\begin{acadtable}{tab:label}{Caption del flotante}
  \begin{tabularx}{\textwidth}{...} ... \end{tabularx}
\end{acadtable}

% Placement forzado: usa [H] cuando la tabla es parte del argumento
\begin{acadtable}[H]{tab:label}{Caption}
  ...
\end{acadtable}
\end{lstlisting}

\subsection{boxtab — tabla no-flotante dentro de tcolorbox}
\label{subsec:boxtab}

Dentro de cualquier entorno basado en \texttt{tcolorbox} (\texttt{definicion},
\texttt{ejemplo}, etc.), los flotantes no están permitidos. El entorno
\texttt{boxtab} resuelve esto:

\begin{lstlisting}[style=latex, frame=none, numbers=none]
\begin{ejemplo}{Tabla dentro de un entorno}
  \begin{boxtab}{tab:label}{Caption de la tabla}
    \begin{tabularx}{\linewidth}{l Y}   % linewidth, no textwidth
      ...
    \end{tabularx}
  \end{boxtab}
\end{ejemplo}
\end{lstlisting}

\subsection{fittable — reescalado automático}
\label{subsec:fittable}

Para tablas con muchas columnas que exceden el ancho disponible,
\texttt{fittable} reescala el contenido hasta que quepa:

\begin{lstlisting}[style=latex, frame=none, numbers=none]
\begin{fittable}
  \begin{tabular}{llllllll}
    % tabla muy ancha que desbordaria la pagina
  \end{tabular}
\end{fittable}
\end{lstlisting}

\begin{observacion}{fittable y legibilidad}
\texttt{fittable} es un recurso de último recurso, no el primer intento.
Antes de usarlo, verificar si la tabla puede restructurarse: eliminar
columnas redundantes, dividir en dos tablas, o usar \texttt{longtable}
para tablas altas. Un texto reescalado al 60\% de su tamaño original
suele ser ilegible en impresión.
\end{observacion}

%───────────────────────────────────────────────────────────────────────────────
\section{Coloreado semántico}
\label{sec:tabla-color}
%───────────────────────────────────────────────────────────────────────────────

Los aliases del dominio \texttt{table-*} en la paleta controlan todos los
colores de tablas:

\begin{acadtable}[H]{tab:table-aliases}{Aliases de color para tablas}
  \begin{tabularx}{\textwidth}{l l X}
    \toprule
    \headerrow
    \thead{Alias} & \thead{Color Ctp} & \thead{Aplicación} \\
    \midrule
    \texttt{table-head}      & Blue!12   & Fondo del encabezado (\texttt{\textbackslash headerrow}) \\
    \texttt{table-head-text} & Text      & Texto del encabezado (\texttt{\textbackslash thead}) \\
    \texttt{table-row-alt}   & Surface0!30 & Filas alternadas (\texttt{\textbackslash rowcolor}) \\
    \texttt{table-rule}      & Overlay0!50 & Color de todas las reglas booktabs \\
    \bottomrule
  \end{tabularx}
\end{acadtable}

Para filas alternadas, \texttt{\textbackslash rowcolor\{table-row-alt\}} se
aplica manualmente en las filas pares. Cuando la tabla tiene muchas filas,
el paquete \texttt{colortbl} permite automatizarlo con
\texttt{\textbackslash rowcolors}:

\begin{lstlisting}[style=latex, frame=none, numbers=none]
% Alternar colores desde la fila 2 (la primera es el encabezado)
\rowcolors{2}{table-row-alt}{white}
\begin{tabularx}{\textwidth}{l Y}
  \toprule
  \headerrow
  \thead{Col A} & \thead{Col B} \\
  \midrule
  dato 1 & valor 1 \\
  dato 2 & valor 2 \\
  dato 3 & valor 3 \\
  \bottomrule
\end{tabularx}
\rowcolors{1}{}{}   % resetear para no afectar tablas posteriores
\end{lstlisting}

%───────────────────────────────────────────────────────────────────────────────
\section{Datos numéricos con siunitx}
\label{sec:tabla-siunitx}
%───────────────────────────────────────────────────────────────────────────────

\begin{definicion}{Columna \texttt{S} de siunitx}
\label{def:columna-s}
El paquete \texttt{siunitx}~\cite{siunitx2023} provee el tipo de columna
\texttt{S} que alinea números por el punto decimal, redondea al número de
dígitos especificado, y formatea incertidumbres. Es el tipo correcto para
cualquier columna de datos numéricos con unidades o cifras significativas.
\end{definicion}

\begin{acadtable}[H]{tab:convergencia}{Convergencia del método de Runge-Kutta 4 para $\equ{y}' = \eqop{-}\eqk{\lambda}\equ{y}$}
  \begin{tabularx}{\textwidth}{Y S[table-format=1.2e-1] S[table-format=1.4e-2] S[table-format=1.2]}
    \toprule
    \headerrow
    \multicolumn{1}{Y}{\thead{$\eqk{h}$}} &
    \multicolumn{1}{c}{\thead{Error $\|\equ{e}\|_\infty$}} &
    \multicolumn{1}{c}{\thead{Razón $\equ{e}_n/\equ{e}_{n-1}$}} &
    \multicolumn{1}{c}{\thead{Orden emp.}} \\
    \midrule
    $1/4$   & 2.31e-3  & \multicolumn{1}{c}{---} & \multicolumn{1}{c}{---} \\
    $1/8$   & 1.48e-4  & 0.0641 & 3.96 \\
    $1/16$  & 9.33e-6  & 0.0630 & 3.99 \\
    $1/32$  & 5.84e-7  & 0.0626 & 4.00 \\
    $1/64$  & 3.65e-8  & 0.0625 & 4.00 \\
    \bottomrule
  \end{tabularx}
  \tablenote{*}{$\eqk{\lambda} = 1$, $\equ{y}(0) = 1$, integración en $\eqdom{[0,1]}$.
  El orden empírico confirma la convergencia $\mathcal{O}(\eqk{h}^4)$ teórica.}
\end{acadtable}

\begin{observacion}{Celdas especiales en columna S}
Cuando una celda de una columna \texttt{S} contiene texto en lugar de un
número (como las entradas \texttt{---} de la primera fila), debe encerrarse
entre llaves: \texttt{\{---\}}. De otro modo \texttt{siunitx} intenta
parsearla como número y produce un error.
\end{observacion}

%───────────────────────────────────────────────────────────────────────────────
\section{Tablas con estructura compleja}
\label{sec:tabla-compleja}
%───────────────────────────────────────────────────────────────────────────────

\subsection{Celdas multifila y multicolumna}
\label{subsec:multirow}

Para tablas donde las categorías abarcan varias filas, el paquete
\texttt{multirow}~\cite{multirow2021} provee \texttt{\textbackslash multirow}
y \LaTeX\ provee \texttt{\textbackslash multicolumn}:

\begin{acadtable}[H]{tab:hiperparametros}{Comparación de hiperparámetros por arquitectura}
  \begin{tabularx}{\textwidth}{l l Y S[table-format=2.2] S[table-format=1.4]}
    \toprule
    \headerrow
    \thead{Modelo} & \thead{Optimizador} & \thead{Regularización} &
    \multicolumn{1}{c}{\thead{Acc. (\%)}} &
    \multicolumn{1}{c}{\thead{Loss}} \\
    \midrule
    \multirow{3}{*}{\textbf{ResNet-18}}
      & SGD       & Weight decay $10^{-4}$ & 93.12 & 0.2341 \\
      & Adam      & Dropout 0.3            & 92.87 & 0.2418 \\
      & AdamW     & Weight decay $10^{-3}$ & 93.54 & 0.2289 \\
    \midrule
    \multirow{2}{*}{\textbf{MobileNetV3}}
      & Adam      & Dropout 0.2            & 90.23 & 0.3102 \\
      & AdamW     & Weight decay $10^{-3}$ & 91.05 & 0.2987 \\
    \bottomrule
  \end{tabularx}
\end{acadtable}

\subsection{Cmidrule para separadores parciales}
\label{subsec:cmidrule}

\texttt{\textbackslash cmidrule\{i-j\}} dibuja una regla horizontal entre
las columnas $i$ y $j$, útil para separar grupos dentro de una tabla sin
usar \texttt{\textbackslash midrule} completo:

\begin{acadtable}[H]{tab:metodos-numericos}{Métodos numéricos para EDOs: orden y estabilidad}
  \begin{tabularx}{\textwidth}{l l C{2cm} X}
    \toprule
    \headerrow
    \thead{Familia} & \thead{Método} & \thead{Orden} & \thead{Observaciones} \\
    \midrule
    \multirow{3}{*}{Explícitos}
      & Euler         & 1 & Primer orden; inestable para $\eqk{\lambda h} > 2$ \\
      & Heun          & 2 & Predictor-corrector; región de estabilidad mayor \\
      & RK4           & 4 & Estándar en precisión; $4$ evaluaciones por paso \\
    \cmidrule{2-4}
    \multirow{2}{*}{Implícitos}
      & Euler implícito & 1 & A-estable; requiere resolver sistema no lineal \\
      & Crank-Nicolson  & 2 & A-estable; exacto para osciladores armónicos \\
    \cmidrule{2-4}
    \multirow{2}{*}{Multipaso}
      & Adams-Bashforth & 4 & Explícito; reutiliza evaluaciones anteriores \\
      & Adams-Moulton   & 5 & Implícito; par con Adams-Bashforth en PECE \\
    \bottomrule
  \end{tabularx}
\end{acadtable}

%───────────────────────────────────────────────────────────────────────────────
\section{Tablas dentro de entornos tcolorbox}
\label{sec:tabla-en-caja}
%───────────────────────────────────────────────────────────────────────────────

Como se explicó en la \cref{def:no-floats}, dentro de entornos
\texttt{tcolorbox} los flotantes no son permitidos. El entorno \texttt{boxtab}
y el ancho \texttt{\textbackslash linewidth} resuelven esto:

\begin{definicion}{Propiedades de la distribución normal multivariante}
\label{def:normal-mv}
Sea $\equ{\mathbf{x}} \sim \mathcal{N}(\eqk{\boldsymbol{\mu}},
\eqten{\boldsymbol{\Sigma}})$ con
$\eqten{\boldsymbol{\Sigma}} \in \eqdom{\mathbb{R}^{d\times d}}$ SPD.

\begin{boxtab}{tab:normal-props}{Propiedades principales}
  \begin{tabularx}{\linewidth}{l X}
    \toprule
    \headerrow
    \thead{Propiedad} & \thead{Expresión} \\
    \midrule
    Media
      & $\eqop{\mathbb{E}}[\equ{\mathbf{x}}] = \eqk{\boldsymbol{\mu}}$ \\
    Covarianza
      & $\eqop{\mathrm{Cov}}[\equ{\mathbf{x}}] = \eqten{\boldsymbol{\Sigma}}$ \\
    Entropía
      & $H = \tfrac{1}{2}\ln\det(2\eqcon{\pi e}\,\eqten{\boldsymbol{\Sigma}})$ \\
    Marginalización
      & Toda marginal es normal; $\equ{x}_i \sim \mathcal{N}(\eqk{\mu_i}, \eqten{\Sigma_{ii}})$ \\
    Transformación lineal
      & $\eqten{A}\equ{\mathbf{x}} + \eqk{\mathbf{b}} \sim
        \mathcal{N}(\eqten{A}\eqk{\boldsymbol{\mu}} + \eqk{\mathbf{b}},\,
        \eqten{A}\eqten{\boldsymbol{\Sigma}}\eqten{A}^\top)$ \\
    \bottomrule
  \end{tabularx}
\end{boxtab}

La densidad de probabilidad es
$\eqfn{p}(\equ{\mathbf{x}}) = (2\eqcon{\pi})^{-d/2}
\det(\eqten{\boldsymbol{\Sigma}})^{-1/2}
\exp\!\left(-\tfrac{1}{2}(\equ{\mathbf{x}}-\eqk{\boldsymbol{\mu}})^\top
\eqten{\boldsymbol{\Sigma}}^{-1}(\equ{\mathbf{x}}-\eqk{\boldsymbol{\mu}})\right)$.
\end{definicion}

%───────────────────────────────────────────────────────────────────────────────
\section{Importar datos desde CSV con pgfplotstable}
\label{sec:pgfplotstable}
%───────────────────────────────────────────────────────────────────────────────

\begin{definicion}{pgfplotstable}
\label{def:pgfplotstable}
El paquete \texttt{pgfplotstable}~\cite{pgfplotstable2023}, incluido en la
distribución de \texttt{pgfplots}, permite cargar archivos \texttt{.csv}
directamente y renderizarlos como tablas \LaTeX\ con formato automático
de números, redondeo y unidades.
\end{definicion}

\begin{lstlisting}[style=latex, frame=none, numbers=none]
% Cargar el archivo una vez (en el preambulo o en el capitulo)
\pgfplotstableread[col sep=comma]{data/resultados.csv}\datatable

% Renderizar como tabla con formato automatico
\pgfplotstabletypeset[
  columns        = {h, error, ratio, order},
  columns/h/.style      = {column name={$h$}},
  columns/error/.style  = {sci, sci e, precision=2,
                            column name={Error $\|\mathbf{e}\|_\infty$}},
  columns/ratio/.style  = {fixed, precision=4,
                            column name={Razon}},
  columns/order/.style  = {fixed, precision=2,
                            column name={Orden emp.}},
  every head row/.style = {before row=\toprule, after row=\midrule},
  every last row/.style = {after row=\bottomrule},
  col sep=comma,
]{\datatable}
\end{lstlisting}

\begin{observacion}{Formato del CSV}
\texttt{pgfplotstable} espera que la primera fila del CSV contenga los
nombres de las columnas sin espacios. Los números deben usar punto decimal
(no coma), independientemente de la configuración de \texttt{siunitx}.
El separador puede ser coma (\texttt{col sep=comma}), tabulador
(\texttt{col sep=tab}) o espacio (\texttt{col sep=space}).
\end{observacion}

%───────────────────────────────────────────────────────────────────────────────
\section{Nota al pie de tabla}
\label{sec:tablenote}
%───────────────────────────────────────────────────────────────────────────────

El comando \texttt{\textbackslash tablenote\{marca\}\{texto\}} añade una
nota inmediatamente debajo de la tabla. No usa \texttt{\textbackslash footnote}
(que envía la nota al pie de página) sino texto inline en tamaño
\texttt{\textbackslash footnotesize}:

\begin{lstlisting}[style=latex, frame=none, numbers=none]
\begin{acadtable}[H]{tab:label}{Caption}
  \begin{tabularx}{\textwidth}{l Y}
    ...
    \bottomrule
  \end{tabularx}
  \tablenote{*}{Resultados promediados sobre 5 semillas aleatorias.}
  \tablenote{\dag}{Modelo pre-entrenado en ImageNet; fine-tuning 50 epochs.}
\end{acadtable}
\end{lstlisting}

Varias llamadas a \texttt{\textbackslash tablenote} se apilan verticalmente.
La marca puede ser cualquier texto: \texttt{*}, \texttt{\dag},
\texttt{a}, \texttt{1}, etc.

%───────────────────────────────────────────────────────────────────────────────
\printbibliography[heading=subbibliography]

%───────────────────────────────────────────────────────────────────────────────
\chapter{Figuras y gráficas}
\label{ch:figuras}
%───────────────────────────────────────────────────────────────────────────────

Las figuras en un documento académico cumplen dos funciones distintas:
las gráficas generadas por \texttt{pgfplots}~\cite{pgfplots2023} son parte
del argumento matemático --- muestran convergencia, distribuciones, campos
de fase --- y deben ser tan rigurosas como una ecuación. Las imágenes
externas importadas con \texttt{graphicx}~\cite{graphicx2021} son apoyo
visual: diagramas de arquitectura, fotografías de montajes, resultados de
simulaciones externas. El tratamiento tipográfico es el mismo pero su
naturaleza es diferente.

%───────────────────────────────────────────────────────────────────────────────
\section{Figuras externas}
\label{sec:figuras-externas}
%───────────────────────────────────────────────────────────────────────────────

\subsection{Figura simple}
\label{subsec:figura-simple}

La estructura canónica de una figura flotante en flujo normal de texto:

\begin{lstlisting}[style=latex, frame=none, numbers=none]
\begin{figure}[htbp]
  \centering
  \includegraphics[width=0.8\textwidth]{figs/nombre.pdf}
  \caption{Descripcion completa que se lee sin ver la figura.}
  \label{fig:etiqueta}
\end{figure}
\end{lstlisting}

\begin{observacion}{Formato de imagen recomendado}
LuaLaTeX acepta PDF, PNG y JPEG vía \texttt{\textbackslash includegraphics}.
Para gráficas vectoriales (diagramas, plots generados externamente), siempre
usar PDF — escala sin pérdida de calidad a cualquier resolución de impresión.
Para fotografías o capturas de pantalla, PNG. Los SVG deben convertirse a
PDF primero (\texttt{inkscape -{}-export-filename=out.pdf in.svg}).
\end{observacion}

\subsection{Figura dentro de tcolorbox}
\label{subsec:figura-en-caja}

Dentro de cualquier entorno basado en \texttt{tcolorbox}, usar \texttt{boxfig}
en lugar de \texttt{figure}. El ancho de referencia cambia de
\texttt{\textbackslash textwidth} a \texttt{\textbackslash linewidth}:

\begin{lstlisting}[style=latex, frame=none, numbers=none]
\begin{ejemplo}{Campo electrico de una carga puntual}
  \begin{boxfig}{fig:campo-e}{Lineas de campo de $\eqvec{\mathbf{E}}$}
    \includegraphics[width=0.7\linewidth]{figs/campo_electrico.pdf}
  \end{boxfig}
  La intensidad decae como $\eqfn{E}(r) \propto \eqk{q}/\equ{r}^2$.
\end{ejemplo}
\end{lstlisting}

\subsection{Subfiguras}
\label{subsec:subfiguras}

El paquete \texttt{subcaption}~\cite{subcaption2022} permite agrupar
múltiples figuras con subcaptions independientes:

\begin{lstlisting}[style=latex, frame=none, numbers=none]
\begin{figure}[htbp]
  \centering
  \begin{subfigure}[b]{0.48\textwidth}
    \centering
    \includegraphics[width=\textwidth]{figs/solucion_t0.pdf}
    \caption{$t = 0$: condicion inicial.}
    \label{fig:calor-t0}
  \end{subfigure}
  \hfill
  \begin{subfigure}[b]{0.48\textwidth}
    \centering
    \includegraphics[width=\textwidth]{figs/solucion_t1.pdf}
    \caption{$t = 0.5$: difusion parcial.}
    \label{fig:calor-t1}
  \end{subfigure}
  \caption{Evolucion de la ecuacion de calor
           $\eqop{\partial_t}\equ{u} = \eqk{\alpha}\eqop{\nabla^2}\equ{u}$
           con condiciones de Dirichlet.}
  \label{fig:calor-evolucion}
\end{figure}
\end{lstlisting}

La referencia a una subfigura usa el label de la subfigura:
\texttt{\textbackslash cref\{fig:calor-t0\}} produce \textit{fig.~7.1a}.
La referencia al grupo usa el label de la figura principal.

\subsection{Figuras embebidas en texto}
\label{subsec:wrapfig}

Para figuras pequeñas que deben fluir con el texto sin interrumpir el
párrafo, \texttt{wrapfig}~\cite{wrapfig2021} provee \texttt{wrapfigure}:

\begin{lstlisting}[style=latex, frame=none, numbers=none]
\begin{wrapfigure}{r}{0.35\textwidth}
  \centering
  \includegraphics[width=0.33\textwidth]{figs/diagrama.pdf}
  \caption{Diagrama del sistema.}
  \label{fig:diagrama}
\end{wrapfigure}
El texto circundante fluye alrededor de la figura. La opcion {r} la
coloca a la derecha; {l} a la izquierda. El ancho (0.35\textwidth)
debe ser ligeramente mayor que el de la imagen para dejar margen.
\end{lstlisting}

\begin{observacion}{Wrapfigure y listas}
\texttt{wrapfigure} no funciona correctamente dentro de entornos de lista
(\texttt{itemize}, \texttt{enumerate}) ni dentro de \texttt{tcolorbox}.
Tampoco debe aparecer al inicio de un capítulo antes de que exista
suficiente texto para fluir. Para esos contextos, usar \texttt{figure}
estándar.
\end{observacion}

%───────────────────────────────────────────────────────────────────────────────
\section{Gráficas con pgfplots}
\label{sec:pgfplots}
%───────────────────────────────────────────────────────────────────────────────

\begin{definicion}{Estilos Catppuccin para pgfplots}
\label{def:pgfplots-estilos}
\texttt{config/pgfplots.tex} define dos estilos con nombre y un ciclo
de colores de 8 series:
\begin{itemize}[leftmargin=2em]
  \item \texttt{catppuccin-clean} --- para publicación: grid mayor, fonts
    pequeñas, leyenda en esquina superior derecha.
  \item \texttt{catppuccin-filled} --- para notas o presentaciones: igual
    que \texttt{clean} pero con fondo \texttt{secondarybg}.
\end{itemize}
El ciclo de colores (\texttt{plot-series-1} a \texttt{plot-series-8})
garantiza que hasta 8 series sean visualmente distinguibles usando tanto
color como marcador distinto.
\end{definicion}

\subsection{Estructura básica de una gráfica}
\label{subsec:pgfplots-basico}

\begin{lstlisting}[style=latex, frame=none, numbers=none]
\begin{figure}[htbp]
  \centering
  \begin{tikzpicture}
    \begin{axis}[
      catppuccin-clean,
      xlabel = {$\eqk{t}$},
      ylabel = {$\equ{y}(\eqk{t})$},
      title  = {Decaimiento exponencial},
      xmin=0, xmax=5,
    ]
      \addplot[plot-series-1, thick] expression[domain=0:5]{exp(-x)};
      \addplot[plot-series-2, thick, dashed] expression[domain=0:5]{exp(-2*x)};
      \legend{$e^{-t}$, $e^{-2t}$}
    \end{axis}
  \end{tikzpicture}
  \caption{Comparacion de tasas de decaimiento.}
  \label{fig:decaimiento}
\end{figure}
\end{lstlisting}

\subsection{Gráfica con datos reales y barras de error}
\label{subsec:pgfplots-errorbar}

Para datos numéricos de simulaciones, \texttt{pgfplots} acepta coordenadas
directas con incertidumbre. El alias \texttt{plot-errorbar} controla el
color de las barras, que debe ser sutil para no competir con la serie:

\begin{lstlisting}[style=latex, frame=none, numbers=none]
\begin{axis}[
  catppuccin-clean,
  xlabel = {Numero de particulas $\eqk{N}$},
  ylabel = {$\langle \equ{E} \rangle / \eqk{N}$},
  xmode  = log,
]
  \addplot+[
    plot-series-1,
    error bars/.cd,
      y dir=both, y explicit,
      error bar style={plot-errorbar, thin},
  ] coordinates {
    (100,  -1.23) +- (0, 0.05)
    (500,  -1.41) +- (0, 0.02)
    (1000, -1.45) +- (0, 0.01)
    (5000, -1.47) +- (0, 0.005)
  };
\end{axis}
\end{lstlisting}

\subsection{Curva de nivel y mapa de calor}
\label{subsec:pgfplots-contour}

Los colormaps \texttt{catppuccin-heat} (magnitudes) y
\texttt{catppuccin-diverge} (residuos con signo) se definen en
\texttt{config/pgfplots.tex}:

\begin{lstlisting}[style=latex, frame=none, numbers=none]
% Mapa de calor de la funcion de densidad de probabilidad 2D
\begin{axis}[
  catppuccin-clean,
  colormap name  = catppuccin-heat,
  colorbar,
  view           = {0}{90},      % proyeccion superior
  xlabel         = {$\equ{x}_1$},
  ylabel         = {$\equ{x}_2$},
]
  \addplot3[surf, shader=interp, samples=40,
            domain=-3:3, domain y=-3:3]
    {exp(-0.5*(x^2 + y^2)) / (2*pi)};
\end{axis}

% Mapa divergente para residuos
\begin{axis}[
  colormap name = catppuccin-diverge,
  colorbar,
]
  \addplot3[surf, domain=-2:2] {x*y};
\end{axis}
\end{lstlisting}

\subsection{Gráfica 2D completa con theming semántico}
\label{subsec:pgfplots-completo}

Un ejemplo de gráfica funcional real con coloreado semántico en labels
y título. Muestra la convergencia del error de discretización
$\equ{e}_h = \|\equ{u}_h - \eqk{u}\|$ en función de $\eqk{h}$~\cite{tufte2001visual}:

\begin{lstlisting}[style=latex, frame=none, numbers=none]
\begin{figure}[htbp]
  \centering
  \begin{tikzpicture}
    \begin{loglogaxis}[
      catppuccin-clean,
      xlabel = {Paso de malla $\eqk{h}$},
      ylabel = {Error $\|\equ{u}_h - \eqk{u}\|_2$},
      title  = {Convergencia de diferencias finitas},
      legend pos = north west,
      xtick     = {0.0625, 0.125, 0.25, 0.5},
      xticklabels = {$1/16$, $1/8$, $1/4$, $1/2$},
    ]
      % Datos numericos: O(h^2)
      \addplot[plot-series-1, mark=*] coordinates {
        (0.5, 1.2e-2) (0.25, 3.1e-3) (0.125, 7.8e-4) (0.0625, 1.9e-4)
      };
      % Referencia teorica: pendiente -2 en log-log
      \addplot[plot-series-5, dashed, no markers]
        coordinates {(0.5, 1.2e-2) (0.0625, 1.9e-4)};
      \legend{FD centradas, $\mathcal{O}(\eqk{h}^2)$}
    \end{loglogaxis}
  \end{tikzpicture}
  \caption{Convergencia del esquema de diferencias finitas centradas
           para la ecuacion de Poisson 1D. La pendiente $-2$ en escala
           log-log confirma el orden teorico $\mathcal{O}(\eqk{h}^2)$.}
  \label{fig:convergencia-fd}
\end{figure}
\end{lstlisting}

\subsection{Subfiguras con pgfplots}
\label{subsec:pgfplots-subfig}

Para comparar resultados en múltiples paneles que comparten un tema
visual coherente:

\begin{lstlisting}[style=latex, frame=none, numbers=none]
\begin{figure}[htbp]
  \centering
  \begin{subfigure}[b]{0.48\textwidth}
    \begin{tikzpicture}
      \begin{axis}[catppuccin-clean, width=\textwidth, height=5cm,
                   title={Espacio de fase}, xlabel={$\equ{x}$},
                   ylabel={$\equ{\dot{x}}$}]
        \addplot[plot-series-1] table {data/fase.dat};
      \end{axis}
    \end{tikzpicture}
    \caption{Oscilador amortiguado: $\eqk{\gamma} < 2\eqk{\omega_0}$}
    \label{fig:fase-subamortiguado}
  \end{subfigure}
  \hfill
  \begin{subfigure}[b]{0.48\textwidth}
    \begin{tikzpicture}
      \begin{axis}[catppuccin-clean, width=\textwidth, height=5cm,
                   title={Serie temporal}, xlabel={$\eqk{t}$},
                   ylabel={$\equ{x}(\eqk{t})$}]
        \addplot[plot-series-2] table {data/tiempo.dat};
      \end{axis}
    \end{tikzpicture}
    \caption{Decaimiento amortiguado en el tiempo}
    \label{fig:tiempo-subamortiguado}
  \end{subfigure}
  \caption{Oscilador armónico subamortiguado: representación dual.}
  \label{fig:oscilador}
\end{figure}
\end{lstlisting}

%───────────────────────────────────────────────────────────────────────────────
\section{Captions: tipografía y colores}
\label{sec:captions}
%───────────────────────────────────────────────────────────────────────────────

El label del caption (``Figura 7.1'', ``Tabla 3.2'') usa el alias de color
del dominio \texttt{caption-*}, distinto por tipo de flotante. Esto permite
identificar visualmente de qué tipo es el elemento referenciado antes de
leer el número:

\begin{acadtable}[H]{tab:caption-colores}{Colores de labels de caption por tipo}
  \begin{tabularx}{\textwidth}{l l l X}
    \toprule
    \headerrow
    \thead{Tipo} & \thead{Alias} & \thead{Color Ctp} & \thead{Razón} \\
    \midrule
    Figura     & \texttt{caption-fig}  & \textcolor{caption-fig}{\textbf{Blue}}
      & = \texttt{link-internal}; coherente con referencias \\
    Tabla      & \texttt{caption-tab}  & \textcolor{caption-tab}{\textbf{Mauve}}
      & = \texttt{env-thm}; distinguible de figura \\
    Código     & \texttt{caption-code} & \textcolor{caption-code}{\textbf{Green}}
      & = \texttt{code-string}; evoca código \\
    Algoritmo  & \texttt{caption-alg}  & \textcolor{caption-alg}{\textbf{Peach}}
      & Cálido; distinguible de los anteriores \\
    Subfigura  & \texttt{caption-sub}  & \textcolor{caption-sub}{\textbf{Sapphire}}
      & = \texttt{env-phys}; un nivel por debajo de la figura \\
    \midrule
    \multicolumn{2}{l}{Texto descriptivo} & \textcolor{caption-text}{\textbf{Subtext0}}
      & Secundario; no compite con el label \\
    \bottomrule
  \end{tabularx}
\end{acadtable}

El texto descriptivo del caption (todo lo que sigue al label) usa
\texttt{caption-text} en gris suave para que el label coloreado ancle
visualmente el flotante sin que el texto completo compita con el contenido.

%───────────────────────────────────────────────────────────────────────────────
\section{linewidth vs textwidth}
\label{sec:linewidth}
%───────────────────────────────────────────────────────────────────────────────

\begin{definicion}{linewidth dentro de tcolorbox}
\label{def:linewidth}
Dentro de un entorno \texttt{tcolorbox}, \texttt{\textbackslash linewidth}
es el ancho del área de contenido de la caja (descontando los márgenes
\texttt{left=6pt} y \texttt{right=6pt} de \texttt{basebox}).
\texttt{\textbackslash textwidth} es siempre el ancho de la columna de texto
completa, independientemente del contexto. Usar \texttt{\textbackslash textwidth}
dentro de una caja produce overflow.
\end{definicion}

\begin{acadtable}[H]{tab:linewidth-contexto}{Ancho correcto según contexto}
  \begin{tabularx}{\textwidth}{l l X}
    \toprule
    \headerrow
    \thead{Contexto} & \thead{Ancho correcto} & \thead{Aplica a} \\
    \midrule
    Flujo normal    & \texttt{\textbackslash textwidth}
      & \texttt{figure}, \texttt{acadtable}, \texttt{tabularx} \\
    Dentro de tcolorbox & \texttt{\textbackslash linewidth}
      & \texttt{boxfig}, \texttt{boxtab}, \texttt{tabularx} \\
    Subfigura       & \texttt{\textbackslash textwidth} del subfigure
      & \texttt{\textbackslash includegraphics}, \texttt{tikzpicture} \\
    Wrapfigure      & ancho declarado en \texttt{wrapfigure}
      & imagen interna \\
    \bottomrule
  \end{tabularx}
\end{acadtable}

%───────────────────────────────────────────────────────────────────────────────
\section{PlaceFloat y FloatBarrier}
\label{sec:float-control}
%───────────────────────────────────────────────────────────────────────────────

El algoritmo de flotantes de \LaTeX\ puede acumular figuras y tablas no
colocadas hasta el final de una sección, dejando páginas desequilibradas.
Dos herramientas del template controlan esto:

\begin{lstlisting}[style=latex, frame=none, numbers=none]
% FloatBarrier: fuerza la colocacion de todos los flotantes pendientes
% antes de continuar. Util al inicio de una nueva seccion.
\FloatBarrier

% Placement [H] (paquete float): coloca el flotante exactamente aqui,
% sin algoritmo de flotantes. Usar cuando la figura es parte del argumento
% y debe aparecer en ese punto exacto del texto.
\begin{figure}[H]
  ...
\end{figure}
\end{lstlisting}

\begin{observacion}{Cuándo usar \texttt{[H]} y cuándo \texttt{[htbp]}}
\texttt{[H]} es para figuras que son parte del argumento narrativo ---
el lector necesita verlas en ese punto exacto para seguir el razonamiento.
\texttt{[htbp]} es para figuras que ilustran o complementan el argumento ---
el lector puede verlas antes o después sin perder el hilo. El uso excesivo
de \texttt{[H]} produce páginas con grandes espacios en blanco cuando la
figura no cabe en el espacio restante. \texttt{\textbackslash FloatBarrier}
al inicio de cada \texttt{\textbackslash section} es un buen balance.
\end{observacion}

%───────────────────────────────────────────────────────────────────────────────
\printbibliography[heading=subbibliography]

%───────────────────────────────────────────────────────────────────────────────
\chapter{Diagramas}
\label{ch:diagramas}
%───────────────────────────────────────────────────────────────────────────────

Los diagramas en un documento académico no son decoración: son un lenguaje
formal con gramática propia. Un autómata finito es una definición matemática
en forma visual. Un árbol de sintaxis abstracta es la estructura semántica
de un programa. Un diagrama conmutativo es una proposición sobre morfismos.
El template~\cite{tikz2023} trata los diagramas con la misma precisión
que las ecuaciones.

%───────────────────────────────────────────────────────────────────────────────
\section{Filosofía de theming en diagramas}
\label{sec:tikz-filosofia}
%───────────────────────────────────────────────────────────────────────────────

\begin{definicion}{Aliases tikz-* como valores por defecto}
\label{def:tikz-aliases}
El dominio \texttt{tikz-*} de la paleta define valores por defecto para
nodos, aristas y labels genéricos. Son el punto de partida — no el único
color posible. Cuando un nodo tiene significado semántico propio
(estado aceptador, colisión, estado inicial), se sobreescribe con el
alias del dominio correspondiente: \texttt{tikz-accept}, \texttt{fb-error},
\texttt{tikz-initial}.
\end{definicion}

\begin{acadtable}[H]{tab:tikz-aliases}{Aliases semánticos para diagramas TikZ}
  \begin{tabularx}{\textwidth}{l l l X}
    \toprule
    \headerrow
    \thead{Alias} & \thead{Color Ctp} & \thead{Rol} & \thead{Uso} \\
    \midrule
    \texttt{tikz-node-bg}     & Surface0  & Fondo genérico
      & Nodo sin semántica propia \\
    \texttt{tikz-node-border} & Overlay0  & Borde genérico
      & Nodo estándar \\
    \texttt{tikz-edge}        & Overlay1  & Arista
      & Conexión sin dirección semántica \\
    \texttt{tikz-label}       & Subtext0  & Label
      & Anotación secundaria \\
    \texttt{tikz-initial}     & Blue      & Estado inicial
      & Autómatas: punto de entrada \\
    \texttt{tikz-accept}      & Green     & Estado aceptador
      & Autómatas: estado final \\
    \texttt{tikz-reject}      & Red       & Estado rechazador
      & Autómatas / error \\
    \texttt{tikz-highlight}   & Yellow    & Nodo destacado
      & Paso activo en explicaciones \\
    \midrule
    \texttt{fb-error}         & Red       & Error / colisión
      & Condición de error en estructuras \\
    \texttt{eq-known}         & Mauve     & Dato conocido
      & Pesos, etiquetas de aristas con valor fijo \\
    \texttt{env-phys}         & Sapphire  & Transformación
      & Cajas de función / operación \\
    \bottomrule
  \end{tabularx}
\end{acadtable}

La regla es consistente con el resto del template: \textbf{nunca} usar
\texttt{Ctp*} directamente en archivos de contenido. Cuando ningún alias
existente cubre el rol semántico necesario, se añade uno nuevo en
\texttt{themes/catppuccin-palette.tex} bajo el dominio \texttt{tikz-*}.

%───────────────────────────────────────────────────────────────────────────────
\section{Autómatas finitos}
\label{sec:automatas}
%───────────────────────────────────────────────────────────────────────────────

\begin{definicion}{Estilo base para autómatas}
\label{def:automata-style}
El estilo canónico para autómatas en el template usa la librería
\texttt{automata} de TikZ con los aliases semánticos del dominio
\texttt{tikz-*}. El estado inicial usa \texttt{tikz-initial} (Blue),
los estados finales \texttt{tikz-accept} (Green), y las aristas
\texttt{tikz-edge} (Overlay1). Las etiquetas de transición aplican
coloreado semántico de ecuaciones donde corresponda.
\end{definicion}

\begin{lstlisting}[style=latex, frame=none, numbers=none]
\tikzset{
  state/.style = {
    circle,
    draw  = tikz-node-border,
    fill  = background,
    text  = maintext,
    thick,
    minimum size = 1cm,
  },
  every edge/.style = {draw=tikz-edge, thick, ->, >=latex},
}
\end{lstlisting}

\subsection{DFA: autómata determinista}
\label{subsec:dfa}

DFA para el lenguaje $L = \{\eqvec{w} \in \eqdom{\{a,b\}^*} \mid
\eqvec{w}$ contiene el sufijo $\eqk{abb}\}$~\cite{sipser2012introduction}:

\begin{figure}[htbp]
  \centering
  \begin{tikzpicture}[
    shorten >=1pt, node distance=2.5cm, on grid, auto, thick,
    state/.style={circle, draw=tikz-node-border, fill=background,
                  text=maintext, minimum size=1cm},
    every edge/.style={draw=tikz-edge, ->, >=latex},
  ]
    \node[state, initial, initial text={}, draw=tikz-initial]
      (A) {$\eqk{A}$};
    \node[state, right=of A] (B) {$\eqk{B}$};
    \node[state, right=of B] (D) {$\eqk{D}$};
    \node[state, accepting, draw=tikz-accept, right=of D]
      (E) {$\eqk{E}$};

    \path[->]
      (A) edge[loop above] node {$\eqk{b}$} (A)
          edge              node {$\eqk{a}$} (B)
      (B) edge[loop above] node {$\eqk{a}$} (B)
          edge              node {$\eqk{b}$} (D)
      (D) edge              node {$\eqk{b}$} (E)
          edge[bend left]   node {$\eqk{a}$} (B)
      (E) edge              node[swap] {$\eqk{a}$} (B)
          edge[bend right]  node[swap] {$\eqk{b}$} (A);
  \end{tikzpicture}
  \caption{DFA mínimo para $L = \{\eqvec{w} \in \eqdom{\{a,b\}^*} \mid
           \eqvec{w}$ termina en $\eqk{abb}\}$.
           Estado inicial en azul; estado aceptador en verde.}
  \label{fig:dfa-abb}
\end{figure}

\subsection{NFA con transiciones $\varepsilon$}
\label{subsec:nfa}

Construcción de Thompson para $(\eqk{a} \eqop{\mid} \eqk{b})\eqop{*}\eqk{abb}$.
Los estados intermedios usan el color por defecto; el estado aceptador
usa \texttt{tikz-accept}:

\begin{figure}[htbp]
  \centering
  \begin{tikzpicture}[
    shorten >=1pt, node distance=2cm, on grid, auto, thick,
    state/.style={circle, draw=tikz-node-border, fill=background,
                  text=maintext, minimum size=0.9cm},
    every edge/.style={draw=tikz-edge, ->, >=latex},
  ]
    \node[state, initial, initial text={}, draw=tikz-initial]
      (0) {$\eqk{0}$};
    \node[state, right=of 0] (1) {$\eqk{1}$};
    \node[state, right=of 1] (2) {$\eqk{2}$};
    \node[state, accepting, draw=tikz-accept, right=of 2]
      (3) {$\eqk{3}$};

    \path[->]
      (0) edge[loop above] node {$\eqk{a}\eqop{,}\eqk{b}$} (0)
          edge              node {$\eqk{a}$}                 (1)
      (1) edge              node {$\eqk{b}$}                 (2)
      (2) edge              node {$\eqk{b}$}                 (3);
  \end{tikzpicture}
  \caption{NFA simplificado para $({\eqk{a} \eqop{|} \eqk{b}})^* \eqk{abb}$.
           El estado aceptador $\eqk{3}$ en verde distingue el criterio
           de aceptación.}
  \label{fig:nfa-abb}
\end{figure}

%───────────────────────────────────────────────────────────────────────────────
\section{Árboles}
\label{sec:arboles}
%───────────────────────────────────────────────────────────────────────────────

\subsection{Árbol binario de búsqueda con TikZ}
\label{subsec:bst}

Los árboles binarios se construyen con TikZ usando el estilo de nodo
estándar. El nodo raíz puede destacarse con \texttt{tikz-highlight}:

\begin{figure}[htbp]
  \centering
  \begin{tikzpicture}[
    level distance  = 1.5cm,
    level 1/.style  = {sibling distance=4cm},
    level 2/.style  = {sibling distance=2cm},
    level 3/.style  = {sibling distance=1.2cm},
    every node/.style = {
      circle, draw=tikz-node-border, fill=tikz-node-bg,
      text=maintext, minimum size=0.8cm, thick,
    },
    edge from parent/.style = {draw=tikz-edge, thick},
  ]
    \node[fill=tikz-highlight] {$\eqk{50}$}
      child { node {$\eqk{30}$}
        child { node {$\eqk{20}$}
          child { node {$\eqk{10}$} }
          child[missing] {}
        }
        child { node {$\eqk{40}$} }
      }
      child { node {$\eqk{70}$}
        child { node {$\eqk{60}$} }
        child { node {$\eqk{80}$}
          child[missing] {}
          child { node {$\eqk{90}$} }
        }
      };
  \end{tikzpicture}
  \caption{BST con raíz $\eqk{50}$ destacada. La propiedad de orden:
           todo nodo izquierdo $< \eqk{k} <$ nodo derecho~\cite{cormen2022algorithms}.}
  \label{fig:bst}
\end{figure}

\subsection{Árbol sintáctico con forest}
\label{subsec:parse-tree}

El paquete \texttt{forest}~\cite{forest2017} es más expresivo que TikZ
para árboles con muchos niveles. Ideal para árboles de derivación gramatical
y árboles de expresión:

\begin{lstlisting}[style=latex, frame=none, numbers=none]
\begin{forest}
  for tree={
    draw          = tikz-node-border,
    fill          = tikz-node-bg,
    text          = maintext,
    rounded corners = 2pt,
    s sep         = 1.5em,
    l sep         = 2em,
    edge          = {draw=tikz-edge, thick},
  }
  [$\eqop{+}$
    [$\eqop{\times}$
      [$\eqk{3}$]
      [$\eqk{4}$]
    ]
    [$\eqk{5}$]
  ]
\end{forest}
\end{lstlisting}

\begin{observacion}{qtree para árboles simples}
Para árboles lingüísticos o de derivación sencillos, \texttt{qtree} ofrece
una sintaxis más compacta:
\begin{lstlisting}[style=latex, frame=none, numbers=none]
\Tree [.S [.NP [.Det la ] [.N funcion ] ]
          [.VP [.V converge ] [.PP [.P en ] [.NP \(L^2\) ] ] ] ]
\end{lstlisting}
\texttt{forest} es preferible cuando se necesita control fino sobre
colores de nodos individuales o layouts no estándar.
\end{observacion}

%───────────────────────────────────────────────────────────────────────────────
\section{Estructuras de datos}
\label{sec:estructuras-datos}
%───────────────────────────────────────────────────────────────────────────────

Los diagramas de estructuras de datos (arrays, tablas hash, listas enlazadas)
comparten un vocabulario visual: celdas rectangulares para memoria, flechas
para punteros, color para estado semántico (libre, ocupado, colisión).

\subsection{Tabla hash con separate chaining}
\label{subsec:hash-chaining}

El diagrama muestra una tabla hash donde el índice $\eqop{h}(\eqk{k}) =
\eqk{k} \eqop{\bmod} \eqk{m}$ mapea claves a buckets. Los buckets
ocupados usan \texttt{tikz-accept} (Green); los vacíos, \texttt{tikz-node-bg}
(Surface0); las colisiones resueltas por chaining extienden la cadena
hacia la derecha~\cite{cormen2022algorithms}:

\begin{figure}[htbp]
  \centering
  \begin{tikzpicture}[font=\small,
    cell/.style = {
      draw=tikz-node-border, minimum width=1.4cm, minimum height=0.7cm,
    },
    occ/.style  = {cell, fill=tikz-accept,   text=background},
    empty/.style= {cell, fill=tikz-node-bg,  text=maintext},
    key/.style  = {cell, fill=tikz-initial,  text=background},
    arr/.style  = {draw=tikz-edge, ->, >=latex, thick},
  ]
    % --- Buckets (índices 0-7) ---
    \foreach \i in {0,...,7} {
      \node[empty] (b\i) at (0, {-\i*0.9}) {};
      \node[left=0.15cm of b\i, text=maintext] {\(\i\)};
    }

    % --- Buckets ocupados ---
    \node[occ] at (b1) {\(\eqk{21}\)};
    \node[occ] at (b3) {\(\eqk{12}\)};

    % --- Chains desde índice 1: 21→31→41 ---
    \node[occ] (n1a) at (2.2, {-1*0.9}) {\(\eqk{31}\)};
    \node[occ] (n1b) at (4.0, {-1*0.9}) {\(\eqk{41}\)};
    \draw[arr] (b1)  -- (n1a);
    \draw[arr] (n1a) -- (n1b);

    % --- Chain desde índice 3: 12→22 ---
    \node[occ] (n3a) at (2.2, {-3*0.9}) {\(\eqk{22}\)};
    \draw[arr] (b3) -- (n3a);

    % --- Etiquetas ---
    \node[above=0.4cm of b0, text=maintext] {\textbf{Índice}};
    \node[above=0.4cm of n1a, text=maintext] {\textbf{Cadena}};
  \end{tikzpicture}
  \caption{Separate chaining: los buckets ocupados (verde) encadenan
           los elementos con el mismo índice hash. Vacíos en gris.}
  \label{fig:chaining}
\end{figure}

\subsection{Linear probing y colisión}
\label{subsec:linear-probing}

En linear probing, una colisión en el índice $\eqop{h}(\eqk{k})$ activa
una búsqueda lineal del siguiente slot libre. El slot de colisión usa
\texttt{tikz-highlight} para visualizar el conflicto; el slot final usa
\texttt{tikz-accept}:

\begin{figure}[htbp]
  \centering
  \begin{tikzpicture}[font=\small,
    cell/.style  = {draw=tikz-node-border, minimum width=0.9cm,
                    minimum height=0.65cm},
    occ/.style   = {cell, fill=tikz-accept,    text=background},
    coll/.style  = {cell, fill=tikz-highlight, text=maintext},
    empty/.style = {cell, fill=tikz-node-bg,   text=maintext},
    arr/.style   = {draw=tikz-edge, ->, >=latex, thick, dashed},
  ]
    % Array de 10 celdas
    \foreach \i in {0,...,9} {
      \node[empty] (a\i) at ({\i*1.1}, 0) {};
      \node[below=0.15cm of a\i, text=tikz-label, font=\scriptsize] {\(\i\)};
    }

    % Valor 15 en índice 5 (sin colisión)
    \node[occ] at (a5) {\(\eqk{15}\)};

    % Valor 25: colisiona en 5, va a 6
    \node[coll] at (a5) {\(\eqk{15}\)};   % ya ocupado
    \node[occ]  at (a6) {\(\eqk{25}\)};

    % Flecha de probing
    \draw[arr] (a5.north) to[bend left=40]
      node[above, text=maintext, font=\scriptsize] {probing}
      (a6.north);

    % Etiquetas semánticas
    \node[above=1.2cm of a5, text=tikz-highlight]
      {\small colisión: $\eqop{h}(\eqk{25}) = \eqk{5}$};
    \node[above=1.2cm of a6, text=tikz-accept]
      {\small colocado en $\eqk{6}$};
  \end{tikzpicture}
  \caption{Linear probing: $\eqk{25}$ colisiona en índice $\eqop{h}(\eqk{25}) = \eqk{5}$
           y se desplaza al siguiente slot libre $\eqk{6}$.}
  \label{fig:linear-probing}
\end{figure}

\subsection{Diagrama de hashing completo}
\label{subsec:hash-completo}

El diagrama más informativo muestra claves, función hash y buckets
como tres columnas conectadas. La función hash usa \texttt{env-phys}
(Sapphire) por ser una transformación matemática; las claves de entrada
usan \texttt{tikz-initial} (Blue):

\begin{figure}[htbp]
  \centering
  \begin{tikzpicture}[font=\small,
    key/.style   = {draw=tikz-node-border, fill=tikz-initial,
                    text=background, minimum width=2.2cm,
                    minimum height=0.7cm},
    fn/.style    = {draw=tikz-node-border, fill=env-phys,
                    text=background, minimum width=1.3cm,
                    minimum height=3.8cm, align=center},
    occ/.style   = {draw=tikz-node-border, fill=tikz-accept,
                    text=background, minimum width=1.4cm,
                    minimum height=0.55cm},
    empty/.style = {draw=tikz-node-border, fill=tikz-node-bg,
                    text=maintext, minimum width=1.4cm,
                    minimum height=0.55cm},
    arr/.style   = {draw=tikz-edge, ->, >=latex, thick},
  ]
    % --- Claves ---
    \node[key] (k1) at (-4.5,  1.2) {John Smith};
    \node[key] (k2) at (-4.5,  0.0) {Lisa Smith};
    \node[key] (k3) at (-4.5, -1.2) {Sandra Dee};

    % --- Función hash ---
    \node[fn] (hf) at (-2, 0) {$\eqop{h}$\\function};

    % --- Buckets ---
    \foreach \i in {0,...,5} {
      \node[empty] (bk\i) at (1, {1.5 - \i*0.6}) {};
      \node[left=0.1cm of bk\i, text=tikz-label,
            font=\scriptsize] {\(\i\)};
    }
    \node[occ] at (bk1) {\footnotesize 521-8976};
    \node[occ] at (bk2) {\footnotesize 521-1234};
    \node[occ] at (bk4) {\footnotesize 521-9655};

    % --- Aristas claves → hash ---
    \draw[arr] (k1) -- (hf);
    \draw[arr] (k2) -- (hf);
    \draw[arr] (k3) -- (hf);

    % --- Aristas hash → buckets ---
    \draw[arr] (hf.east) to[out=0, in=180] (bk1.west);
    \draw[arr] (hf.east) to[out=0, in=180] (bk2.west);
    \draw[arr] (hf.east) to[out=0, in=180] (bk4.west);

    % --- Etiquetas ---
    \node[above=0.4cm of k1,   text=maintext] {\textbf{Claves}};
    \node[above=0.4cm of hf,   text=maintext] {\textbf{Función hash}};
    \node[above=0.4cm of bk0,  text=maintext] {\textbf{Buckets}};
  \end{tikzpicture}
  \caption{Hashing: claves (azul) transformadas por
           $\eqop{h}$ (sapphire) a buckets (verde=ocupado, gris=libre).
           La función hash tiene rol de transformación, análogo a un
           operador físico.}
  \label{fig:hashing}
\end{figure}

%───────────────────────────────────────────────────────────────────────────────
\section{Grafos}
\label{sec:grafos}
%───────────────────────────────────────────────────────────────────────────────

\subsection{Grafo dirigido ponderado}
\label{subsec:grafo-dirigido}

Los pesos de aristas son datos conocidos del problema — usan
\texttt{eq-known} (Mauve). Los nodos son genéricos salvo el nodo fuente
que usa \texttt{tikz-initial} y el nodo destino \texttt{tikz-accept}:

\begin{figure}[htbp]
  \centering
  \begin{tikzpicture}[
    node distance = 2.5cm,
    vertex/.style = {circle, draw=tikz-node-border, fill=tikz-node-bg,
                     text=maintext, minimum size=0.9cm, thick},
    every edge/.style = {draw=tikz-edge, ->, >=latex, thick},
    wlabel/.style = {text=eq-known, font=\small},
  ]
    \node[vertex, draw=tikz-initial, fill=tikz-initial!15] (s) {$\eqk{s}$};
    \node[vertex, right=of s]  (u) {$\eqk{u}$};
    \node[vertex, right=of u]  (v) {$\eqk{v}$};
    \node[vertex, below=of u]  (x) {$\eqk{x}$};
    \node[vertex, draw=tikz-accept, fill=tikz-accept!15,
          right=of x]          (t) {$\eqk{t}$};

    \path[->]
      (s) edge node[above, wlabel] {$\eqk{6}$} (u)
      (s) edge node[left,  wlabel] {$\eqk{7}$} (x)
      (u) edge node[above, wlabel] {$\eqk{5}$} (v)
      (u) edge node[right, wlabel] {$\eqk{8}$} (x)
      (u) edge node[right, wlabel] {$\eqk{-4}$}(t)
      (x) edge node[right, wlabel] {$\eqk{9}$} (v)
      (x) edge node[below, wlabel] {$\eqk{7}$} (t)
      (v) edge node[right, wlabel] {$\eqk{-2}$}(t);
  \end{tikzpicture}
  \caption{Grafo dirigido ponderado. Fuente $\eqk{s}$ (azul),
           destino $\eqk{t}$ (verde). Pesos en malva
           como datos conocidos del problema de camino mínimo~\cite{cormen2022algorithms}.}
  \label{fig:grafo-ponderado}
\end{figure}

\subsection{Nodo activo durante ejecución de algoritmo}
\label{subsec:grafo-highlight}

Para ilustrar la ejecución paso a paso de Dijkstra u otro algoritmo
sobre un grafo, \texttt{tikz-highlight} marca el nodo activo en cada
paso. El patrón es sobreescribir el estilo del nodo localmente:

\begin{lstlisting}[style=latex, frame=none, numbers=none]
% Nodo activo en el paso actual del algoritmo
\node[vertex, fill=tikz-highlight] (u) {$\eqk{u}$};

% Arista relajada en este paso
\path[->, draw=tikz-highlight, very thick]
  (s) edge node[above, text=tikz-highlight] {$\eqk{6}$} (u);

% Nodo ya resuelto (distancia final)
\node[vertex, fill=tikz-accept!20, draw=tikz-accept] (s) {$\eqk{s}$};
\end{lstlisting}

%───────────────────────────────────────────────────────────────────────────────
\section{Diagramas conmutativos}
\label{sec:tikz-cd}
%───────────────────────────────────────────────────────────────────────────────

\begin{definicion}{tikz-cd}
\label{def:tikzcd}
El paquete \texttt{tikz-cd}~\cite{tikzcd2021} extiende TikZ para producir
diagramas conmutativos con sintaxis matricial. Cada celda es un nodo;
las flechas se declaran con \texttt{\textbackslash arrow[dirección]\{etiqueta\}}.
Las etiquetas de morfismos son objetos matemáticos y aplican el coloreado
semántico de ecuaciones.
\end{definicion}

\subsection{Diagrama conmutativo: teorema del isomorfismo}
\label{subsec:cd-isomorfismo}

\begin{figure}[htbp]
  \centering
  \begin{tikzcd}[
    every arrow/.append style = {draw=tikz-edge, thick},
    every label/.append style = {text=maintext},
  ]
    \eqdom{G}
      \arrow[r, "\eqop{\varphi}"]
      \arrow[d, "\eqop{\pi}"', two heads]
    & \eqdom{H} \\
    \eqdom{G} \eqop{/} \eqdom{\ker\varphi}
      \arrow[ur, dashed, "\eqop{\tilde{\varphi}}"', "\sim"']
  \end{tikzcd}
  \caption{Primer teorema de isomorfismo: $\eqdom{G}/\eqdom{\ker\varphi}
           \cong \eqop{\mathrm{im}}(\eqop{\varphi})$.
           $\eqop{\pi}$ es la proyección canónica (flecha doble);
           $\eqop{\tilde{\varphi}}$ es el isomorfismo inducido (punteado).}
  \label{fig:cd-isomorfismo}
\end{figure}

\subsection{Diagrama conmutativo: sucesión exacta corta}
\label{subsec:cd-ses}

\begin{figure}[htbp]
  \centering
  \begin{tikzcd}[
    every arrow/.append style = {draw=tikz-edge, thick},
  ]
    \eqdom{0}
      \arrow[r]
    & \eqdom{A}
      \arrow[r, "\eqop{f}"]
      \arrow[d, "\eqop{\alpha}"']
    & \eqdom{B}
      \arrow[r, "\eqop{g}"]
      \arrow[d, "\eqop{\beta}"']
    & \eqdom{C}
      \arrow[r]
      \arrow[d, "\eqop{\gamma}"]
    & \eqdom{0} \\
    \eqdom{0}
      \arrow[r]
    & \eqdom{A'}
      \arrow[r, "\eqop{f'}"']
    & \eqdom{B'}
      \arrow[r, "\eqop{g'}"']
    & \eqdom{C'}
      \arrow[r]
    & \eqdom{0}
  \end{tikzcd}
  \caption{Diagrama del lema de los cinco: morfismo entre dos sucesiones
           exactas cortas. La conmutatividad implica que si $\eqop{\alpha}$
           y $\eqop{\gamma}$ son isomorfismos, también lo es
           $\eqop{\beta}$~\cite{mac2013categories}.}
  \label{fig:cd-five-lemma}
\end{figure}

\begin{observacion}{Estilos de flecha en tikz-cd}
Las flechas más comunes en diagramas conmutativos y su declaración:
\begin{lstlisting}[style=latex, frame=none, numbers=none]
\arrow[r]                  % flecha estandar ->
\arrow[r, two heads]       % epimorfico ->>
\arrow[r, hook]            % monomorfico (gancho)
\arrow[r, dashed]          % morfismo inducido
\arrow[r, Rightarrow]      % transformacion natural =>
\arrow[r, no head]         % relacion de equivalencia
\end{lstlisting}
\end{observacion}

%───────────────────────────────────────────────────────────────────────────────
\section{Patrones de código reutilizables}
\label{sec:tikz-patrones}
%───────────────────────────────────────────────────────────────────────────────

El patrón de definición de estilos al inicio del entorno
\texttt{tikzpicture} es el más limpio y portable: evita contaminar el
ámbito global y hace el diagrama autocontenido.

\begin{lstlisting}[style=latex, frame=none, numbers=none]
\begin{tikzpicture}[
  % Definir estilos locales que consumen aliases semanticos
  vertex/.style = {
    circle,
    draw    = tikz-node-border,
    fill    = tikz-node-bg,
    text    = maintext,
    thick,
    minimum size = 1cm,
  },
  edge/.style = {
    draw    = tikz-edge,
    ->,
    >=latex,
    thick,
  },
  % Variantes semanticas por tipo de nodo
  initial/.style  = {vertex, draw=tikz-initial},
  accepted/.style = {vertex, draw=tikz-accept},
  active/.style   = {vertex, fill=tikz-highlight},
  error/.style    = {vertex, fill=fb-error, text=background},
]
  % ... contenido del diagrama
\end{tikzpicture}
\end{lstlisting}

\begin{observacion}{Reutilización entre figuras}
Si varios diagramas de un capítulo comparten el mismo vocabulario visual,
los estilos se pueden declarar una sola vez con
\texttt{\textbackslash tikzset\{...\}} al inicio del capítulo.
Para estilos que aplican a todo el documento (como el estilo base de
autómatas), declararlos en \texttt{config/} o en un bloque de
\texttt{main.tex} después de cargar los paquetes.
\end{observacion}

%───────────────────────────────────────────────────────────────────────────────
\printbibliography[heading=subbibliography]


\begin{appendices}
  %───────────────────────────────────────────────────────────────────────────────
\chapter{Style Guide de Catppuccin Académico}
\label{ap:styleguide}
%───────────────────────────────────────────────────────────────────────────────

Este apéndice es el documento normativo del sistema de color del template.
Define la arquitectura, las reglas de decisión, el inventario completo de
aliases y el protocolo de extensión. Todo color nuevo que se añada al
template debe pasar por este documento primero.

%───────────────────────────────────────────────────────────────────────────────
\section{Arquitectura de tres capas}
\label{ap:sec:arquitectura}
%───────────────────────────────────────────────────────────────────────────────

\begin{definicion}{Las tres capas del sistema de color}
\label{ap:def:tres-capas}
El sistema de color opera en tres niveles de abstracción con una regla
de visibilidad estricta:

\textbf{Nivel 0 — Paleta cruda (\texttt{Ctp*}):} Valores hexadecimales
absolutos del flavor activo. Existen 26 colores: \texttt{CtpRosewater},
\texttt{CtpFlamingo}, \texttt{CtpPink}, \texttt{CtpMauve}, \texttt{CtpRed},
\texttt{CtpMaroon}, \texttt{CtpPeach}, \texttt{CtpYellow}, \texttt{CtpGreen},
\texttt{CtpTeal}, \texttt{CtpSky}, \texttt{CtpSapphire}, \texttt{CtpBlue},
\texttt{CtpLavender}, \texttt{CtpText}, \texttt{CtpSubtext1},
\texttt{CtpSubtext0}, \texttt{CtpOverlay2}, \texttt{CtpOverlay1},
\texttt{CtpOverlay0}, \texttt{CtpSurface2}, \texttt{CtpSurface1},
\texttt{CtpSurface0}, \texttt{CtpBase}, \texttt{CtpMantle},
\texttt{CtpCrust}.

\textbf{Nivel 1 — Aliases estructurales:} Roles de UI genéricos sin
semántica de dominio. Ejemplos: \texttt{background}, \texttt{surface0},
\texttt{overlay0}, \texttt{maintext}. Describen propiedades visuales de
elementos de interfaz.

\textbf{Nivel 2 — Aliases semánticos:} Roles de dominio del documento.
Ejemplos: \texttt{env-def}, \texttt{doc-chapter}, \texttt{code-keyword}.
Describen la intención del elemento en el documento.
\end{definicion}

La regla de visibilidad es \textbf{absoluta}:

\begin{estrella}{Invariante de capas}
\label{ap:star:invariante}
\texttt{Ctp*} solo aparece en \texttt{themes/catppuccin-palette.tex}.\\
Los archivos de \texttt{config/} consumen aliases semánticos.\\
Los archivos de contenido (\texttt{chapters/}, \texttt{environments/})
consumen comandos del template.

Cuando este invariante se mantiene, cambiar el flavor del documento
entero es descomentar una línea en \texttt{catppuccin-palette.tex}.
Cuando se viola, el cambio de flavor requiere buscar y editar en
múltiples archivos.
\end{estrella}

%───────────────────────────────────────────────────────────────────────────────
\section{Árbol de decisión para asignar colores}
\label{ap:sec:decision}
%───────────────────────────────────────────────────────────────────────────────

Cuando se necesita un color nuevo en cualquier archivo de config o de
contenido, la secuencia de preguntas es:

\begin{enumerate}[label=\textbf{\arabic*.}, leftmargin=2em]
  \item \textbf{¿Estoy en \texttt{catppuccin-palette.tex}?}
    Si sí, puedo usar \texttt{Ctp*} libremente para definir aliases.
    Si no, nunca uso \texttt{Ctp*}: avanzo al paso 2.

  \item \textbf{¿Ya existe un alias semántico para este rol?}
    Buscar en la paleta con el prefijo del dominio (\texttt{env-},
    \texttt{code-}, \texttt{tikz-}, etc.). Si existe, lo consumo.
    Si no existe, avanzo al paso 3.

  \item \textbf{¿El color tiene semántica de dominio o es infraestructura?}
    Si es semántica de dominio (``el borde de un estado aceptador de
    autómata'', ``el color de una keyword''), añado un alias semántico
    nuevo en la paleta bajo el prefijo correspondiente.
    Si es infraestructura genérica (``un fondo de superficie'', ``texto
    secundario''), uso un alias estructural existente (\texttt{surface0},
    \texttt{subtext0}).

  \item \textbf{¿Es un cambio independiente de otros elementos?}
    Si el color de este elemento siempre cambia junto con otro, comparten
    alias. Si pueden cambiar independientemente, tienen aliases separados.
\end{enumerate}

\begin{observacion}{La regla de granularidad mínima}
Un elemento nuevo merece su propio alias semántico si y solo si su
cambio de color es \textit{independiente} de otros elementos. Si el
label de figura y el label de tabla siempre deben ser del mismo color,
comparten \texttt{caption-label}. Si pueden diferir (como es el caso en
este template), tienen \texttt{caption-fig} y \texttt{caption-tab}
separados.
\end{observacion}

%───────────────────────────────────────────────────────────────────────────────
\section{Convención de nombres}
\label{ap:sec:nombres}
%───────────────────────────────────────────────────────────────────────────────

\begin{definicion}{Convención kebab-case}
\label{ap:def:kebab}
Todos los aliases siguen la convención:
\[
  \texttt{\{dominio\}-\{elemento\}[-\{variante\}]}
\]
en kebab-case. El dominio identifica la categoría; el elemento describe
el rol dentro del dominio; la variante (\texttt{frame}, \texttt{back},
\texttt{text}) solo aparece cuando el elemento tiene múltiples roles de
color. Los aliases estructurales (Nivel 1) no llevan prefijo de dominio.
\end{definicion}

\begin{acadtable}[H]{ap:tab:dominios}{Dominios activos y sus prefijos}
  \begin{tabularx}{\textwidth}{l l X}
    \toprule
    \headerrow
    \thead{Prefijo} & \thead{Dominio} & \thead{Cubre} \\
    \midrule
    (sin prefijo)    & Infraestructura     & Fondos, superficies, texto base \\
    \texttt{doc-}    & Documento           & Jerarquía de títulos \\
    \texttt{link-}   & Hipervínculos       & Referencias, URLs, citas \\
    \texttt{env-}    & Entornos            & Definición, teorema, ejemplo, etc. \\
    \texttt{eq-}     & Ecuaciones          & Coloreado semántico de matemáticas \\
    \texttt{code-}   & Código fuente       & Listings, syntax highlighting \\
    \texttt{alg-}    & Algoritmos          & algorithm2e, pseudocódigo \\
    \texttt{plot-}   & Gráficas            & pgfplots, series, ejes \\
    \texttt{tikz-}   & Diagramas TikZ      & Nodos, aristas, autómatas \\
    \texttt{table-}  & Tablas              & Encabezados, filas, reglas \\
    \texttt{caption-}& Captions            & Labels por tipo de flotante \\
    \texttt{fb-}     & Feedback            & Éxito, error, advertencia, info \\
    \texttt{diff-}   & Diffs               & Cambios insertados/eliminados \\
    \texttt{rainbow-}& Secuencia           & N colores distinguibles ordenados \\
    \texttt{frontmatter-} & Portada        & Elementos de la página de título \\
    \bottomrule
  \end{tabularx}
\end{acadtable}

%───────────────────────────────────────────────────────────────────────────────
\section{Inventario completo de aliases}
\label{ap:sec:inventario}
%───────────────────────────────────────────────────────────────────────────────

\subsection{Infraestructura visual (Nivel 1)}
\label{ap:subsec:infraestructura}

\begin{acadtable}[H]{ap:tab:infra}{Aliases estructurales: fondos, superficies y tipografía}
  \begin{tabularx}{\textwidth}{l l X}
    \toprule
    \headerrow
    \thead{Alias} & \thead{Color Ctp} & \thead{Fundamento Style Guide} \\
    \midrule
    \texttt{background}   & Base      & Background Pane \\
    \texttt{secondarybg}  & Mantle    & Secondary Panes \\
    \texttt{crustbg}      & Crust     & Secondary Panes (más oscuro) \\
    \texttt{surface0/1/2} & Surface 0/1/2 & Surface Elements \\
    \texttt{overlay0/1/2} & Overlay 0/1/2 & Overlays \\
    \texttt{maintext}     & Text      & Body Copy \\
    \texttt{subtext0/1}   & Subtext 0/1 & Sub-Headlines, Labels \\
    \texttt{subtletext}   & Overlay 1 & Subtle — texto terciario \\
    \bottomrule
  \end{tabularx}
\end{acadtable}

\subsection{Estructura del documento (\texttt{doc-})}

\begin{acadtable}[H]{ap:tab:doc}{Aliases de jerarquía de títulos}
  \begin{tabularx}{\textwidth}{l l X}
    \toprule
    \headerrow
    \thead{Alias} & \thead{Color Ctp} & \thead{Razón} \\
    \midrule
    \texttt{doc-chapter}       & Pink      & Máximo contraste; warm; ancla el capítulo \\
    \texttt{doc-section}       & Sky       & Cool; segundo nivel; contrasta con Pink \\
    \texttt{doc-subsection}    & Lavender  & Suave; tercer nivel \\
    \texttt{doc-subsubsection} & Peach     & Warm suave; cuarto nivel \\
    \texttt{doc-paragraph}     & Subtext0  & Casi texto; diferencia mínima \\
    \bottomrule
  \end{tabularx}
\end{acadtable}

\subsection{Hipervínculos (\texttt{link-})}

Mapeo directo del Catppuccin Style Guide: Links=Blue, Followed=Lavender, Hover=Sky.

\begin{acadtable}[H]{ap:tab:link}{Aliases de hipervínculos}
  \begin{tabularx}{\textwidth}{l l X}
    \toprule
    \headerrow
    \thead{Alias} & \thead{Color Ctp} & \thead{Uso en hyperref} \\
    \midrule
    \texttt{link-internal} & Blue     & \texttt{linkcolor}: refs internas, TOC \\
    \texttt{link-external} & Peach    & \texttt{urlcolor}: URLs externas \\
    \texttt{link-cite}     & Green    & \texttt{citecolor}: citas bibliográficas \\
    \texttt{link-visited}  & Lavender & Followed Links \\
    \texttt{link-hover}    & Sky      & On Hover Links \\
    \bottomrule
  \end{tabularx}
\end{acadtable}

\subsection{Entornos teoremáticos (\texttt{env-})}

Cada entorno tiene \texttt{env-X} (frame) y \texttt{env-X-back} (fondo,
generado por \texttt{\textbackslash createback\{env-X\}}).

\begin{acadtable}[H]{ap:tab:env}{Aliases de entornos teoremáticos}
  \begin{tabularx}{\textwidth}{l l X}
    \toprule
    \headerrow
    \thead{Alias} & \thead{Color Ctp} & \thead{Entorno y razón} \\
    \midrule
    \texttt{env-def}  & Blue     & Definición: epistémico, fundacional \\
    \texttt{env-thm}  & Mauve    & Teorema: el más prominente del sistema \\
    \texttt{env-prop} & Green    & Proposición: verdad demostrada \\
    \texttt{env-ex}   & Peach    & Ejemplo: warm, invita a leer \\
    \texttt{env-obs}  & Overlay0 & Observación: secundaria, no compite \\
    \texttt{env-phys} & Sapphire & Conexión física: evoca ciencia \\
    \texttt{env-star} & Yellow   & Estrella Polar: máxima prominencia \\
    \texttt{env-comp} & Pink     & Ejercicio Computacional: activo \\
    \bottomrule
  \end{tabularx}
\end{acadtable}

\subsection{Matemáticas semánticas (\texttt{eq-})}

Tres ejes ortogonales. Ver \cref{ch:ecuaciones} para la documentación completa.

\begin{acadtable}[H]{ap:tab:eq}{Aliases de coloreado semántico de ecuaciones}
  \begin{tabularx}{\textwidth}{l l l X}
    \toprule
    \headerrow
    \thead{Alias} & \thead{Macro} & \thead{Color Ctp} & \thead{Eje y rol} \\
    \midrule
    \texttt{eq-op}      & \texttt{\textbackslash eqop}  & Blue      & Sintáctico: operadores \\
    \texttt{eq-fn}      & \texttt{\textbackslash eqfn}  & Green     & Sintáctico: funciones/mapas \\
    \texttt{eq-dm}      & \texttt{\textbackslash eqdm}  & Teal      & Sintáctico: diferenciales \\
    \texttt{eq-dom}     & \texttt{\textbackslash eqdom} & Sapphire  & Sintáctico: dominios/conjuntos \\
    \midrule
    \texttt{eq-unknown} & \texttt{\textbackslash equ}   & Peach     & Epistémico: incógnita \\
    \texttt{eq-known}   & \texttt{\textbackslash eqk}   & Mauve     & Epistémico: dato/fijado \\
    \texttt{eq-const}   & \texttt{\textbackslash eqcon} & Lavender  & Epistémico: constante universal \\
    \midrule
    \texttt{eq-vec}     & \texttt{\textbackslash eqvec} & Red       & Tipo: vector/campo \\
    \texttt{eq-ten}     & \texttt{\textbackslash eqten} & Rosewater & Tipo: tensor/matriz \\
    \bottomrule
  \end{tabularx}
\end{acadtable}

\subsection{Código fuente (\texttt{code-})}

Mapeo directo del Catppuccin Style Guide: Language Defaults (Code Editors).

\begin{acadtable}[H]{ap:tab:code}{Aliases de syntax highlighting}
  \begin{tabularx}{\textwidth}{l l X}
    \toprule
    \headerrow
    \thead{Alias} & \thead{Color Ctp} & \thead{Clase sintáctica} \\
    \midrule
    \texttt{code-keyword} & Mauve    & Keywords del lenguaje \\
    \texttt{code-string}  & Green    & Strings y literales \\
    \texttt{code-comment} & Overlay2 & Comentarios \\
    \texttt{code-operator}& Sky      & Operadores \\
    \texttt{code-type}    & Yellow   & Tipos / clases / interfaces \\
    \texttt{code-fn}      & Blue     & Funciones / métodos \\
    \texttt{code-mod}     & Teal     & Módulos / namespaces / enums \\
    \texttt{code-macro}   & Rosewater& Macros / decoradores / pragmas \\
    \texttt{code-builtin} & Red      & Builtins / atoms / symbols \\
    \texttt{code-const}   & Peach    & Constantes / números \\
    \texttt{code-param}   & Maroon   & Parámetros \\
    \texttt{code-delim}   & Overlay2 & Delimitadores (\texttt{\{\}[]()}) \\
    \texttt{code-bg}      & Base     & Fondo del bloque \\
    \texttt{code-border}  & Surface0 & Frame / regla \\
    \texttt{code-linenum} & Overlay1 & Números de línea \\
    \bottomrule
  \end{tabularx}
\end{acadtable}

\subsection{Algoritmos, plots, TikZ, tablas, captions y feedback}

Las tablas de aliases para los dominios \texttt{alg-}, \texttt{plot-},
\texttt{tikz-}, \texttt{table-}, \texttt{caption-} y \texttt{fb-}
se presentan en las secciones correspondientes de los capítulos:
\cref{ch:codigo} (algoritmos), \cref{ch:figuras} (plots y captions),
\cref{ch:diagramas} (TikZ), \cref{ch:tablas} (tablas).

%───────────────────────────────────────────────────────────────────────────────
\section{Protocolo de extensión}
\label{ap:sec:extension}
%───────────────────────────────────────────────────────────────────────────────

\begin{definicion}{Protocolo para añadir un dominio nuevo}
\label{ap:def:extension}
Cuando el template necesita soporte para un paquete o tipo de elemento
nuevo, el protocolo es exactamente este — no más, no menos:
\begin{enumerate}[label=\textbf{\arabic*.}, leftmargin=2em]
  \item Añadir una sección comentada en \texttt{themes/catppuccin-palette.tex}
    con los \texttt{\textbackslash colorlet} del nuevo dominio.
  \item Si los entornos necesitan fondos, llamar
    \texttt{\textbackslash createback\{alias\}} para cada uno.
  \item Crear \texttt{config/\{dominio\}.tex} consumiendo únicamente
    los aliases recién definidos. Nunca \texttt{Ctp*}.
  \item Registrar \texttt{\textbackslash input\{config/\{dominio\}.tex\}}
    en \texttt{main.tex} en el paso 8 (configs de dominio).
\end{enumerate}
\end{definicion}

El ejemplo canónico para un dominio \texttt{circuitikz} hipotético:

\begin{lstlisting}[style=latex, frame=none, numbers=none]
% En themes/catppuccin-palette.tex:
% - Circuitos (circuitikz) 
% Roles: pasivos, activos, fuentes, cables, labels
\colorlet{circ-passive}{CtpBlue}    % resistencias, capacitores
\colorlet{circ-active}{CtpGreen}    % transistores, diodos
\colorlet{circ-source}{CtpYellow}   % fuentes de voltaje/corriente
\colorlet{circ-wire}{CtpOverlay1}   % cables y conexiones
\colorlet{circ-label}{CtpSubtext0}  % labels de componentes

% En config/circuitikz.tex:
% Nunca Ctp* aqui
\ctikzset{
  resistors/scale=0.8,
  color=circ-wire,
}
\tikzset{
  component/.style={color=circ-passive, fill=circ-passive!10},
  source/.style   ={color=circ-source},
}

% En main.tex (paso 8):
\input{config/circuitikz.tex}
\end{lstlisting}

%───────────────────────────────────────────────────────────────────────────────
\section{El mecanismo \textbackslash createback}
\label{ap:sec:createback}
%───────────────────────────────────────────────────────────────────────────────

\begin{definicion}{\textbackslash createback}
\label{ap:def:createback}
El macro \texttt{\textbackslash createback\{alias\}} genera
\texttt{\{alias\}-back} como variante de fondo clara para
\texttt{tcolorbox}. La luminosidad se adapta automáticamente al flavor:
\begin{itemize}[leftmargin=2em]
  \item Tema claro (Latte):
    \texttt{\{alias\}-back} = \texttt{white!97!\{alias\}}
  \item Tema oscuro (Frappé, Macchiato, Mocha):
    \texttt{\{alias\}-back} = \texttt{CtpBase!95!\{alias\}}
\end{itemize}
\end{definicion}

Esto garantiza que todos los entornos teoremáticos tengan fondos
consistentes sin importar el flavor. El toggle \texttt{\textbackslash ifdarktheme}
se establece manualmente en \texttt{catppuccin-palette.tex} junto a la
línea de \texttt{\textbackslash usepackage[Flavor]}.

%───────────────────────────────────────────────────────────────────────────────
\section{Fundamento en el Style Guide oficial de Catppuccin}
\label{ap:sec:fundamento}
%───────────────────────────────────────────────────────────────────────────────

Las asignaciones de este template no son arbitrarias. Cada color de dominio
tiene una correspondencia directa con las categorías del Style Guide oficial
de Catppuccin. La siguiente tabla traza esa correspondencia para los dominios
de mayor visibilidad:

\begin{acadtable}[H]{ap:tab:sg-mapping}{Correspondencia con el Style Guide oficial}
  \begin{tabularx}{\textwidth}{l l X}
    \toprule
    \headerrow
    \thead{Alias del template} & \thead{Categoría Style Guide} & \thead{Descripción SG} \\
    \midrule
    \texttt{link-internal}   & Links, URLs      & ``Blue for links'' \\
    \texttt{link-visited}    & Followed Links   & ``Lavender for followed links'' \\
    \texttt{link-hover}      & On Hover Links   & ``Sky for hover'' \\
    \texttt{code-keyword}    & Keyword          & ``Mauve for keywords'' \\
    \texttt{code-string}     & Strings          & ``Green for strings'' \\
    \texttt{code-comment}    & Comments         & ``Overlay 2 for comments'' \\
    \texttt{code-operator}   & Operators        & ``Sky for operators'' \\
    \texttt{code-type}       & Classes, Types   & ``Yellow for types'' \\
    \texttt{code-fn}         & Methods          & ``Blue for methods'' \\
    \texttt{fb-success}      & Success          & ``Green for success'' \\
    \texttt{fb-warning}      & Warnings         & ``Yellow for warnings'' \\
    \texttt{fb-error}        & Errors           & ``Red for errors'' \\
    \texttt{fb-info}         & Information      & ``Teal for information'' \\
    \texttt{env-thm}         & Highest priority & ``Mauve for most prominent'' \\
    \texttt{rainbow-1..6}    & Rainbow Highlights & Secuencia de colores ordenada \\
    \bottomrule
  \end{tabularx}
\end{acadtable}

Las extensiones propias del template (dominios \texttt{eq-}, \texttt{doc-},
\texttt{tikz-}, \texttt{table-}, \texttt{caption-}, \texttt{plot-}) no
tienen correspondencia directa en el Style Guide porque ese no cubre
documentos matemáticos. Las asignaciones se derivaron de los principios
del Style Guide: colores fríos para elementos epistémicos neutros (Blue
para operadores, Sapphire para dominios), colores cálidos para elementos
con carga semántica fuerte (Peach para incógnitas, Red para vectores).

\begin{notasbib}
El Style Guide oficial de Catppuccin está disponible en
\texttt{https://github.com/catppuccin/catppuccin/blob/main/docs/style-guide.md}.
La paleta completa con valores hexadecimales para cada flavor se documenta
en \texttt{https://catppuccin.com/palette}. El paquete
\texttt{catppuccinpalette} para \LaTeX\ se mantiene en CTAN y en
\texttt{https://github.com/catppuccin/latex}.
\end{notasbib}

  %═══════════════════════════════════════════════════════════════════════════════
\chapter{Galería de código fuente}
\label{ap:codigo-galeria}
%═══════════════════════════════════════════════════════════════════════════════

Este apéndice presenta ejemplos de código en los diez lenguajes configurados
en la plantilla. Cada estilo mapea las categorías sintácticas del lenguaje a
la paleta Catppuccin según la tabla del \cref{sec:codigo-breve}. Los estilos
se invocan con \verb|style=python|, \verb|style=rust|, etc.

\dificultad{1}{1}

%───────────────────────────────────────────────────────────────────────────────
\section{Python}
%───────────────────────────────────────────────────────────────────────────────

\begin{lstlisting}[style=python, caption={Descenso de gradiente con momentum de Nesterov}]
import numpy as np
from typing import Callable, Tuple

class SGDMomentum:
    """SGD optimizer with Nesterov momentum."""

    def __init__(self, lr: float = 1e-3, momentum: float = 0.9):
        self.lr = lr
        self.momentum = momentum
        self.velocity = None

    def step(self, params: np.ndarray, grad: np.ndarray) -> np.ndarray:
        if self.velocity is None:
            self.velocity = np.zeros_like(params)
        # Nesterov momentum update
        self.velocity = self.momentum * self.velocity - self.lr * grad
        return params + self.momentum * self.velocity - self.lr * grad

def train(model: Callable, X: np.ndarray, y: np.ndarray,
          epochs: int = 100) -> list:
    optimizer = SGDMomentum(lr=0.01)
    losses = []
    for epoch in range(epochs):
        y_hat = model(X)
        loss = np.mean((y_hat - y) ** 2)
        losses.append(loss)
        if loss < 1e-6:
            print(f"Converged at epoch {epoch}")
            break
    return losses
\end{lstlisting}

%───────────────────────────────────────────────────────────────────────────────
\section{C}
%───────────────────────────────────────────────────────────────────────────────

\begin{lstlisting}[style=c, caption={Producto matriz-vector}]
#include <stdio.h>
#include <stdlib.h>
#include <string.h>
#include <math.h>

void matvec(const double *A, const double *x, double *y,
            size_t m, size_t n) {
    memset(y, 0, m * sizeof(double));
    for (size_t i = 0; i < m; i++)
        for (size_t j = 0; j < n; j++)
            y[i] += A[i * n + j] * x[j];
}

double norm2(const double *v, size_t n) {
    double s = 0.0;
    for (size_t i = 0; i < n; i++)
        s += v[i] * v[i];
    return sqrt(s);
}

int main(void) {
    double A[] = {1.0, 2.0, 3.0, 4.0, 5.0, 6.0};
    double x[] = {1.0, 1.0, 1.0};
    double y[2];
    matvec(A, x, y, 2, 3);
    printf("||y|| = %.6f\n", norm2(y, 2));
    return EXIT_SUCCESS;
}
\end{lstlisting}

%───────────────────────────────────────────────────────────────────────────────
\section{C++}
%───────────────────────────────────────────────────────────────────────────────

\begin{lstlisting}[style=cpp, caption={Integrador de Monte Carlo con templates}]
#include <iostream>
#include <vector>
#include <random>
#include <functional>

template<typename T>
class MonteCarloIntegrator {
private:
    std::mt19937 rng_;
    std::uniform_real_distribution<T> dist_;
public:
    explicit MonteCarloIntegrator(unsigned seed = 42)
        : rng_(seed), dist_(0.0, 1.0) {}

    T integrate(std::function<T(T)> f, T a, T b, size_t N) {
        T sum = 0.0;
        for (size_t i = 0; i < N; ++i) {
            T x = a + (b - a) * dist_(rng_);
            sum += f(x);
        }
        return (b - a) * sum / static_cast<T>(N);
    }
};

int main() {
    MonteCarloIntegrator<double> mc;
    auto result = mc.integrate(
        [](double x) { return std::sin(x); },
        0.0, M_PI, 1'000'000
    );
    std::cout << "Result: " << result << std::endl;
    return 0;
}
\end{lstlisting}

%───────────────────────────────────────────────────────────────────────────────
\section{Rust}
%───────────────────────────────────────────────────────────────────────────────

\begin{lstlisting}[style=rust, caption={Solucionador RK4 con traits}]
trait OdeSolver {
    fn step(&self, f: &dyn Fn(f64, &[f64]) -> Vec<f64>,
            t: f64, y: &[f64], dt: f64) -> Vec<f64>;
}

struct RungeKutta4;

impl OdeSolver for RungeKutta4 {
    fn step(&self, f: &dyn Fn(f64, &[f64]) -> Vec<f64>,
            t: f64, y: &[f64], dt: f64) -> Vec<f64> {
        let n = y.len();
        let k1 = f(t, y);
        let y2: Vec<f64> = (0..n).map(|i| y[i] + 0.5*dt*k1[i]).collect();
        let k2 = f(t + 0.5*dt, &y2);
        let y3: Vec<f64> = (0..n).map(|i| y[i] + 0.5*dt*k2[i]).collect();
        let k3 = f(t + 0.5*dt, &y3);
        let y4: Vec<f64> = (0..n).map(|i| y[i] + dt*k3[i]).collect();
        let k4 = f(t + dt, &y4);
        (0..n).map(|i| {
            y[i] + dt/6.0 * (k1[i] + 2.0*k2[i] + 2.0*k3[i] + k4[i])
        }).collect()
    }
}

fn main() {
    let rk4 = RungeKutta4;
    let harmonic = |_t: f64, y: &[f64]| -> Vec<f64> {
        vec![y[1], -y[0]]
    };
    let mut y = vec![1.0, 0.0];
    for i in 0..1000 {
        y = rk4.step(&harmonic, i as f64 * 0.01, &y, 0.01);
    }
    println!("y(10) = {:.6}, y'(10) = {:.6}", y[0], y[1]);
}
\end{lstlisting}

%───────────────────────────────────────────────────────────────────────────────
\section{Go}
%───────────────────────────────────────────────────────────────────────────────

\begin{lstlisting}[style=go, caption={Problema de Basilea con goroutines}]
package main

import (
    "fmt"
    "math"
    "sync"
)

type Result struct {
    Index int
    Value float64
}

func partialSum(start, end int, ch chan<- Result, wg *sync.WaitGroup) {
    defer wg.Done()
    sum := 0.0
    for k := start; k < end; k++ {
        sum += 1.0 / math.Pow(float64(k), 2)
    }
    ch <- Result{Index: start, Value: sum}
}

func main() {
    const N = 1_000_000
    const workers = 8
    chunk := N / workers

    ch := make(chan Result, workers)
    var wg sync.WaitGroup

    for i := 0; i < workers; i++ {
        wg.Add(1)
        s := i*chunk + 1
        e := s + chunk
        if i == workers-1 { e = N + 1 }
        go partialSum(s, e, ch, &wg)
    }
    go func() { wg.Wait(); close(ch) }()

    total := 0.0
    for r := range ch { total += r.Value }
    fmt.Printf("pi ~ %.10f\n", math.Sqrt(6.0*total))
}
\end{lstlisting}

%───────────────────────────────────────────────────────────────────────────────
\section{Julia}
%───────────────────────────────────────────────────────────────────────────────

\begin{lstlisting}[style=julia, caption={Euler--Maruyama con dispatch múltiple}]
using Random, LinearAlgebra, Statistics

abstract type AbstractSDE end

struct OrnsteinUhlenbeck <: AbstractSDE
    theta::Float64
    mu::Float64
    sigma::Float64
end

drift(sde::OrnsteinUhlenbeck, x, t) = sde.theta * (sde.mu - x)
diffusion(sde::OrnsteinUhlenbeck, x, t) = sde.sigma

function euler_maruyama(sde::AbstractSDE, x0::Float64,
                        dt::Float64, T::Float64; seed::Int=42)
    rng = MersenneTwister(seed)
    N = Int(T / dt)
    x = zeros(N + 1)
    x[1] = x0
    @inbounds for i in 1:N
        dW = sqrt(dt) * randn(rng)
        x[i+1] = x[i] + drift(sde, x[i], (i-1)*dt)*dt +
                         diffusion(sde, x[i], (i-1)*dt)*dW
    end
    return range(0.0, T; length=N+1), x
end

ou = OrnsteinUhlenbeck(2.0, 0.0, 1.0)
t, x = euler_maruyama(ou, 1.0, 1e-3, 10.0)
@info "Stationary variance" exact=ou.sigma^2/(2*ou.theta) numerical=var(x[end-1000:end])
\end{lstlisting}

%───────────────────────────────────────────────────────────────────────────────
\section{MATLAB}
%───────────────────────────────────────────────────────────────────────────────

\begin{lstlisting}[style=matlab, caption={Elementos finitos 1D}]
function [u, x] = fem1d(N, f_rhs, a, b)
% FEM1D  Solves -u''(x) = f(x) on [a,b], u(a)=u(b)=0
%        Linear finite elements with N elements.

h = (b - a) / N;
x = linspace(a, b, N+1)';
K = sparse(N+1, N+1);
F = zeros(N+1, 1);

for e = 1:N
    Ke = (1/h) * [1 -1; -1 1];
    xm = (x(e) + x(e+1)) / 2;
    Fe = h/2 * [f_rhs(xm); f_rhs(xm)];
    idx = [e, e+1];
    K(idx, idx) = K(idx, idx) + Ke;
    F(idx) = F(idx) + Fe;
end

interior = 2:N;
u = zeros(N+1, 1);
u(interior) = K(interior, interior) \ F(interior);

u_exact = @(x) sin(pi*x) / pi^2;
err = sqrt(h * sum((u - u_exact(x)).^2));
fprintf('L2 error: %.2e (h = %.2e)\n', err, h);
end
\end{lstlisting}

%───────────────────────────────────────────────────────────────────────────────
\section{Fortran}
%───────────────────────────────────────────────────────────────────────────────

\begin{lstlisting}[style=fortran, caption={Factorización LU con pivoteo parcial}]
module linear_algebra
  implicit none
  integer, parameter :: dp = selected_real_kind(15, 307)
contains
  subroutine lu_factor(A, n, ipiv, info)
    integer, intent(in)     :: n
    real(dp), intent(inout) :: A(n, n)
    integer, intent(out)    :: ipiv(n), info
    integer  :: i, j, k, imax
    real(dp) :: maxval, temp

    info = 0
    do k = 1, n
      maxval = abs(A(k, k));  imax = k
      do i = k + 1, n
        if (abs(A(i, k)) > maxval) then
          maxval = abs(A(i, k));  imax = i
        end if
      end do
      ipiv(k) = imax
      if (maxval < epsilon(1.0_dp)) then
        info = k;  return
      end if
      if (imax /= k) then
        do j = 1, n
          temp = A(k,j); A(k,j) = A(imax,j); A(imax,j) = temp
        end do
      end if
      do i = k + 1, n
        A(i,k) = A(i,k) / A(k,k)
        do j = k + 1, n
          A(i,j) = A(i,j) - A(i,k) * A(k,j)
        end do
      end do
    end do
  end subroutine lu_factor
end module linear_algebra
\end{lstlisting}

%───────────────────────────────────────────────────────────────────────────────
\section{Assembler (x86-64)}
%───────────────────────────────────────────────────────────────────────────────

\begin{lstlisting}[style=asm, caption={Producto punto con SSE}]
; dot_product_sse(float *a, float *b, int n) -> float
; System V AMD64: rdi=a, rsi=b, edx=n; returns in xmm0
section .text
global dot_product_sse

dot_product_sse:
    xorps   xmm0, xmm0
    test    edx, edx
    jle     .done
    mov     ecx, edx
    shr     ecx, 2
    jz      .remainder
.loop4:
    movups  xmm1, [rdi]
    movups  xmm2, [rsi]
    mulps   xmm1, xmm2
    addps   xmm0, xmm1
    add     rdi, 16
    add     rsi, 16
    dec     ecx
    jnz     .loop4
    ; horizontal reduction
    movaps  xmm1, xmm0
    shufps  xmm1, xmm1, 0x4E
    addps   xmm0, xmm1
    movaps  xmm1, xmm0
    shufps  xmm1, xmm1, 0xB1
    addss   xmm0, xmm1
.remainder:
    and     edx, 3
    jz      .done
.loop1:
    movss   xmm1, [rdi]
    mulss   xmm1, [rsi]
    addss   xmm0, xmm1
    add     rdi, 4
    add     rsi, 4
    dec     edx
    jnz     .loop1
.done:
    ret
\end{lstlisting}

%───────────────────────────────────────────────────────────────────────────────
\section{Mathematica}
%───────────────────────────────────────────────────────────────────────────────

\begin{lstlisting}[style=mathematica, caption={EDP de calor: simbólica + numérica}]
(* 1D heat equation: du/dt = k * d^2u/dx^2 *)

k = 0.01;  L = Pi;

uExact[x_, t_, nTerms_] := Module[{n},
    Sum[(2/(n*Pi)) * Sin[n*x] * Exp[-k*n^2*t],
        {n, 1, nTerms, 2}]
]

sol = NDSolve[{
    D[u[x,t], t] == k * D[u[x,t], {x,2}],
    u[0,t] == 0,  u[L,t] == 0,
    u[x,0] == If[x < L/2, 2*x/L, 2*(1 - x/L)]
}, u, {x, 0, L}, {t, 0, 10}];

error = NIntegrate[
    (u[x,5] /. sol[[1]] - uExact[x, 5, 50])^2,
    {x, 0, L}] // Sqrt;

Print["L2 error at t=5: ", ScientificForm[error, 3]]
\end{lstlisting}

%───────────────────────────────────────────────────────────────────────────────
\section*{Referencia rápida de estilos}
\addcontentsline{toc}{section}{Referencia rápida de estilos}
%───────────────────────────────────────────────────────────────────────────────

\begin{acadtable}{tab:estilos-listings}{Estilos de listings disponibles}
  \begin{tabularx}{\textwidth}{l l Y}
    \toprule
    \headerrow
    \thead{Estilo} & \thead{Lenguaje} & \thead{Uso} \\
    \midrule
    \texttt{python}      & Python         & \verb|style=python| \\
    \rowcolor{table-row-alt}
    \texttt{c}           & C              & \verb|style=c| \\
    \texttt{cpp}         & C++            & \verb|style=cpp| \\
    \rowcolor{table-row-alt}
    \texttt{rust}        & Rust           & \verb|style=rust| \\
    \texttt{go}          & Go             & \verb|style=go| \\
    \rowcolor{table-row-alt}
    \texttt{julia}       & Julia          & \verb|style=julia| \\
    \texttt{matlab}      & MATLAB/Octave  & \verb|style=matlab| \\
    \rowcolor{table-row-alt}
    \texttt{fortran}     & Fortran (2008) & \verb|style=fortran| \\
    \texttt{asm}         & x86-64 (Intel) & \verb|style=asm| \\
    \rowcolor{table-row-alt}
    \texttt{mathematica} & Wolfram        & \verb|style=mathematica| \\
    \bottomrule
  \end{tabularx}
\end{acadtable}

\begin{notasbib}
La configuración de \texttt{listings} sigue la guía de estilo oficial de
Catppuccin (\url{https://catppuccin.com/style-guide}). Los lenguajes Rust
y Go se definen desde cero con \verb|\lstdefinelanguage| porque el paquete
\texttt{listings} no los incluye nativamente.
Para resaltado más avanzado (funciones definidas por el usuario, type
inference), considerar \texttt{minted} con Pygments como alternativa.
\end{notasbib}

\begin{notasbib}
La configuración de \texttt{listings} sigue la guía de estilo oficial de
Catppuccin. Los lenguajes Rust y Go se definen desde cero con
\texttt{\textbackslash lstdefinelanguage} porque \texttt{listings}
no los incluye nativamente.
\end{notasbib}

  %───────────────────────────────────────────────────────────────────────────────
\chapter{Referencia de comandos}
\label{ap:referencia}
%───────────────────────────────────────────────────────────────────────────────

Inventario completo de comandos del template organizados por dominio.
Para cada comando se indica la firma, el archivo de origen y una descripción
concisa. Los comandos de \LaTeX\ estándar y de paquetes externos no se
incluyen — solo los definidos o configurados por este template.

%───────────────────────────────────────────────────────────────────────────────
\section{Entornos teoremáticos}
\label{ap:ref:entornos}
%───────────────────────────────────────────────────────────────────────────────

Definidos en \texttt{environments/instances.tex} vía
\texttt{environments/generator.tex}.
Todos tienen la firma \texttt{\textbackslash begin\{nombre\}\{Título\}...\textbackslash end\{nombre\}}.

\begin{acadtable}[H]{ap:tab:ref-entornos}{Entornos teoremáticos}
  \begin{tabularx}{\textwidth}{l l X}
    \toprule
    \headerrow
    \thead{Entorno} & \thead{Contador} & \thead{Uso} \\
    \midrule
    \texttt{definicion}   & por capítulo & Vocabulario formal. Color: \textcolor{env-def}{\textbf{Blue}} \\
    \texttt{teorema}      & por capítulo & Resultado principal demostrado. Color: \textcolor{env-thm}{\textbf{Mauve}} \\
    \texttt{proposicion}  & por capítulo & Resultado intermedio. Color: \textcolor{env-prop}{\textbf{Green}} \\
    \texttt{ejemplo}      & por capítulo & Instancia concreta. Color: \textcolor{env-ex}{\textbf{Peach}} \\
    \texttt{observacion}  & — & Nota técnica o matiz. Color: \textcolor{env-obs}{\textbf{Overlay0}} \\
    \texttt{fisicaconexion} & — & Interpretación física. Color: \textcolor{env-phys}{\textbf{Sapphire}} \\
    \texttt{estrella}     & por capítulo & Resultado pivote del capítulo. Color: \textcolor{env-star}{\textbf{Yellow}} \\
    \texttt{ejerciciocomp}& por capítulo & Tarea computacional. Color: \textcolor{env-comp}{\textbf{Pink}} \\
    \bottomrule
  \end{tabularx}
\end{acadtable}

\begin{observacion}{Referencia a entornos con contador}
Los entornos con contador registran \texttt{\textbackslash label} internamente
cuando se usa \texttt{\textbackslash refstepcounter}. El label debe ir
\textit{dentro} del entorno. \texttt{\textbackslash cref\{def:hilbert\}}
produce \textit{def.~1.3} usando el nombre configurado en
\texttt{config/packages.tex}.
\end{observacion}

%───────────────────────────────────────────────────────────────────────────────
\section{Entornos especiales}
\label{ap:ref:especiales}
%───────────────────────────────────────────────────────────────────────────────

Definidos en \texttt{environments/specials.tex}.

\begin{acadtable}[H]{ap:tab:ref-especiales}{Entornos y comandos especiales}
  \begin{tabularx}{\textwidth}{l X}
    \toprule
    \headerrow
    \thead{Comando / Entorno} & \thead{Descripción} \\
    \midrule
    \texttt{\textbackslash dificultad\{m\}\{f\}}
      & Indicador de dificultad. \texttt{m}: matemática (1--5),
        \texttt{f}: física (1--5). Genera una caja compacta con estrellas. \\
    \texttt{notasbib}
      & Entorno de notas bibliográficas. Genera una sección
        \textit{Notas bibliográficas} no numerada con caja sutil. \\
    \texttt{ejercicio}
      & Ejercicio matemático simple sin caja. Contador por capítulo.
        Formato: \textit{Ejercicio 3.2.} \\
    \bottomrule
  \end{tabularx}
\end{acadtable}

%───────────────────────────────────────────────────────────────────────────────
\section{Entornos de tabla}
\label{ap:ref:tablas}
%───────────────────────────────────────────────────────────────────────────────

Definidos en \texttt{config/tables.tex}.

\begin{acadtable}[H]{ap:tab:ref-tablas}{Entornos y comandos de tablas}
  \begin{tabularx}{\textwidth}{l X}
    \toprule
    \headerrow
    \thead{Comando / Entorno} & \thead{Descripción} \\
    \midrule
    \texttt{acadtable[pos]\{label\}\{caption\}}
      & Tabla flotante con caption y label integrados.
        Posición por defecto \texttt{[htbp]}; usar \texttt{[H]}
        para colocación fija. Ancho: \texttt{\textbackslash textwidth}. \\
    \texttt{boxtab\{label\}\{caption\}}
      & Tabla no-flotante para usar \textit{dentro} de tcolorbox.
        Usa \texttt{\textbackslash captionof} + \texttt{\textbackslash phantomsection}.
        Ancho: \texttt{\textbackslash linewidth}. \\
    \texttt{fittable}
      & Reescala su contenido si excede \texttt{\textbackslash linewidth}.
        Recurso de último recurso: prefiere reestructurar la tabla. \\
    \texttt{\textbackslash thead\{texto\}}
      & Encabezado de columna: bold + color \texttt{table-head-text}. \\
    \texttt{\textbackslash headerrow}
      & Aplica \texttt{\textbackslash rowcolor\{table-head\}} a la fila actual. \\
    \texttt{\textbackslash tablenote\{marca\}\{texto\}}
      & Nota al pie de tabla en \texttt{\textbackslash footnotesize}.
        Multiple llamadas se apilan. \\
    \midrule
    \texttt{L\{w\}}, \texttt{C\{w\}}, \texttt{R\{w\}}
      & Columnas de ancho fijo con wrapping y alineación izquierda,
        centro, derecha. Para \texttt{tabular}. \\
    \texttt{X}, \texttt{Y}, \texttt{Z}
      & Columnas proporcionales izquierda, centro, derecha.
        Solo para \texttt{tabularx}. \\
    \bottomrule
  \end{tabularx}
\end{acadtable}

%───────────────────────────────────────────────────────────────────────────────
\section{Entornos de figura}
\label{ap:ref:figuras}
%───────────────────────────────────────────────────────────────────────────────

Definidos en \texttt{config/tcolorbox.tex}.

\begin{acadtable}[H]{ap:tab:ref-figuras}{Entornos de figura}
  \begin{tabularx}{\textwidth}{l X}
    \toprule
    \headerrow
    \thead{Entorno} & \thead{Descripción} \\
    \midrule
    \texttt{boxfig\{label\}\{caption\}}
      & Figura no-flotante para usar \textit{dentro} de tcolorbox.
        El contenido (imagen o TikZ) va antes del caption.
        Ancho de referencia: \texttt{\textbackslash linewidth}. \\
    \midrule
    \multicolumn{2}{l}{\textit{Entornos estándar de LaTeX (no definidos por el template):}} \\
    \texttt{figure[pos]}
      & Figura flotante estándar. Usar fuera de tcolorbox.
        Ancho: \texttt{\textbackslash textwidth}. \\
    \texttt{subfigure[align]\{ancho\}}
      & Subfigura dentro de \texttt{figure}. Requiere \texttt{subcaption}. \\
    \texttt{wrapfigure\{pos\}\{ancho\}}
      & Figura embebida en texto. No usar dentro de listas o tcolorbox. \\
    \bottomrule
  \end{tabularx}
\end{acadtable}

%───────────────────────────────────────────────────────────────────────────────
\section{Macros de ecuaciones semánticas}
\label{ap:ref:ecuaciones}
%───────────────────────────────────────────────────────────────────────────────

Definidos en \texttt{config/equations.tex}.
\textbf{Regla}: siempre dentro de modo matemático (\texttt{\$...\$},
\texttt{\textbackslash[...\textbackslash]}, \texttt{equation}, etc.).

\begin{acadtable}[H]{ap:tab:ref-eq-basicos}{Macros de coloreado semántico básicos}
  \begin{tabularx}{\textwidth}{l l X}
    \toprule
    \headerrow
    \thead{Macro} & \thead{Color} & \thead{Rol semántico} \\
    \midrule
    \texttt{\textbackslash eqop\{·\}} & \textcolor{eq-op}{\textbf{Blue}}
      & Operador sintáctico: $\eqop{\nabla}$, $\eqop{\partial}$,
        $\eqop{\sum}$, $\eqop{\delta}$ \\
    \texttt{\textbackslash eqfn\{·\}} & \textcolor{eq-fn}{\textbf{Green}}
      & Función/mapa: $\eqfn{f}(\cdot)$, $\eqfn{\mathcal{L}}$,
        $\eqfn{\sigma}$ \\
    \texttt{\textbackslash eqdm\{·\}} & \textcolor{eq-dm}{\textbf{Teal}}
      & Diferencial/medida: $\eqdm{dx}$, $\eqdm{dt}$, $\eqdm{dW}$ \\
    \texttt{\textbackslash eqdom\{·\}}& \textcolor{eq-dom}{\textbf{Sapphire}}
      & Dominio/conjunto: $\eqdom{\Omega}$, $\eqdom{\mathbb{R}^n}$,
        $\eqdom{\Sigma}$ \\
    \midrule
    \texttt{\textbackslash equ\{·\}}  & \textcolor{eq-unknown}{\textbf{Peach}}
      & Incógnita/buscado: $\equ{u}$, $\equ{\theta}$, $\equ{\psi}$ \\
    \texttt{\textbackslash eqk\{·\}}  & \textcolor{eq-known}{\textbf{Mauve}}
      & Conocido/parámetro: $\eqk{\lambda}$, $\eqk{m}$, $\eqk{q_0}$ \\
    \texttt{\textbackslash eqcon\{·\}}& \textcolor{eq-const}{\textbf{Lavender}}
      & Constante universal: $\eqcon{\hbar}$, $\eqcon{k_B}$, $\eqcon{\pi}$ \\
    \midrule
    \texttt{\textbackslash eqvec\{·\}}& \textcolor{eq-vec}{\textbf{Red}}
      & Vector/campo: $\eqvec{\mathbf{E}}$, $\eqvec{\mathbf{u}}$ \\
    \texttt{\textbackslash eqten\{·\}}& \textcolor{eq-ten}{\textbf{Rosewater}}
      & Tensor/matriz: $\eqten{\sigma_{ij}}$, $\eqten{W}$ \\
    \bottomrule
  \end{tabularx}
\end{acadtable}

\begin{acadtable}[H]{ap:tab:ref-eq-compuestos}{Macros compuestos}
  \begin{tabularx}{\textwidth}{l X l}
    \toprule
    \headerrow
    \thead{Macro} & \thead{Descripción} & \thead{Produce} \\
    \midrule
    \texttt{\textbackslash eqpdv\{f\}\{t\}}
      & Derivada parcial: numerador = incógnita, denominador = variable indep.
      & $\eqpdv{f}{t}$ \\
    \texttt{\textbackslash eqgrad\{u\}}
      & Gradiente coloreado de una incógnita
      & $\eqgrad{u}$ \\
    \texttt{\textbackslash eqdvg\{v\}}
      & Divergencia (evita colisión con \texttt{\textbackslash div} de physics)
      & $\eqdvg{\mathbf{v}}$ \\
    \texttt{\textbackslash eqSDE\{X\}\{a\}\{b\}}
      & SDE completa con color: $dX = a\,dt + b\,dW$
      & $\eqSDE{X}{a}{b}$ \\
    \bottomrule
  \end{tabularx}
\end{acadtable}

\begin{acadtable}[H]{ap:tab:ref-eq-toggle}{Toggle de coloreado}
  \begin{tabularx}{\textwidth}{l X}
    \toprule
    \headerrow
    \thead{Comando} & \thead{Efecto} \\
    \midrule
    \texttt{\textbackslash eqcolorstrue}  & Activa coloreado (por defecto) \\
    \texttt{\textbackslash eqcolorsfalse} & Desactiva: todas las macros pasan
      el argumento sin modificar. Usar para impresión B/N o en grupos locales. \\
    \bottomrule
  \end{tabularx}
\end{acadtable}

%───────────────────────────────────────────────────────────────────────────────
\section{Macros de notación matemática}
\label{ap:ref:notation}
%───────────────────────────────────────────────────────────────────────────────

Definidos en \texttt{config/notation.tex}.

\begin{acadtable}[H]{ap:tab:ref-notation}{Vocabulario matemático del template}
  \begin{tabularx}{\textwidth}{l l X}
    \toprule
    \headerrow
    \thead{Macro} & \thead{Produce} & \thead{Dominio} \\
    \midrule
    \texttt{\textbackslash R}, \texttt{\textbackslash N},
    \texttt{\textbackslash Z}, \texttt{\textbackslash C},
    \texttt{\textbackslash Q}, \texttt{\textbackslash F}
      & $\R$, $\N$, $\Z$, $\C$, $\Q$, $\F$ & Conjuntos numéricos \\
    \texttt{\textbackslash T}   & $A\T$  & Transpuesta \\
    \texttt{\textbackslash HH}  & $A\HH$ & Adjunto hermítico ($\dagger$) \\
    \texttt{\textbackslash Id}  & $\Id$  & Matriz identidad \\
    \texttt{\textbackslash diag}, \texttt{\textbackslash spec}
      & $\diag(\cdot)$, $\spec(\cdot)$ & Operadores matriciales \\
    \texttt{\textbackslash argmin}, \texttt{\textbackslash argmax}
      & $\argmin$, $\argmax$ & Optimización \\
    \texttt{\textbackslash Lap}, \texttt{\textbackslash Div}
      & $\Lap$, $\Div$ & Aliases de physics \\
    \midrule
    \texttt{\textbackslash E\{·\}}  & $\E{X}$  & Valor esperado $\mathbb{E}[\,\cdot\,]$ \\
    \texttt{\textbackslash Prob\{·\}}& $\Prob{A}$ & Probabilidad $\mathbb{P}(\cdot)$ \\
    \texttt{\textbackslash Ep\{·\}} & $\Ep{f}$ & Valor esperado estilo física $\langle\cdot\rangle$ \\
    \texttt{\textbackslash KL}, \texttt{\textbackslash Var}, \texttt{\textbackslash Cov}
      & $\KL$, $\Var$, $\Cov$ & Estadísticos \\
    \texttt{\textbackslash Normal}, \texttt{\textbackslash Uniform}
      & $\Normal$, $\Uniform$ & Distribuciones \\
    \midrule
    \texttt{\textbackslash Loss}  & $\Loss$ & Función de pérdida $\mathcal{L}$ \\
    \texttt{\textbackslash Lscr}  & $\Lscr$ & Pérdida script $\mathscr{L}$ \\
    \texttt{\textbackslash D}     & $\D$    & Dataset $\mathcal{D}$ \\
    \texttt{\textbackslash Params}& $\Params$ & Espacio de parámetros $\Theta$ \\
    \texttt{\textbackslash params}& $\params$ & Parámetros $\theta$ \\
    \texttt{\textbackslash score[v]}& $\score$ & Score function $\nabla_v\log p$ \\
    \texttt{\textbackslash SDE\{X\}\{a\}\{b\}}
      & $\SDE{X}{a}{b}$ & SDE sin color (inline) \\
    \midrule
    \texttt{\textbackslash Ham}   & $\Ham$  & Hamiltoniano $\hat{H}$ \\
    \texttt{\textbackslash Hscr}  & $\Hscr$ & Espacio de Hilbert $\mathscr{H}$ \\
    \texttt{\textbackslash Fscr}  & $\Fscr$ & Transformada de Fourier $\mathscr{F}$ \\
    \texttt{\textbackslash kB}    & $\kB$   & Constante de Boltzmann $k_B$ \\
    \texttt{\textbackslash keff}  & $\keff$ & Factor efectivo $k_{\mathrm{eff}}$ \\
    \midrule
    \texttt{\textbackslash OK}, \texttt{\textbackslash nOK}
      & $\OK$, $\nOK$ & Marcas de verificación y error \\
    \bottomrule
  \end{tabularx}
\end{acadtable}

%───────────────────────────────────────────────────────────────────────────────
\section{Comandos de pgfplots}
\label{ap:ref:pgfplots}
%───────────────────────────────────────────────────────────────────────────────

Estilos y ciclos definidos en \texttt{config/pgfplots.tex}.

\begin{acadtable}[H]{ap:tab:ref-pgfplots}{Estilos y colormaps de pgfplots}
  \begin{tabularx}{\textwidth}{l X}
    \toprule
    \headerrow
    \thead{Nombre} & \thead{Descripción} \\
    \midrule
    \texttt{catppuccin} (cycle list)
      & 8 series con colores \texttt{plot-series-1..8} y marcadores
        distintos. Activado por defecto en todos los ejes. \\
    \texttt{catppuccin-clean} (estilo)
      & Grid mayor, fonts pequeñas, leyenda en esquina superior derecha.
        Para publicación. \\
    \texttt{catppuccin-filled} (estilo)
      & Como \texttt{clean} pero con fondo \texttt{secondarybg}.
        Para notas y presentaciones. \\
    \texttt{catppuccin-heat} (colormap)
      & Secuencial: magnitudes positivas. \texttt{background} → series frías
        → series cálidas. \\
    \texttt{catppuccin-diverge} (colormap)
      & Divergente: residuos con signo. \texttt{plot-series-1} ←→
        \texttt{maintext} ←→ \texttt{plot-series-5}. \\
    \bottomrule
  \end{tabularx}
\end{acadtable}

%───────────────────────────────────────────────────────────────────────────────
\section{Prefijos de etiquetas recomendados}
\label{ap:ref:labels}
%───────────────────────────────────────────────────────────────────────────────

\begin{acadtable}[H]{ap:tab:ref-labels}{Convención de prefijos para \texttt{\textbackslash label}}
  \begin{tabularx}{\textwidth}{l l l}
    \toprule
    \headerrow
    \thead{Tipo} & \thead{Prefijo} & \thead{Ejemplo} \\
    \midrule
    Capítulo       & \texttt{ch:}     & \texttt{ch:ecuaciones} \\
    Sección        & \texttt{sec:}    & \texttt{sec:tres-ejes} \\
    Subsección     & \texttt{subsec:} & \texttt{subsec:eje-epistemico} \\
    Ecuación       & \texttt{eq:}     & \texttt{eq:schrodinger} \\
    Figura         & \texttt{fig:}    & \texttt{fig:convergencia} \\
    Tabla          & \texttt{tab:}    & \texttt{tab:col-tipos} \\
    Algoritmo      & \texttt{alg:}    & \texttt{alg:cg} \\
    Listing/Código & \texttt{lst:}    & \texttt{lst:sgd} \\
    Apéndice       & \texttt{ap:}     & \texttt{ap:styleguide} \\
    Definición     & \texttt{def:}    & \texttt{def:hilbert} \\
    Teorema        & \texttt{thm:}    & \texttt{thm:lp-completo} \\
    Proposición    & \texttt{prop:}   & \texttt{prop:norma-continua} \\
    Ejemplo        & \texttt{ex:}     & \texttt{ex:elbo} \\
    Estrella       & \texttt{star:}   & \texttt{star:espectral} \\
    Ejercicio comp.& \texttt{comp:}   & \texttt{comp:triangular} \\
    \bottomrule
  \end{tabularx}
\end{acadtable}

%───────────────────────────────────────────────────────────────────────────────
\section{Nombres en español de cleveref}
\label{ap:ref:cref}
%───────────────────────────────────────────────────────────────────────────────

Configurados en \texttt{config/packages.tex}.
\texttt{\textbackslash cref} produce la forma corta; \texttt{\textbackslash Cref}
la capitalizada.

\begin{acadtable}[H]{ap:tab:ref-cref}{Nombres de cleveref en español}
  \begin{tabularx}{\textwidth}{l l l}
    \toprule
    \headerrow
    \thead{Tipo} & \thead{\textbackslash cref} & \thead{\textbackslash Cref} \\
    \midrule
    Ecuación    & ec.~(1.1)      & Ecuación~(1.1) \\
    Figura      & fig.~1.1       & Figura~1.1 \\
    Tabla       & tabla~1.1      & Tabla~1.1 \\
    Capítulo    & cap.~1         & Capítulo~1 \\
    Sección     & sec.~1.1       & Sección~1.1 \\
    Algoritmo   & alg.~1         & Algoritmo~1 \\
    Apéndice    & ap.~A          & Apéndice~A \\
    Código      & cód.~1         & Código~1 \\
    Definición  & def.~1.1       & Definición~1.1 \\
    Teorema     & teo.~1.1       & Teorema~1.1 \\
    Proposición & prop.~1.1      & Proposición~1.1 \\
    Ejemplo     & ej.~1.1        & Ejemplo~1.1 \\
    Estrella    & ~1.1          & Estrella Polar~1.1 \\
    Ej. comp.   & ej. comp.~1.1  & Ejercicio Computacional~1.1 \\
    \bottomrule
  \end{tabularx}
\end{acadtable}

%───────────────────────────────────────────────────────────────────────────────
\section{Checklist de extensión del template}
\label{ap:ref:checklist}
%───────────────────────────────────────────────────────────────────────────────

Lista de verificación al añadir cualquier componente nuevo:

\begin{enumerate}[label=$\square$, leftmargin=2em]
  \item ¿Los colores nuevos están en \texttt{catppuccin-palette.tex}
    bajo el prefijo de dominio correcto?
  \item ¿Se llamó \texttt{\textbackslash createback\{\}} para aliases
    que se usan como fondos de tcolorbox?
  \item ¿El archivo de config consume solo aliases (sin \texttt{Ctp*})?
  \item ¿Se registró el \texttt{\textbackslash input} en \texttt{main.tex}
    en el paso 8?
  \item ¿Si el dominio tiene entornos flotantes, existe la variante
    no-flotante para uso dentro de tcolorbox?
  \item ¿Las tablas con columnas \texttt{S} de siunitx envuelven los
    headers en \texttt{\textbackslash multicolumn\{1\}\{c\}\{...\}}?
  \item ¿Las macros de ecuaciones siempre van dentro de modo matemático?
  \item ¿Los \texttt{\textbackslash label} dentro de tcolorbox están
    \textit{dentro} del entorno, no en el argumento del título?
\end{enumerate}

  %───────────────────────────────────────────────────────────────────────────────
\chapter{Galería de referencia implementada}
\label{ch:galeria}
%───────────────────────────────────────────────────────────────────────────────

% INSTRUCCIÓN PARA LA IA QUE LEE ESTE CAPÍTULO:
% Este capítulo es tu referencia de implementación. Cada sección muestra
% un elemento del template implementado correctamente con coloreado semántico,
% theming correcto y citas bibliográficas reales. Úsalo como fuente de
% copy-paste y como verificación antes de generar cualquier contenido.
%
% REGLAS ABSOLUTAS que este capítulo ejemplifica:
%   1. Nunca Ctp* fuera de catppuccin-palette.tex — todos los colores
%      usan aliases semánticos (env-*, tikz-*, plot-*, eq-*, etc.)
%   2. Las macros \eqop, \equ, \eqk, etc. SIEMPRE dentro de modo matemático
%   3. Cada capítulo termina con \printbibliography[heading=subbibliography]
%   4. Los flotantes (figure, acadtable) NO van dentro de tcolorbox —
%      usar boxfig y boxtab en su lugar

Este capítulo documenta cada elemento del template con ejemplos
funcionales, coloreado correcto y referencias bibliográficas reales.
Está diseñado para ser usado directamente como referencia de implementación.

%───────────────────────────────────────────────────────────────────────────────
\section{Entornos teoremáticos}
\label{sec:gal:entornos}
%───────────────────────────────────────────────────────────────────────────────

% Los ocho entornos tienen firma uniforme:
%   \begin{nombre}{Título de esta instancia}
%   \label{tipo:etiqueta}
%   Contenido con ecuaciones coloreadas...
%   \end{nombre}
%
% NUNCA uses un entorno en lugar de otro por estética.
% Cada uno tiene un rol epistémico preciso — ver README.md.

\subsection{Nivel 1 — Fundamentos}

\begin{definicion}{Espacio de Hilbert}
\label{def:gal:hilbert}
Un \textbf{espacio de Hilbert}~\cite{reed1980}
$(\eqdom{\mathcal{H}}, \langle\cdot,\cdot\rangle)$ es un espacio vectorial
sobre $\mathbb{K}$ con producto interior
$\langle \cdot, \cdot \rangle : \eqdom{\mathcal{H}} \times \eqdom{\mathcal{H}}
\to \mathbb{K}$ tal que la norma inducida
$\|\equ{x}\| = \sqrt{\langle \equ{x}, \equ{x} \rangle}$
hace a $\eqdom{\mathcal{H}}$ completo.
\end{definicion}

\begin{teorema}{Teorema de representación de Riesz~\cite{reed1980}}
\label{thm:gal:riesz}
Sea $\eqdom{\mathcal{H}}$ un espacio de Hilbert y
$\eqfn{\varphi} : \eqdom{\mathcal{H}} \to \mathbb{K}$ un funcional lineal
continuo. Existe un único $\eqk{y} \in \eqdom{\mathcal{H}}$ tal que
\[
  \eqfn{\varphi}(\equ{x}) = \langle \equ{x}, \eqk{y} \rangle
  \qquad \forall\, \equ{x} \in \eqdom{\mathcal{H}}.
\]
Además $\|\eqfn{\varphi}\| = \|\eqk{y}\|$.
\end{teorema}

\begin{proposicion}{Identidad de polarización}
\label{prop:gal:polarizacion}
En un espacio de Hilbert real~\cite{munkres2000}, el producto interior
se recupera de la norma:
\[
  \langle \eqk{x}, \eqk{y} \rangle
  = \frac{1}{4}\left(\|\eqk{x}+\eqk{y}\|^2 - \|\eqk{x}-\eqk{y}\|^2\right).
\]
\end{proposicion}

\subsection{Nivel 2 — Apoyo}

\begin{ejemplo}{$\ell^2$ como espacio de Hilbert}
\label{ex:gal:l2}
El espacio $\eqdom{\ell^2} = \{(\equ{x}_n) \subset \mathbb{R} \mid
\eqop{\sum}_n \equ{x}_n^2 < \infty\}$ con producto interior
$\langle \equ{x}, \equ{y} \rangle = \eqop{\sum}_{n=1}^\infty \equ{x}_n \equ{y}_n$
es un espacio de Hilbert separable~\cite{reed1980}. La base ortonormal
canónica $\{\eqk{e}_n\}$ es una base de Schauder pero no algebraica (Hamel).
\end{ejemplo}

\begin{observacion}{Completitud es indispensable}
El espacio $C[0,1]$ con la norma $\eqdom{L^2}$ tiene producto interior pero
\textit{no} es de Hilbert — su completación es $\eqdom{L^2[0,1]}$.
La convergencia débil $\eqk{e}_n \rightharpoonup 0$ en $\eqdom{\ell^2}$
no implica convergencia en norma: $\|\eqk{e}_n\| = 1$ para todo $n$.
Este fenómeno~\cite{reed1980} es fuente frecuente de errores en argumentos
de compacidad en espacios de dimensión infinita.
\end{observacion}

\begin{fisicaconexion}{Estados cuánticos y espacios de Hilbert~\cite{griffiths2018}}
Los estados de un sistema cuántico son vectores en
$\eqdom{\mathcal{H}} = \eqdom{L^2(\mathbb{R}^3)}$.
La norma $\|\equ{\psi}\|^2 = \eqop{\int}_{\eqdom{\mathbb{R}^3}}
|\equ{\psi}(\eqvec{\mathbf{r}})|^2 \eqdm{d^3r} = 1$
es la condición de normalización. La evolución temporal
$\equ{\psi}(t) = e^{-\eqcon{i}\eqop{\hat{H}}t/\eqcon{\hbar}}\eqk{\psi_0}$
es un operador unitario sobre $\eqdom{\mathcal{H}}$.
\end{fisicaconexion}

\subsection{Nivel 3 — Especiales}

\begin{estrella}{Teorema espectral para operadores autoadjuntos}
\label{star:gal:espectral}
Sea $\eqop{A} : \eqdom{\mathcal{H}} \to \eqdom{\mathcal{H}}$ autoadjunto
acotado~\cite{reed1980}. Existe una medida espectral única
$\eqfn{E} : \mathcal{B}(\mathbb{R}) \to \mathcal{L}(\eqdom{\mathcal{H}})$
tal que
\[
  \eqop{A} = \eqop{\int}_{\sigma(\eqop{A})} \equ{\lambda}\,
  \eqdm{dE}(\equ{\lambda}),
\]
y $\eqfn{f}(\eqop{A}) = \eqop{\int} \eqfn{f}(\equ{\lambda})\,
\eqdm{dE}(\equ{\lambda})$ define un cálculo funcional consistente.
\end{estrella}

\begin{ejerciciocomp}{Verificación numérica del método de potencias}
\label{comp:gal:potencias}
Implementa la iteración de potencias para hallar el eigenvalor dominante
$\equ{\lambda_{\max}}$ de una matriz $\eqten{A} \in \eqdom{\mathbb{R}^{n\times n}}$
SPD aleatoria. Verifica que $|\equ{\lambda}_k - \equ{\lambda_{\max}}|$
converge con tasa $|\eqk{\lambda_2}/\eqk{\lambda_1}|^k$
comparando con los eigenvalores exactos de \texttt{np.linalg.eigh}.
Usa $\eqk{n} \in \{50, 200, 500\}$ y reporta la convergencia en una tabla.
\end{ejerciciocomp}

%───────────────────────────────────────────────────────────────────────────────
\section{Matemáticas semánticas}
\label{sec:gal:ecuaciones}
%───────────────────────────────────────────────────────────────────────────────

% REGLA CRÍTICA PARA LA IA:
% Las macros \eqop, \equ, \eqk, \eqfn, \eqdm, \eqdom, \eqvec, \eqten, \eqcon
% SOLO funcionan dentro de modo matemático: $...$, \[...\], equation, align.
% Fuera de modo matemático el color no se aplica — error frecuente de IA.

\dificultad{4}{5}

\subsection{Los tres ejes combinados: Schrödinger}

% Eje 1 (sintáctico): \eqop para operadores, \eqfn para funciones, \eqdm para medidas
% Eje 2 (epistémico): \equ para incógnitas, \eqk para datos, \eqcon para constantes
% Eje 3 (tipo): \eqvec para vectores, \eqten para tensores
\begin{equation}
\label{eq:gal:schrodinger}
\eqop{i}\eqcon{\hbar}\,\eqop{\frac{\partial}{\partial t}}\equ{\Psi}
= \eqop{\hat{H}}\equ{\Psi}
= \left(\eqop{-}\frac{\eqcon{\hbar}^2}{2\eqk{m}}\eqop{\nabla}^2
  + \eqfn{V}(\eqvec{\mathbf{r}},t)\right)\equ{\Psi}.
\end{equation}

El hamiltoniano $\eqop{\hat{H}}$ y $\eqop{\nabla}^2$ son operadores~\cite{griffiths2018}.
$\equ{\Psi}$ es la incógnita. $\eqfn{V}$ es el potencial (función conocida del
sistema). $\eqcon{\hbar}$ y $\eqk{m}$ son constante y parámetro respectivamente.

\subsection{Dualidad epistémica: física vs. ML}

% Esta es la demostración canónica del eje epistémico.
% La estructura algebraica es idéntica; solo el color cambia.
\textbf{Lectura física}~\cite{griffiths2018} —
parámetros $\eqk{\theta}$ conocidos, estado $\equ{x}$ buscado:
\begin{equation}
\label{eq:gal:physics-reading}
\equ{x^*} = \eqop{\argmin}_{\equ{x}}\;
\eqfn{\mathcal{L}}\!\left(\eqfn{f}(\equ{x};\,\eqk{\theta}),\; \eqk{y}\right).
\end{equation}

\textbf{Lectura ML}~\cite{goodfellow2016} —
dato $\eqk{x}$ conocido, parámetros $\equ{\theta}$ aprendidos:
\begin{equation}
\label{eq:gal:ml-reading}
\equ{\theta^*} = \eqop{\argmin}_{\equ{\theta}}\;
\eqfn{\mathcal{L}}\!\left(\eqfn{f}(\eqk{x};\,\equ{\theta}),\; \eqk{y}\right).
\end{equation}

\subsection{Ecuaciones de Maxwell (eje de tipo)}

\begin{align}
\eqop{\nabla \cdot} \eqvec{\mathbf{E}}
  &= \frac{\eqk{\rho}}{\eqcon{\varepsilon_0}},
  \label{eq:gal:gauss-e}\\
\eqop{\nabla \times} \eqvec{\mathbf{E}}
  &= \eqop{-}\eqop{\frac{\partial}{\partial t}}\eqvec{\mathbf{B}},
  \label{eq:gal:faraday}\\
\eqop{\nabla \cdot} \eqvec{\mathbf{B}}
  &= 0,
  \label{eq:gal:gauss-b}\\
\eqop{\nabla \times} \eqvec{\mathbf{B}}
  &= \eqcon{\mu_0}\!\left(\eqvec{\mathbf{J}}
    + \eqcon{\varepsilon_0}\eqop{\frac{\partial}{\partial t}}\eqvec{\mathbf{E}}\right).
  \label{eq:gal:ampere}
\end{align}

Los campos $\eqvec{\mathbf{E}}$, $\eqvec{\mathbf{B}}$, $\eqvec{\mathbf{J}}$
usan \texttt{\textbackslash eqvec}~\cite{jackson1999}. Las constantes
$\eqcon{\varepsilon_0}$, $\eqcon{\mu_0}$ usan \texttt{\textbackslash eqcon}.
La fuente $\eqk{\rho}$ es un dato del problema (\texttt{\textbackslash eqk}).

\subsection{Macros compuestos}

Proceso de Ornstein-Uhlenbeck~\cite{oksendal2003}:
\begin{equation}
\label{eq:gal:ou}
\eqSDE{X_t}{-\eqk{\theta}\equ{X_t}}{\eqk{\sigma}}.
\end{equation}

Fokker-Planck para la densidad $\equ{p}(\eqvec{\mathbf{x}},t)$~\cite{oksendal2003}:
\begin{equation}
\label{eq:gal:fokker-planck}
\eqpdv{p}{t}
= \eqop{-}\eqop{\nabla_x \cdot}\!\left(\eqfn{f}\,\equ{p}\right)
+ \frac{\eqk{\sigma}^2}{2}\,\eqop{\nabla_x^2}\,\equ{p}.
\end{equation}

Score matching~\cite{hyvarinen2005} (ML — $\equ{\theta}$ es la incógnita):
\begin{equation}
\label{eq:gal:score}
\eqfn{\mathcal{L}}(\equ{\theta})
= \eqcon{\tfrac{1}{2}}\,
  \eqop{\mathbb{E}}_{\eqk{x} \sim \eqdom{p_{\text{data}}}}
  \left[
    \norm{\eqfn{s}_{\equ{\theta}}(\eqk{x})
    - \eqop{\nabla}_{\eqk{x}} \log \eqk{p}(\eqk{x})}^2
  \right].
\end{equation}

Tensor de Einstein ($\eqten{g_{\mu\nu}}$ es el tensor métrico):
\begin{equation}
\label{eq:gal:einstein}
\equ{R_{\mu\nu}} - \eqcon{\tfrac{1}{2}}\eqten{g_{\mu\nu}}\equ{R}
+ \eqcon{\Lambda}\eqten{g_{\mu\nu}}
= \eqcon{\frac{8\pi G}{c^4}}\eqk{T_{\mu\nu}}.
\end{equation}

\subsection{Toggle B/N — versión impresión}

% \eqcolorsfalse desactiva todo el coloreado
% Las macros se vuelven transparentes — no hay que editar ninguna ecuación
{%
\eqcolorsfalse
\begin{equation}
\eqop{i}\eqcon{\hbar}\,\eqop{\frac{\partial}{\partial t}}\equ{\Psi}
= \left(\eqop{-}\frac{\eqcon{\hbar}^2}{2\eqk{m}}\eqop{\nabla}^2
  + \eqfn{V}\right)\equ{\Psi}.
\tag{\ref{eq:gal:schrodinger}$'$}
\end{equation}
}

%───────────────────────────────────────────────────────────────────────────────
\section{Algoritmos con coloreado semántico}
\label{sec:gal:algoritmos}
%───────────────────────────────────────────────────────────────────────────────

% Los algoritmos NO están envueltos en tcolorbox — eso rompe el paginado.
% Las expresiones matemáticas dentro del algoritmo usan coloreado semántico normal.

\begin{algorithm}[H]
  \DontPrintSemicolon
  \SetKwInOut{Input}{Entrada}\SetKwInOut{Output}{Salida}
  \Input{$\eqfn{f}:\eqdom{\mathbb{R}^n}\to\mathbb{R}$,
         $\eqop{\nabla}\eqfn{f}$,
         $\eqk{x_0}$, $\eqk{\eta}$, $\eqk{\varepsilon}$}
  \Output{Mínimo aproximado $\equ{x^*}$}
  $\equ{x} \gets \eqk{x_0}$\;
  \While{$\|\eqop{\nabla}\eqfn{f}(\equ{x})\| > \eqk{\varepsilon}$}{
    $\equ{x} \gets \equ{x} - \eqk{\eta}\,\eqop{\nabla}\eqfn{f}(\equ{x})$\;
  }
  \Return{$\equ{x}$}
  \caption{Descenso de gradiente~\cite{goodfellow2016}}
  \label{alg:gal:gradient}
\end{algorithm}

\begin{algorithm}[H]
  \DontPrintSemicolon
  \SetKwInOut{Input}{Entrada}\SetKwInOut{Output}{Salida}
  \Input{SDE $\eqdm{d}\equ{X} = \eqfn{\mu}(\equ{X},t)\eqdm{dt}
         + \eqfn{\sigma}(\equ{X},t)\eqdm{dW}$,
         $\eqk{X_0}$, $\eqk{\Delta t}$, $\eqk{T}$}
  \Output{Trayectoria $\{\equ{X}_k\}_{k=0}^{\eqk{N}}$}
  $\equ{X} \gets \eqk{X_0}$\;
  \For{$k = 0$ \KwTo $\eqk{N}-1$}{
    $t \gets k\,\eqk{\Delta t}$\;
    $\equ{\xi} \gets \mathcal{N}(0,1)$\tcp*{ruido gaussiano}
    $\equ{X} \gets \equ{X}
      + \eqfn{\mu}(\equ{X},t)\,\eqk{\Delta t}
      + \eqfn{\sigma}(\equ{X},t)\sqrt{\eqk{\Delta t}}\,\equ{\xi}$\;
  }
  \Return{trayectoria}
  \caption{Euler-Maruyama~\cite{oksendal2003}}
  \label{alg:gal:em}
\end{algorithm}

%───────────────────────────────────────────────────────────────────────────────
\section{Código fuente}
\label{sec:gal:codigo}
%───────────────────────────────────────────────────────────────────────────────

% Siempre especificar style=lenguaje.
% Para código LaTeX del template, usar style=latex.

\begin{lstlisting}[style=python,
  caption={Descenso de gradiente estocástico con momentum~\cite{goodfellow2016}},
  label={lst:gal:sgd}]
import numpy as np

def sgd_momentum(grad_fn, theta0, lr=1e-3, beta=0.9, n_iter=1000):
    """SGD con momentum de Polyak."""
    theta = theta0.copy()
    v     = np.zeros_like(theta)
    for _ in range(n_iter):
        g  = grad_fn(theta)
        v  = beta * v + (1 - beta) * g
        theta -= lr * v
    return theta
\end{lstlisting}

\begin{lstlisting}[style=julia,
  caption={Iteración de potencias en Julia~\cite{leveque2007}},
  label={lst:gal:power}]
function power_iteration(A::AbstractMatrix; tol=1e-10, maxiter=1000)
    n      = size(A, 1)
    v      = normalize(randn(n))
    lambda = zero(eltype(A))
    for _ in 1:maxiter
        w      = A * v
        lambda = dot(v, w)         # cociente de Rayleigh
        v      = normalize(w)
    end
    return lambda, v
end
\end{lstlisting}

\begin{lstlisting}[style=latex,
  caption={Patrón correcto de ecuación semántica},
  label={lst:gal:eq-pattern}]
% Siempre dentro de modo matematico
\[
  \eqop{i}\eqcon{\hbar}\,\eqop{\partial_t}\equ{\Psi}
  = \eqop{\hat{H}}\equ{\Psi}
\]

% En tabla (celda con modo matematico)
Operador  & $\eqop{\nabla}$    & Azul     \\
Incognita & $\equ{\psi}$       & Peach    \\
Dato      & $\eqk{m}$          & Mauve    \\

% INCORRECTO: fuera de modo matematico - no colorea
\eqop{\nabla}
\end{lstlisting}

%───────────────────────────────────────────────────────────────────────────────
\section{Tablas}
\label{sec:gal:tablas}
%───────────────────────────────────────────────────────────────────────────────

% acadtable: flotante en flujo normal de texto → \textwidth
% boxtab: no-flotante dentro de tcolorbox → \linewidth
% Columnas S de siunitx: encabezados SIEMPRE en \multicolumn{1}{c}{\thead{...}}
% Celdas con texto en col S: \multicolumn{1}{c}{texto}

\begin{acadtable}[H]{tab:gal:optimizadores}{Métodos de optimización de primer orden~\cite{goodfellow2016}}
  \begin{tabularx}{\textwidth}{l Y Y l}
    \toprule
    \headerrow
    \thead{Método} & \thead{Convergencia} & \thead{Memoria} & \thead{Hessiano} \\
    \midrule
    Gradiente      & $\mathcal{O}(1/k)$          & $\mathcal{O}(n)$   & No \\
    \rowcolor{table-row-alt}
    Momentum       & $\mathcal{O}(1/k^2)$         & $\mathcal{O}(n)$   & No \\
    L-BFGS         & Superlineal                   & $\mathcal{O}(mn)$  & Aprox. \\
    \rowcolor{table-row-alt}
    Adam           & $\mathcal{O}(1/\sqrt{k})$    & $\mathcal{O}(n)$   & No \\
    Newton         & Cuadrática                    & $\mathcal{O}(n^2)$ & Exacto \\
    \bottomrule
  \end{tabularx}
\end{acadtable}

\begin{acadtable}[H]{tab:gal:convergencia}{Convergencia de RK4 para $\equ{y}' = \eqop{-}\eqk{\lambda}\equ{y}$~\cite{leveque2007}}
  \begin{tabularx}{\textwidth}{Y S[table-format=1.2e-1] S[table-format=1.4e-2] S[table-format=1.2]}
    \toprule
    \headerrow
    \multicolumn{1}{Y}{\thead{$\eqk{h}$}} &
    \multicolumn{1}{c}{\thead{$\|\equ{e}\|_\infty$}} &
    \multicolumn{1}{c}{\thead{Razón}} &
    \multicolumn{1}{c}{\thead{Orden emp.}} \\
    \midrule
    $1/4$  & 2.31e-3 & \multicolumn{1}{c}{---} & \multicolumn{1}{c}{---} \\
    $1/8$  & 1.48e-4 & 0.0641 & 3.96 \\
    $1/16$ & 9.33e-6 & 0.0630 & 3.99 \\
    $1/32$ & 5.84e-7 & 0.0626 & 4.00 \\
    \bottomrule
  \end{tabularx}
  \tablenote{*}{$\eqk{\lambda}=1$, $\equ{y}(0)=1$, $t\in[0,1]$.}
\end{acadtable}

% boxtab dentro de un entorno tcolorbox:
\begin{definicion}{Media y varianza muestral}
\label{def:gal:estadisticos}
Dado un conjunto $\{\eqk{x_i}\}_{i=1}^{\eqk{n}}$ de observaciones i.i.d.:
\begin{boxtab}{tab:gal:stats}{Estimadores muestrales}
  \begin{tabularx}{\linewidth}{l X}
    \toprule
    \headerrow
    \thead{Estimador} & \thead{Fórmula} \\
    \midrule
    Media
      & $\equ{\bar{x}} = \eqop{\frac{1}{\eqk{n}}}
          \eqop{\sum}_{i=1}^{\eqk{n}} \eqk{x_i}$ \\
    Varianza
      & $\equ{s^2} = \eqop{\frac{1}{\eqk{n}-1}}
          \eqop{\sum}_{i=1}^{\eqk{n}} (\eqk{x_i} - \equ{\bar{x}})^2$ \\
    Error estándar
      & $\equ{\text{SE}} = \equ{s}/\sqrt{\eqk{n}}$ \\
    \bottomrule
  \end{tabularx}
\end{boxtab}
El estimador $\equ{s^2}$ es insesgado: $\eqop{\mathbb{E}}[\equ{s^2}] = \eqk{\sigma^2}$.
\end{definicion}

%───────────────────────────────────────────────────────────────────────────────
\section{Figuras y gráficas}
\label{sec:gal:figuras}
%───────────────────────────────────────────────────────────────────────────────

% CRÍTICO: nunca Ctp* en pgfplots — usar plot-series-N, plot-axis, etc.
% Estilos: catppuccin-clean, catppuccin-filled
% Cycle list: catppuccin (8 series, plot-series-1..8)

\begin{figure}[htbp]
  \centering
  \begin{tikzpicture}
    \begin{axis}[
      catppuccin-clean,
      width=0.85\textwidth, height=6cm,
      xlabel={Época},
      ylabel={Pérdida $\eqfn{\mathcal{L}}(\equ{\theta})$},
      title={Convergencia de optimizadores~\cite{goodfellow2016}},
      ymode=log, xmin=0, xmax=100,
      legend pos=north east, no markers,
    ]
      \addplot+[thick, domain=0:100, samples=80]
        {2*exp(-0.02*x)+0.10}; \addlegendentry{SGD}
      \addplot+[thick, domain=0:100, samples=80]
        {2*exp(-0.04*x)+0.05}; \addlegendentry{Momentum}
      \addplot+[thick, domain=0:100, samples=80]
        {2*exp(-0.07*x)+0.02}; \addlegendentry{Adam}
      \addplot+[thick, domain=0:100, samples=80]
        {2*exp(-0.05*x)+0.03}; \addlegendentry{RMSprop}
    \end{axis}
  \end{tikzpicture}
  \caption{Pérdida $\eqfn{\mathcal{L}}(\equ{\theta})$ por épocas.
           Colores del cycle list \texttt{catppuccin} (aliases
           \texttt{plot-series-1..4}, sin \texttt{Ctp*}).}
  \label{fig:gal:optimizers}
\end{figure}

\begin{figure}[htbp]
  \centering
  \begin{tikzpicture}
    \begin{loglogaxis}[
      catppuccin-clean,
      width=0.85\textwidth, height=6cm,
      xlabel={Paso de malla $\eqk{h}$},
      ylabel={Error $\|\equ{u}_h - \eqk{u}\|_2$},
      title={Convergencia FD centradas~\cite{leveque2007}},
      legend pos=north west,
      xtick={0.0625,0.125,0.25,0.5},
      xticklabels={$1/16$,$1/8$,$1/4$,$1/2$},
    ]
      \addplot[plot-series-1, mark=*] coordinates {
        (0.5,1.2e-2)(0.25,3.1e-3)(0.125,7.8e-4)(0.0625,1.9e-4)};
        \addlegendentry{FD centradas}
      \addplot[plot-series-5, dashed, no markers]
        coordinates {(0.5,1.2e-2)(0.0625,1.9e-4)};
        \addlegendentry{$\mathcal{O}(\eqk{h}^2)$}
    \end{loglogaxis}
  \end{tikzpicture}
  \caption{Convergencia $\mathcal{O}(\eqk{h}^2)$ del esquema de FD centradas
           para la ecuación de Poisson 1D.}
  \label{fig:gal:fd}
\end{figure}

\begin{figure}[htbp]
  \centering
  \begin{tikzpicture}
    \begin{axis}[
      catppuccin-clean,
      width=0.85\textwidth, height=5.5cm,
      xlabel={$\equ{x}$},
      ylabel={$\equ{\psi}_n(\equ{x})$},
      title={Eigenfunciones del oscilador armónico~\cite{griffiths2018}},
      domain=-4:4, samples=150, ymin=-1.2, ymax=1.5,
      no markers, legend pos=south east,
    ]
      \addplot+[thick] {0.7511*exp(-x^2/2)};
        \addlegendentry{$\equ{\psi_0}$}
      \addplot+[thick] {0.5311*2*x*exp(-x^2/2)};
        \addlegendentry{$\equ{\psi_1}$}
      \addplot+[thick] {0.2659*(4*x^2-2)*exp(-x^2/2)};
        \addlegendentry{$\equ{\psi_2}$}
      \addplot+[thick] {0.1083*(8*x^3-12*x)*exp(-x^2/2)};
        \addlegendentry{$\equ{\psi_3}$}
    \end{axis}
  \end{tikzpicture}
  \caption{Primeras eigenfunciones $\equ{\psi}_n = H_n e^{-\equ{x}^2/2}$
           del oscilador armónico cuántico.}
  \label{fig:gal:harmonic}
\end{figure}

%───────────────────────────────────────────────────────────────────────────────
\section{Diagramas TikZ}
\label{sec:gal:diagramas}
%───────────────────────────────────────────────────────────────────────────────

% CRÍTICO: nunca Ctp* en TikZ — usar tikz-*, env-*, fb-*, eq-* según el rol.
% tikz-initial (Blue): estado/nodo inicial
% tikz-accept (Green): estado aceptador / nodo destino
% tikz-node-border / background / maintext: nodo genérico
% eq-known (Mauve): pesos de aristas (datos conocidos del problema)
% env-phys (Sapphire): caja de función / transformación

\subsection{Autómata finito determinista}

\begin{figure}[htbp]
  \centering
  \begin{tikzpicture}[
    shorten >=1pt, node distance=2.5cm, on grid, auto, thick,
    state/.style={circle, draw=tikz-node-border, fill=background,
                  text=maintext, minimum size=1cm},
    every edge/.style={draw=tikz-edge, ->, >=latex},
  ]
    \node[state, initial, initial text={}, draw=tikz-initial]
      (A) {$\eqk{A}$};
    \node[state, right=of A] (B) {$\eqk{B}$};
    \node[state, right=of B] (C) {$\eqk{C}$};
    \node[state, accepting, draw=tikz-accept, right=of C]
      (D) {$\eqk{D}$};
    \path[->]
      (A) edge[loop above] node {$\eqk{a}$} (A)
          edge              node {$\eqk{b}$} (B)
      (B) edge[loop above] node {$\eqk{b}$} (B)
          edge              node {$\eqk{a}$} (C)
      (C) edge              node {$\eqk{b}$} (D)
          edge[bend left]   node[swap] {$\eqk{a}$} (A)
      (D) edge[loop right]  node {$\eqk{a},\eqk{b}$} (D);
  \end{tikzpicture}
  \caption{DFA para $L = \{\eqvec{w} \in \eqdom{\{a,b\}^*} \mid
           \eqvec{w}$ contiene $\eqk{ab}\}$~\cite{sipser2012}.
           Estado inicial \texttt{tikz-initial} (azul);
           aceptador \texttt{tikz-accept} (verde).}
  \label{fig:gal:dfa}
\end{figure}

\subsection{Grafo dirigido ponderado}

\begin{figure}[htbp]
  \centering
  \begin{tikzpicture}[
    node distance=2.5cm,
    vertex/.style={circle, draw=tikz-node-border, fill=tikz-node-bg,
                   text=maintext, minimum size=0.9cm, thick},
    every edge/.style={draw=tikz-edge, ->, >=latex, thick},
    wlabel/.style={text=eq-known, font=\small},
  ]
    \node[vertex, draw=tikz-initial, fill=tikz-initial!15] (s) {$\eqk{s}$};
    \node[vertex, right=of s]  (u) {$\eqk{u}$};
    \node[vertex, right=of u]  (v) {$\eqk{v}$};
    \node[vertex, below=of u]  (x) {$\eqk{x}$};
    \node[vertex, draw=tikz-accept, fill=tikz-accept!15,
          right=of x]          (t) {$\eqk{t}$};
    \path[->]
      (s) edge node[above, wlabel] {$\eqk{6}$} (u)
      (s) edge node[left,  wlabel] {$\eqk{7}$} (x)
      (u) edge node[above, wlabel] {$\eqk{5}$} (v)
      (u) edge node[right, wlabel] {$\eqk{-4}$}(t)
      (x) edge node[below, wlabel] {$\eqk{7}$} (t)
      (v) edge node[right, wlabel] {$\eqk{-2}$}(t);
  \end{tikzpicture}
  \caption{Grafo ponderado~\cite{cormen2022}: pesos en \texttt{eq-known} (malva)
           porque son datos del problema de camino mínimo.}
  \label{fig:gal:grafo}
\end{figure}

\subsection{Tabla hash con separate chaining}

\begin{figure}[htbp]
  \centering
  \begin{tikzpicture}[font=\small,
    cell/.style  = {draw=tikz-node-border, minimum width=1.3cm,
                    minimum height=0.65cm},
    occ/.style   = {cell, fill=tikz-accept,  text=background},
    empty/.style = {cell, fill=tikz-node-bg, text=maintext},
    fn/.style    = {draw=tikz-node-border, fill=env-phys, text=background,
                    minimum width=1.2cm, minimum height=3.2cm, align=center},
    key/.style   = {cell, fill=tikz-initial, text=background,
                    minimum width=2cm},
    arr/.style   = {draw=tikz-edge, ->, >=latex, thick},
  ]
    \node[key] (k1) at (-4.5, 0.9) {\texttt{k1}};
    \node[key] (k2) at (-4.5, 0.0) {\texttt{k2}};
    \node[key] (k3) at (-4.5,-0.9) {\texttt{k3}};
    \node[fn]  (hf) at (-2.4, 0)   {$\eqop{h}$};
    \foreach \i in {0,...,4} {
      \node[empty] (b\i) at (0.5, {0.9 - \i*0.65}) {};
      \node[left=0.1cm of b\i, text=tikz-label, font=\scriptsize] {$\i$};
    }
    \node[occ] at (b1) {\texttt{k1}};
    \node[occ] at (b3) {\texttt{k2}};
    \node[occ] (c1) at (2.1, {0.9 - 1*0.65}) {\texttt{k3}};
    \draw[arr] (b1) -- (c1);
    \draw[arr] (k1) -- (hf);
    \draw[arr] (k2) -- (hf);
    \draw[arr] (k3) -- (hf);
    \draw[arr] (hf.east) to[out=0,in=180] (b1.west);
    \draw[arr] (hf.east) to[out=0,in=180] (b3.west);
    \node[above=0.3cm of k1, text=maintext, font=\small] {\textbf{Claves}};
    \node[above=0.3cm of hf, text=maintext, font=\small] {\textbf{Hash}};
    \node[above=0.3cm of b0, text=maintext, font=\small] {\textbf{Buckets}};
  \end{tikzpicture}
  \caption{Separate chaining~\cite{cormen2022}: función hash $\eqop{h}$
           (\texttt{env-phys}, sapphire) mapea claves (\texttt{tikz-initial},
           azul) a buckets (verde=ocupado, gris=libre).}
  \label{fig:gal:hashing}
\end{figure}

\subsection{Diagrama conmutativo}

\begin{figure}[htbp]
  \centering
  \begin{tikzcd}[
    every arrow/.append style={draw=tikz-edge, thick},
  ]
    \eqdom{G}
      \arrow[r, "\eqop{\varphi}"]
      \arrow[d, "\eqop{\pi}"', two heads]
    & \eqdom{H} \\
    \eqdom{G}\,\eqop{/}\,\eqdom{\ker\varphi}
      \arrow[ur, dashed, "\eqop{\tilde{\varphi}}"', "\sim"']
  \end{tikzcd}
  \caption{Primer teorema de isomorfismo~\cite{mac1998}:
           $\eqdom{G}/\eqdom{\ker\varphi} \cong \eqop{\mathrm{im}}(\eqop{\varphi})$.}
  \label{fig:gal:isomorfismo}
\end{figure}

%───────────────────────────────────────────────────────────────────────────────
\section{Elementos auxiliares}
\label{sec:gal:auxiliares}
%───────────────────────────────────────────────────────────────────────────────

\dificultad{3}{4}

\begin{ejercicio}
Demuestre que la pérdida MSE para regresión lineal conduce a las ecuaciones
normales $\eqten{X}^\top\eqten{X}\equ{\hat{\beta}} = \eqten{X}^\top\eqk{y}$.
\end{ejercicio}

\begin{ejercicio}
Sea $\eqfn{f}(\equ{x}) = e^{-\equ{x}^2}$. Calcule la transformada de
Fourier $\eqop{\mathcal{F}}[\eqfn{f}]$ y verifique que es un punto fijo
de $\eqop{\mathcal{F}}$ (salvo escala).
\end{ejercicio}

\begin{notasbib}
La teoría de espacios de Hilbert y el análisis funcional riguroso siguen
a Reed y Simon~\cite{reed1980}. Los métodos de score matching y su variante
denoising~\cite{hyvarinen2005,vincent2011} son fundamentales en los
modelos de difusión modernos. La formulación variacional de EDPs y los
métodos de diferencias finitas se tratan en~\cite{leveque2007}.
\end{notasbib}

%───────────────────────────────────────────────────────────────────────────────
\printbibliography[heading=subbibliography]

\end{appendices}

\end{document}
